% Integration activities / Evaluation / Remediation (Chapter 1) — Anglais 1ère A — Cours signé OnBuch+
% Programme officiel MINESEC. Compiler avec : tectonic ENA13-integration-evaluation-remediation-13.tex
\newcommand{\DOCMATIERE}{Anglais}
\newcommand{\DOCNIVEAU}{1ère A}
\newcommand{\DOCMODULE}{Chapter 1 : Family and Social Life}
\newcommand{\DOCLECON}{ENA13}
\newcommand{\DOCTITRE}{Integration activities / Evaluation / Remediation (Chapter 1)}
% OnBuch+ — préambule « cours signés » ENRICHI (compilé avec tectonic / XeLaTeX)
% Système visuel « Pop » d'OnBuch+ : fond crème (jamais blanc), bordures encre
% épaisses, ombres « dures » décalées, titres Archivo Black en capitales.
% Version MATHÉMATIQUES : graphes (pgfplots/tikz), boîtes pédagogiques
% (définition, propriété, méthode, attention, le savais-tu, exercice résolu,
% exercice, TP, bilan), page de garde et signature OnBuch+ ancrée en bas.
%
% Macros attendues dans main.tex AVANT \input{preamble} :
%   \DOCMATIERE  \DOCNIVEAU  \DOCMODULE  \DOCLECON (numéro)  \DOCTITRE
\documentclass[11pt]{article}

\usepackage{fontspec}
\usepackage[english,french]{babel} % français = langue principale ; l'anglais via \eng{..} / textebox
\usepackage[a4paper,margin=2cm,top=3.4cm,bottom=2.8cm,headheight=1.4cm,footskip=1.3cm]{geometry}
\usepackage{amsmath,amssymb,amsfonts}
\usepackage{mathtools} % \xmapsto, \coloneqq... souvent utilisés par les modèles
\usepackage{cancel} % simplifications barrées (bilans de forces)
% \abs{z} (module, valeur absolue) et \norm{u} (norme), utilisés par les modèles
\providecommand{\abs}{}\let\abs\relax\DeclarePairedDelimiter{\abs}{\lvert}{\rvert}
\providecommand{\norm}{}\let\norm\relax\DeclarePairedDelimiter{\norm}{\lVert}{\rVert}
\usepackage[table]{xcolor} % [table] : fournit aussi \rowcolors (alternance de lignes)
\usepackage{pagecolor}
\usepackage{tikz}
\usetikzlibrary{arrows.meta,positioning,calc,shapes.geometric,shapes.misc,shapes.arrows,decorations.pathmorphing,decorations.pathreplacing,decorations.markings,patterns,shadows,mindmap,trees,fit,backgrounds,matrix,intersections,angles,quotes}
\usepackage{pgfplots}
\usepackage[version=4]{mhchem} % équations nucléaires \ce{^{235}_{92}U}
\usepackage[european,siunitx]{circuitikz} % schémas électriques
\pgfplotsset{compat=1.18}
\usepgfplotslibrary{fillbetween,groupplots} % aires entre courbes (calcul intégral)
\usepackage{enumitem}
\usepackage{fancyhdr}
% [most] charge toutes les bibliothèques tcolorbox (listings, minted, raster,
% documentation...) et gonfle inutilement la mémoire TeX sur un document long
% (~300 boîtes) : on ne charge que ce qui est réellement utilisé.
\usepackage[skins,breakable]{tcolorbox}
\usepackage{graphicx}
\usepackage{colortbl}
\usepackage{booktabs}
\usepackage{tabularx}
\usepackage{diagbox} % case d'en-tête diagonale des tableaux à double entrée
% Colonnes extensibles centrée / gauche / droite (C, L, R), souvent utilisées par les modèles
\newcolumntype{C}{>{\centering\arraybackslash}X}
\newcolumntype{L}{>{\raggedright\arraybackslash}X}
\newcolumntype{R}{>{\raggedleft\arraybackslash}X}
\usepackage{array}
\usepackage{multicol}
\usepackage{multirow}
\usepackage{siunitx}
% parse-numbers=false : accepte aussi \SI{2,5 \times 10^{-3}}{...}
\sisetup{locale=FR,output-decimal-marker={,},group-minimum-digits=4,parse-numbers=false}
\DeclareSIUnit\ohm{\ensuremath{\symup{\Omega}}} % U+2126 absent de la police
\DeclareSIUnit\division{div} % réglages d'oscilloscope (ms/div, V/div)
\DeclareSIUnit\tr{tr} % tours (vitesse de rotation, ex. \SI{1500}{\tr\per\minute})
\DeclareSIUnit\molar{M} % concentration molaire (ex. \SI{3}{\molar}, \SI{1}{\micro\molar})
\DeclareSIUnit\an{an} % année (le modèle confond parfois avec \year, un compteur TeX) — ex. \SI{5}{\kilo\meter\per\an}
\DeclareSIUnit\FCFA{FCFA} % franc CFA (ex. \SI{500}{\FCFA\per\kilo\gram})
\DeclareSIUnit\degreeBrix{\textdegree Brix} % teneur en sucre d'un sirop/jus (réfractométrie)
\DeclareSIUnit\month{mois} % le modèle utilise parfois \month, un compteur TeX, comme unité de durée
\usepackage{qrcode}
% \degree : commande fréquemment utilisée par les modèles pour un angle ou une
% température en dehors de \SI{}{} (non fournie par les paquets chargés).
\providecommand{\degree}{^{\circ}}
% \qed (fin de démonstration), utilisé par les modèles sans amsthm
\providecommand{\qed}{\hfill$\square$}
% \barre : double barre des valeurs interdites dans un tableau de variations (tkz-tab)
\providecommand{\barre}{\ensuremath{\|}}
% environnement proof (amsthm non chargé) utilisé par les modèles
\ifdefined\proof\else\newenvironment{proof}[1][Démonstration]{\par\noindent\textit{#1.}\ }{\qed\par}\fi
\usepackage{lastpage}
\usepackage{hyperref}
\hypersetup{hidelinks}

% ============================================================
% Palette « Pop » OnBuch+ — mêmes hex que la classe P (plus_ui.dart)
% ============================================================
\definecolor{poporange}{HTML}{F59321}
\definecolor{popdark}{HTML}{E07A0C}
\definecolor{popcream}{HTML}{FFF6E8}
\definecolor{popink}{HTML}{1C1714}
\definecolor{poppurple}{HTML}{7A5AE0}
\definecolor{popgreen}{HTML}{2E9E6A}
\definecolor{poppink}{HTML}{E0457B}
\definecolor{popblue}{HTML}{2F7DE1}
\definecolor{popgold}{HTML}{C98A12}
\definecolor{popmuted}{HTML}{8A7F76}
\definecolor{popline}{HTML}{EADFD2}
\definecolor{popgray}{HTML}{8A7F76}
\colorlet{popgrey}{popgray}
% Alias tolérants (le modèle nomme parfois les couleurs différemment)
\colorlet{popwhite}{white}
\colorlet{popblack}{popink}
\colorlet{popred}{poppink}
\colorlet{popyellow}{popgold}
% Teintes claires (fonds de tableaux/encadrés) — le modèle nomme parfois
% les couleurs avec un suffixe L (ex. popblueL) sans les avoir définies.
\colorlet{poporangeL}{poporange!20}
\colorlet{popdarkL}{popdark!20}
\colorlet{poppurpleL}{poppurple!20}
\colorlet{popgreenL}{popgreen!20}
\colorlet{poppinkL}{poppink!20}
\colorlet{popblueL}{popblue!20}
\colorlet{popgoldL}{popgold!20}
\colorlet{popredL}{popred!20}
\colorlet{popgrayL}{popgray!20}
% Teintes claires pour fonds de tableaux / zones
\colorlet{popblueL}{popblue!10!white}
\colorlet{poporangeL}{poporange!14!white}
\colorlet{popgreenL}{popgreen!12!white}
\colorlet{poppurpleL}{poppurple!12!white}
\colorlet{poppinkL}{poppink!12!white}
\colorlet{popgoldL}{popgold!14!white}

\pagecolor{popcream}

% ============================================================
% Typographie
% ============================================================
\defaultfontfeatures{Ligatures=TeX}
\setmainfont{PlusJakartaSans-Medium.ttf}[Path=../../../fonts/,BoldFont=PlusJakartaSans-Bold.ttf]
% Maths en sans-serif (Fira Math) avec chiffres et lettres droites de Plus
% Jakarta, pour que les formules se fondent dans le texte.
\usepackage[math-style=ISO,bold-style=ISO]{unicode-math}
\setmathfont{FiraMath-Regular.otf}[Path=../../../fonts/]
\setmathfont{PlusJakartaSans-Medium.ttf}[Path=../../../fonts/,range={up/num,up/{Latin,latin},\mathup}]
\usepackage{icomma} % virgule décimale sans espace en mode math
% Caractères mathématiques Unicode absents des polices de texte (Plus Jakarta,
% Archivo Black) : rendus par leur équivalent mathématique, en texte comme en
% formule — sinon ils s'affichent en carré vide (ex. « Arithmétique dans ℤ »).
\usepackage{newunicodechar}
\newunicodechar{ℝ}{\ensuremath{\mathbb{R}}}
\newunicodechar{ℤ}{\ensuremath{\mathbb{Z}}}
\newunicodechar{ℂ}{\ensuremath{\mathbb{C}}}
\newunicodechar{ℕ}{\ensuremath{\mathbb{N}}}
\newunicodechar{ℚ}{\ensuremath{\mathbb{Q}}}
\newunicodechar{𝔻}{\ensuremath{\mathbb{D}}}
\newunicodechar{α}{\ensuremath{\alpha}}
\newunicodechar{β}{\ensuremath{\beta}}
\newunicodechar{γ}{\ensuremath{\gamma}}
\newunicodechar{δ}{\ensuremath{\delta}}
\newunicodechar{ε}{\ensuremath{\varepsilon}}
\newunicodechar{θ}{\ensuremath{\theta}}
\newunicodechar{λ}{\ensuremath{\lambda}}
\newunicodechar{μ}{\ensuremath{\mu}}
\newunicodechar{σ}{\ensuremath{\sigma}}
\newunicodechar{τ}{\ensuremath{\tau}}
\newunicodechar{φ}{\ensuremath{\varphi}}
\newunicodechar{ψ}{\ensuremath{\psi}}
\newunicodechar{ω}{\ensuremath{\omega}}
\newunicodechar{Σ}{\ensuremath{\Sigma}}
% majuscules produites par \MakeUppercase dans les titres (α-aminés -> Α)
\newunicodechar{Α}{\ensuremath{\alpha}}
\newunicodechar{Β}{\ensuremath{\beta}}
\newunicodechar{Γ}{\ensuremath{\Gamma}}
\newunicodechar{Δ}{\ensuremath{\Delta}}
\newunicodechar{Ε}{\ensuremath{\varepsilon}}
\newunicodechar{Θ}{\ensuremath{\Theta}}
\newunicodechar{Λ}{\ensuremath{\Lambda}}
\newunicodechar{Μ}{\ensuremath{\mu}}
\newunicodechar{Τ}{\ensuremath{\tau}}
\newunicodechar{Φ}{\ensuremath{\Phi}}
\newunicodechar{Ψ}{\ensuremath{\Psi}}
\newunicodechar{Ω}{\ensuremath{\Omega}}
\newunicodechar{↦}{\ensuremath{\mapsto}}
\newunicodechar{∈}{\ensuremath{\in}}
\newunicodechar{∉}{\ensuremath{\notin}}
\newunicodechar{∧}{\ensuremath{\wedge}}
\newunicodechar{⇔}{\ensuremath{\Leftrightarrow}}
\newunicodechar{⟺}{\ensuremath{\Longleftrightarrow}}
\newunicodechar{⇒}{\ensuremath{\Rightarrow}}
\newunicodechar{⟹}{\ensuremath{\Longrightarrow}}
\newunicodechar{∘}{\ensuremath{\circ}}
\newunicodechar{∪}{\ensuremath{\cup}}
\newunicodechar{∩}{\ensuremath{\cap}}
\newunicodechar{≡}{\ensuremath{\equiv}}
\newunicodechar{⊂}{\ensuremath{\subset}}
\newunicodechar{⊥}{\ensuremath{\perp}}
\newunicodechar{∃}{\ensuremath{\exists}}
\newunicodechar{∀}{\ensuremath{\forall}}
\newunicodechar{∛}{\ensuremath{\sqrt[3]{\,}}}
\newunicodechar{∜}{\ensuremath{\sqrt[4]{\,}}}
\newunicodechar{⃗}{\ensuremath{{}^{\to}}}
\newunicodechar{ⁿ}{\ifmmode^{n}\else\textsuperscript{\ensuremath{n}}\fi}
\newunicodechar{ˣ}{\ifmmode^{x}\else\textsuperscript{\ensuremath{x}}\fi}
\newunicodechar{ᵇ}{\ifmmode^{b}\else\textsuperscript{\ensuremath{b}}\fi}
\newunicodechar{ʳ}{\ifmmode^{r}\else\textsuperscript{\ensuremath{r}}\fi}
\newunicodechar{ᵘ}{\ifmmode^{u}\else\textsuperscript{\ensuremath{u}}\fi}
\newunicodechar{ʸ}{\ifmmode^{y}\else\textsuperscript{\ensuremath{y}}\fi}
\newunicodechar{ᵃ}{\ifmmode^{a}\else\textsuperscript{\ensuremath{a}}\fi}
\newunicodechar{ᵉ}{\ifmmode^{e}\else\textsuperscript{\ensuremath{e}}\fi}
\newunicodechar{ᵏ}{\ifmmode^{k}\else\textsuperscript{\ensuremath{k}}\fi}
\newunicodechar{ᵖ}{\ifmmode^{p}\else\textsuperscript{\ensuremath{p}}\fi}
\newunicodechar{ᴺ}{\ifmmode^{N}\else\textsuperscript{\ensuremath{N}}\fi}
\newunicodechar{ᶜ}{\ifmmode^{c}\else\textsuperscript{\ensuremath{c}}\fi}
\newunicodechar{ʰ}{\ifmmode^{h}\else\textsuperscript{\ensuremath{h}}\fi}
\newunicodechar{ᵠ}{\ifmmode^{q}\else\textsuperscript{\ensuremath{q}}\fi}
\newunicodechar{ₙ}{\ifmmode_{n}\else\textsubscript{\ensuremath{n}}\fi}
\newunicodechar{ₐ}{\ifmmode_{a}\else\textsubscript{\ensuremath{a}}\fi}
\newunicodechar{ᵢ}{\ifmmode_{i}\else\textsubscript{\ensuremath{i}}\fi}
\newunicodechar{ᵣ}{\ifmmode_{r}\else\textsubscript{\ensuremath{r}}\fi}
\newunicodechar{ₖ}{\ifmmode_{k}\else\textsubscript{\ensuremath{k}}\fi}
\newunicodechar{ᵦ}{\ifmmode_{\beta}\else\textsubscript{\ensuremath{\beta}}\fi}
\newunicodechar{ₓ}{\ifmmode_{x}\else\textsubscript{\ensuremath{x}}\fi}
\newfontfamily\popdisplay{ArchivoBlack-Regular.ttf}[Path=../../../fonts/]
\newfontfamily\poplogo{SpaceGrotesk-Bold.ttf}[Path=../../../fonts/]
\newfontfamily\popsemi{PlusJakartaSans-SemiBold.ttf}[Path=../../../fonts/]
\newfontfamily\popxbold{PlusJakartaSans-ExtraBold.ttf}[Path=../../../fonts/]
\newfontfamily\popxboldls{PlusJakartaSans-ExtraBold.ttf}[Path=../../../fonts/,LetterSpace=3.0]

\setlist[enumerate]{leftmargin=1.7em,itemsep=5pt,topsep=4pt}
\setlist[itemize]{leftmargin=1.7em,itemsep=4pt,topsep=4pt,label={\color{poporange}\textbullet}}
\setlength{\parindent}{0pt}
\setlength{\parskip}{5pt}
\renewcommand{\arraystretch}{1.25}

% pgfplots : style Pop par défaut
\pgfplotsset{
  every axis/.append style={axis line style={popink,line width=1pt},tick style={popink},
    grid style={popline,line width=0.5pt},label style={font=\small},tick label style={font=\footnotesize}},
  popcurve/.style={poporange,line width=1.8pt,no markers,smooth},
}
% Même style disponible dans un \draw TikZ simple (hors pgfplots)
\tikzset{popcurve/.style={poporange,line width=1.8pt}}
\tikzset{popgrid/.style={popline,very thin}}
\colorlet{popcurve}{poporange} % le nom du style est parfois utilisé comme couleur

% ============================================================
% Pill / squiggle
% ============================================================
\newcommand{\poppill}[3]{%
  \tikz[baseline=-0.55ex]\node[draw=popink,line width=1.2pt,rounded corners=2.6pt,%
    fill=#2,inner xsep=6.5pt,inner ysep=3pt]{%
    \popxboldls\fontsize{7}{7}\selectfont\color{#3}\MakeUppercase{#1}};%
}
\newcommand{\popsquiggle}[1]{%
  \tikz[baseline=-0.4ex]{
    \draw[#1,line width=2.6pt,line cap=round]
      plot [smooth] coordinates {(0,0.09) (0.34,0.01) (0.68,0.21) (1.02,0.01) (1.36,0.10)};
  }%
}
% Mot-clé surligné (marqueur orange sous le texte)
\newcommand{\cle}[1]{{\popxbold\color{popdark}#1}}

% ============================================================
% En-tête / pied
% ============================================================
\pagestyle{fancy}
\fancyhf{}
\renewcommand{\headrulewidth}{0pt}
\renewcommand{\footrulewidth}{0pt}
\lhead{\raisebox{-0.15ex}{\poplogo\fontsize{13}{13}\selectfont\color{popink}OnBuch\color{poporange}+}}
\rhead{\poppill{\DOCMATIERE\ \textbullet\ \DOCNIVEAU\ \textbullet\ Leçon \DOCLECON}{white}{popink}}
\lfoot{\popxbold\fontsize{7.6}{7.6}\selectfont\color{popmuted}\MakeUppercase{Cours signé OnBuch+}\ \textbullet\ {\popsemi\DOCTITRE}}
\rfoot{\tikz[baseline=-0.6ex]\node[rounded corners=8.5pt,fill=popink,minimum height=17pt,minimum width=17pt,inner xsep=5pt,inner sep=0]{\popxbold\fontsize{8}{8}\selectfont\color{popcream}\thepage\,/\,\pageref*{LastPage}};}

% ============================================================
% Titres
% ============================================================
\newcommand{\chaptitre}[2]{%
  \begin{tcolorbox}[enhanced,colback=poporange,colframe=popink,boxrule=2.6pt,arc=11pt,
      drop shadow=popink,left=15pt,right=15pt,top=12pt,bottom=11pt,before skip=3pt,after skip=18pt]
    \poppill{#2}{popink}{popcream}\par\vspace{9pt}
    {\raggedright\hyphenpenalty=10000\exhyphenpenalty=10000\popdisplay\fontsize{22}{24}\selectfont\color{popcream}\MakeUppercase{#1}\par}
  \end{tcolorbox}
}
% Section numérotée : pastille orange avec le numéro + titre Archivo
\newcounter{popsec}
\newcommand{\coursec}[1]{%
  \stepcounter{popsec}\setcounter{popsub}{0}%
  \par\vspace{16pt}\noindent
  \begin{minipage}{\linewidth}
  \tikz[baseline=-0.7ex]\node[draw=popink,line width=1.6pt,fill=poporange,rounded corners=4pt,minimum size=22pt,inner sep=0]{\popdisplay\fontsize{12}{12}\selectfont\color{popcream}\thepopsec};%
  \hspace{8pt}{\popdisplay\fontsize{15}{17}\selectfont\color{popink}\MakeUppercase{#1}}\par
  \vspace{4pt}\noindent{\color{poporange}\rule{36pt}{3.6pt}}
  \end{minipage}\par\nopagebreak\vspace{8pt}
}
\newcounter{popsub}
\newcommand{\courssub}[1]{%
  \stepcounter{popsub}%
  \par\vspace{9pt}\noindent{\popxbold\fontsize{11.5}{13}\selectfont\color{popink}{\color{poporange}\thepopsec.\thepopsub}\hspace{6pt}#1}\par\nopagebreak\vspace{4pt}
}

% ============================================================
% Boîtes pédagogiques — même recette : fond blanc, bordure épaisse colorée,
% ombre dure encre, badge plein. Titre optionnel [#1].
% ============================================================
\tcbset{popbase/.style={breakable,enhanced,colback=white,boxrule=2.3pt,arc=9pt,
  shadow={3pt}{-3pt}{0pt}{popink},left=14pt,right=14pt,top=11pt,bottom=11pt,before skip=11pt,after skip=15pt,
  fontupper=\popsemi\fontsize{10.6}{14.5}\selectfont\color{popink},
  fontlower=\popsemi\fontsize{10.6}{14.5}\selectfont\color{popink}}}
% Coche dessinée (le glyphe ✓ n'existe pas dans les polices du document)
\newcommand{\popcheck}[1]{\tikz[baseline=-0.5ex]\draw[#1,line width=1.6pt,line cap=round,line join=round](-0.13,0.0)--(-0.04,-0.09)--(0.13,0.1);}
\providecommand{\coche}{\popcheck{popgreen}}
\providecommand{\resultat}[1]{\par\smallskip\noindent\fcolorbox{popgreen}{popgreenL}{\parbox{\dimexpr\linewidth-2\fboxsep-2\fboxrule\relax}{\textbf{Résultat.} #1}}\par\smallskip}
\newcommand{\popboxhead}[3]{\poppill{#1}{#2}{white}\ifx\relax#3\relax\else\hspace{6pt}{\popxbold\fontsize{10.6}{12}\selectfont\color{popink}#3}\fi\par\vspace{7pt}}

\newtcolorbox{aretenir}[1][]{popbase,colframe=popgreen,before upper={\popboxhead{À retenir}{popgreen}{#1}}}
% Texte en anglais (document de lecture, dialogue, extrait) : cadre
% d'étude. Un titre optionnel (source, auteur) s'affiche dans l'en-tête.
\newtcolorbox{textebox}[1][]{popbase,colframe=popblue,before upper={\popboxhead{Text}{popblue}{#1}\begin{otherlanguage}{english}},after upper={\end{otherlanguage}}}
% Marques ✓ / ✗ (forme correcte / fautive) souvent utilisées par les modèles
\providecommand{\cross}{\textcolor{poppink}{\ensuremath{\times}}}
\providecommand{\xmark}{\cross}\providecommand{\crossmark}{\cross}\providecommand{\cmark}{\ensuremath{\checkmark}}
% Ligne de réponse / justification à compléter (inventée par les modèles)
\providecommand{\justification}[1]{\par\noindent\textit{Justification~:} \dotfill\par}
% \textipa (package tipa non chargé) : la police ne couvre pas l'API -> simple passage
\providecommand{\textipa}[1]{#1}
% Soulignement (ulem non chargé) : \uline, \ul
\providecommand{\uline}[1]{\underline{#1}}\providecommand{\ul}[1]{\underline{#1}}
% \glossaire{\item ...} inventée par les modèles : liste de glossaire
\providecommand{\glossaire}[1]{\begin{itemize}#1\end{itemize}}
% \alert (beamer) inventée par les modèles : texte mis en évidence
\providecommand{\alert}[1]{\textbf{\textcolor{poppink}{#1}}}
% variantes ASCII d'ordinaux écrites en macro
\providecommand{\iere}{\textsuperscript{re}}\providecommand{\ieme}{\textsuperscript{e}}\providecommand{\ere}{\textsuperscript{re}}\providecommand{\eme}{\textsuperscript{e}}\providecommand{\er}{\textsuperscript{er}}\providecommand{\ie}{\textsuperscript{e}}
% Ordinaux français écrits en macro par les modèles : 1\ière, 3\ième
\newcommand{\ière}{\textsuperscript{re}}\newcommand{\ième}{\textsuperscript{e}}
% \exemple (inventée par les modèles) : étiquette « Exemple : »
\providecommand{\exemple}{\textbf{Exemple~:}\ }
% \square utilisable aussi hors mode math (cases à cocher)
\AtBeginDocument{\let\squareMath\square\renewcommand{\square}{\ensuremath{\squareMath}}}
% Mot ou phrase en anglais à l'intérieur d'un paragraphe français
\newcommand{\eng}[1]{\ifmmode\text{\textit{#1}}\else\begingroup\selectlanguage{english}\itshape #1\endgroup\fi}
\let\esp\eng % alias toléré
% flèches d'intonation inventées par les modèles
\providecommand{\textsearrow}{}\renewcommand{\textsearrow}{\ensuremath{\searrow}}\providecommand{\textnearrow}{}\renewcommand{\textnearrow}{\ensuremath{\nearrow}}
% \fr{..} : mot français (inventé par les modèles)
\providecommand{\fr}[1]{\textit{#1}}
\newtcolorbox{exemplebox}[1][]{popbase,colframe=poporange,before upper={\popboxhead{Exemple}{poporange}{#1}}}
\newtcolorbox{definition}[1][]{popbase,colframe=popblue,before upper={\popboxhead{Définition}{popblue}{#1}}}
\newtcolorbox{propriete}[1][]{popbase,colframe=poppurple,before upper={\popboxhead{Propriété}{poppurple}{#1}}}
\newtcolorbox{methode}[1][]{popbase,colframe=popgold,before upper={\popboxhead{Méthode}{popgold}{#1}}}
\newtcolorbox{attention}[1][]{popbase,colframe=poppink,before upper={\popboxhead{Attention}{poppink}{#1}}}
\newtcolorbox{experience}[1][]{popbase,colframe=popblue,colback=popblueL,before upper={\popboxhead{Expérience}{popblue}{#1}}}
% « Le savais-tu ? » : carte violette pleine (contexte, culture, Cameroun)
\newtcolorbox{savaistu}[1][]{popbase,colback=poppurple,colframe=popink,
  fontupper=\popsemi\fontsize{10.4}{14.2}\selectfont\color{white},
  before upper={\poppill{Le savais-tu\,?}{white}{poppurple}\ifx\relax#1\relax\else\hspace{6pt}{\popxbold\color{white}#1}\fi\par\vspace{7pt}}}
% Exercice résolu : énoncé en haut, \tcblower puis la solution sur fond crème
\newcounter{popexo}\newcounter{popexr}
\newtcolorbox{exoresolu}[1][]{popbase,colframe=poporange,colbacklower=popcream,
  segmentation style={popink,line width=1.2pt,dashed},
  before upper={\stepcounter{popexr}\popboxhead{Exercice résolu \thepopexr}{poporange}{#1}},
  before lower={\poppill{Solution}{popink}{popcream}\par\vspace{6pt}}}
% Exercice (non résolu en place) — la correction est dans la partie Corrigés
\newtcolorbox{exercice}[1][]{popbase,colframe=popink,
  before upper={\stepcounter{popexo}\popboxhead{Exercice \thepopexo}{popink}{#1}}}
% Corrigé
\newtcolorbox{corrige}[1][]{popbase,colframe=popgreen,colback=popgreenL,
  before upper={\popboxhead{Corrigé}{popgreen}{#1}}}
% Objectifs / prérequis / bilan
\newtcolorbox{objectifs}{popbase,colframe=popink,colback=poporangeL,
  before upper={\popboxhead{Objectifs de la leçon}{popink}{}}}
\newtcolorbox{prerequis}{popbase,colframe=popink,colback=popblueL,
  before upper={\popboxhead{Prérequis}{popblue}{}}}
\newtcolorbox{bilan}{popbase,colframe=popink,colback=white,boxrule=2.8pt,
  before upper={\popboxhead{Fiche bilan}{poporange}{L'essentiel à connaître pour l'examen}}}
% Figure encadrée avec légende
\newtcolorbox{popfigure}[1][]{enhanced,breakable=false,colback=white,colframe=popline,boxrule=1.4pt,arc=8pt,
  left=10pt,right=10pt,top=10pt,bottom=8pt,before skip=10pt,after skip=12pt,halign=center,
  lower separated=false,halign lower=center,
  fontlower=\popsemi\fontsize{9}{11.5}\selectfont\color{popmuted},
  #1}
\newcommand{\legende}[1]{\par\vspace{6pt}{\popsemi\fontsize{9}{11.5}\selectfont\color{popmuted}#1}}

% ============================================================
% Signature OnBuch+
% ============================================================
\newcommand{\poplogoinline}[1]{{\poplogo\fontsize{#1}{#1}\selectfont\color{popink}OnBuch\color{poporange}+}}
\newcommand{\signatureonbuch}{%
  \begin{tcolorbox}[enhanced,colback=popink,colframe=popink,boxrule=2.2pt,arc=10pt,
      drop shadow=poporange,left=14pt,right=14pt,top=10pt,bottom=10pt,before skip=0pt,after skip=0pt]
    \tikz[baseline=-0.7ex]\node[circle,fill=popgreen,draw=popcream,line width=1.3pt,minimum size=22pt,inner sep=0]{\popcheck{white}};%
    \hspace{9pt}\begin{minipage}[c]{0.62\linewidth}
      {\popxbold\fontsize{12}{13}\selectfont\color{popcream}Signé\ \textbullet\ {\poplogo OnBuch\color{poporange}+}}\par\vspace{2pt}
      {\raggedright\popsemi\fontsize{8}{10.5}\selectfont\color{popline}\MakeUppercase{Équipe pédagogique OnBuch+}\\\MakeUppercase{Conforme au programme officiel MINESEC}\par}
    \end{minipage}\hfill
    {\poplogo\fontsize{22}{22}\selectfont\color{popcream}OnBuch\color{poporange}+}
  \end{tcolorbox}%
}

% ============================================================
% Page de garde
% #1 titre · #2 matière · #3 niveau · #4 module · #5 n° leçon · #6 durée ·
% #7 description / objectifs (texte ou itemize)
% ============================================================
\newcommand{\pagegarde}[7]{%
  \thispagestyle{empty}
  \newgeometry{a4paper,margin=1.7cm,top=1.6cm,bottom=1.4cm}%
  \begin{tikzpicture}[remember picture,overlay]
    \fill[poporange!18] ([xshift=-3.2cm,yshift=-3.6cm]current page.north east) circle (4.6cm);
    \fill[poppurple!14] ([xshift=2.4cm,yshift=7.6cm]current page.south west) circle (3.8cm);
  \end{tikzpicture}%
  {\poplogo\fontsize{30}{30}\selectfont\color{popink}OnBuch\color{poporange}+}\hfill
  \raisebox{0.6ex}{\poppill{Cours signé \textbullet\ Premium}{popink}{popcream}}\par
  \vspace{1pt}\popsquiggle{poppurple}\par
  \vspace{8pt}
  \begin{tcolorbox}[enhanced,colback=poporange,colframe=popink,boxrule=2.8pt,arc=13pt,
      drop shadow=popink,left=18pt,right=18pt,top=12pt,bottom=13pt,after skip=10pt]
    \poppill{#2 \textbullet\ #3}{popink}{popcream}\hspace{5pt}\poppill{#4}{white}{popink}\par\vspace{8pt}
    {\popdisplay\fontsize{14}{14}\selectfont\color{popink}LEÇON #5}\par\vspace{4pt}
    {\raggedright\hyphenpenalty=10000\exhyphenpenalty=10000\popdisplay\fontsize{25}{27}\selectfont\color{popcream}\MakeUppercase{#1}\par}
    \vspace{8pt}
    {\popxbold\fontsize{9.5}{9.5}\selectfont\color{popink}\MakeUppercase{Durée conseillée : #6}}
  \end{tcolorbox}
  \begin{tcolorbox}[enhanced,colback=white,colframe=popink,boxrule=2.2pt,arc=10pt,
      drop shadow=popink,left=16pt,right=16pt,top=10pt,bottom=10pt,after skip=8pt]
    \poppill{Dans cette leçon}{poppurple}{white}\par\vspace{5pt}
    \popsemi\fontsize{9.3}{12.6}\selectfont\color{popink}#7
  \end{tcolorbox}
  \noindent\begin{minipage}[t]{0.487\linewidth}
    \begin{tcolorbox}[enhanced,colback=popgreen,colframe=popink,boxrule=2.2pt,arc=10pt,
        drop shadow=popink,left=11pt,right=11pt,top=10pt,bottom=10pt,height=3.1cm,valign=center]
      \tcbox[colback=white,colframe=popink,boxrule=1.3pt,arc=5pt,boxsep=0pt,nobeforeafter,
        left=3pt,right=3pt,top=3pt,bottom=3pt,valign=center]{\qrcode[height=1.9cm]{https://play.google.com/store/apps/details?id=com.luvvix.onbuch&hl=fr}}%
      \hspace{8pt}\begin{minipage}[c]{0.5\linewidth}
        {\popdisplay\fontsize{11}{12.5}\selectfont\color{white}\MakeUppercase{Télécharge OnBuch}}\par\vspace{4pt}
        {\raggedright\popsemi\fontsize{8.4}{11}\selectfont\color{white}Gratuit \textbullet\ Google Play\\Tuteur IA, annales, concours, résultats\par}
      \end{minipage}
    \end{tcolorbox}
  \end{minipage}\hfill
  \begin{minipage}[t]{0.487\linewidth}
    \begin{tcolorbox}[enhanced,colback=poppurple,colframe=popink,boxrule=2.2pt,arc=10pt,
        drop shadow=popink,left=11pt,right=11pt,top=10pt,bottom=10pt,height=3.1cm,valign=center]
      \tikz[baseline=-0.7ex]\node[circle,fill=white,draw=popink,line width=1.3pt,minimum size=1.6cm,inner sep=0]{\popdisplay\fontsize{24}{24}\selectfont\color{poppurple}+};%
      \hspace{8pt}\begin{minipage}[c]{0.6\linewidth}
        {\popdisplay\fontsize{11}{12.5}\selectfont\color{white}\MakeUppercase{Passe à OnBuch+}}\par\vspace{4pt}
        {\raggedright\popsemi\fontsize{8.4}{11}\selectfont\color{white}Tous les cours signés, Tuteur IA illimité, fiches PDF — onglet OnBuch+ de l'app\par}
      \end{minipage}
    \end{tcolorbox}
  \end{minipage}
  \vspace{8pt}
  \popcontenu
  \vfill
  \signatureonbuch
  \restoregeometry
  \newpage
}

% Rangée « Ce cours contient » (page de garde)
\newcommand{\poptile}[3]{% #1 couleur #2 gros texte #3 légende
  \begin{tcolorbox}[enhanced,colback=#1,colframe=popink,boxrule=2pt,arc=9pt,shadow={2.5pt}{-2.5pt}{0pt}{popink},
      left=6pt,right=6pt,top=9pt,bottom=9pt,halign=center,valign=center,height=1.75cm,width=\linewidth,nobeforeafter]
    {\popdisplay\fontsize{12.5}{13}\selectfont\color{white}\MakeUppercase{#2}}\par\vspace{3pt}
    {\centering\popsemi\fontsize{7.8}{9.5}\selectfont\color{white}#3\par}
  \end{tcolorbox}}
\newcommand{\popcontenu}{%
  \par\noindent{\popxboldls\fontsize{7.5}{7.5}\selectfont\color{popmuted}CE COURS SIGNÉ CONTIENT}\par\vspace{7pt}
  \noindent\begin{minipage}[t]{0.235\linewidth}\poptile{popblue}{Cours}{complet, illustré, pas à pas}\end{minipage}\hfill
  \begin{minipage}[t]{0.235\linewidth}\poptile{poporange}{Méthodes}{et exercices résolus}\end{minipage}\hfill
  \begin{minipage}[t]{0.235\linewidth}\poptile{poppink}{Exercices}{type examen + corrigés}\end{minipage}\hfill
  \begin{minipage}[t]{0.235\linewidth}\poptile{popgreen}{Fiche bilan}{l'essentiel à retenir}\end{minipage}\par}

% Sommaire
\newtcolorbox{sommaire}{enhanced,colback=white,colframe=popink,boxrule=2.2pt,arc=10pt,
  drop shadow=popink,left=18pt,right=18pt,top=13pt,bottom=13pt,before skip=6pt,after skip=20pt,
  before upper={\poppill{Sommaire}{poppurple}{white}\par\vspace{9pt}\popsemi\fontsize{10.6}{14.5}\selectfont\color{popink}}}

% Signature de fin de cours — poussée tout en bas de la dernière page
\newcommand{\findecours}{%
  \par\vfill\nopagebreak
  \begin{center}\popsquiggle{poporange}\end{center}\vspace{4pt}
  \signatureonbuch
}
\AtBeginDocument{\def\llbracket{\mathopen{[\mkern-3.5mu[}}\def\rrbracket{\mathclose{]\mkern-3.5mu]}}} % intervalles d'entiers (glyphe ⟦ absent de la police)
% Glyphes absents de Fira Math : redéfinis à partir de glyphes disponibles
\AtBeginDocument{%
  \def\setminus{\mathbin{\backslash}}\let\smallsetminus\setminus
  \def\vdots{\mathord{\vbox{\baselineskip3.2pt\lineskiplimit0pt\kern4pt\hbox{.}\hbox{.}\hbox{.}}}}%
  \def\ddots{\mathinner{\mkern1mu\raise6.5pt\hbox{.}\mkern2mu\raise3.6pt\hbox{.}\mkern2mu\raise0.7pt\hbox{.}\mkern1mu}}%
  \def\triangle{\mathord{\scalebox{0.9}{$\symup{\Delta}$}}}%
  \def\nabla{\mathord{\raisebox{\depth}{\scalebox{1}[-1]{$\symup{\Delta}$}}}}%
  \def\checkmark{\popcheck{popdark}}%
  \def\star{\mathbin{\ast}}\def\bigstar{\mathord{\ast}}%
  \def\diamond{\mathbin{\scalebox{0.6}{$\Diamond$}}}%
  \def\bigcup{\mathop{\mathchoice{\raisebox{-0.3em}{\scalebox{1.7}{$\cup$}}}{\scalebox{1.25}{$\cup$}}{\cup}{\cup}}\displaylimits}%
  \def\bigcap{\mathop{\mathchoice{\raisebox{-0.3em}{\scalebox{1.7}{$\cap$}}}{\scalebox{1.25}{$\cap$}}{\cap}{\cap}}\displaylimits}%
  \def\lesseqqgtr{\mathrel{\lessgtr}}\def\bigoplus{\mathop{\oplus}\displaylimits}%
}
\providecommand{\ssi}{si et seulement si} % abréviation fréquente
\AtBeginDocument{\providecommand{\wideparen}{\overparen}} % arc de cercle

%% FIGFIX-BEGIN v3 — correctifs globaux des figures (docs/onbuch-generation/tools/figfix)
\usepackage{adjustbox}\usepackage{etoolbox}
% toute figure TikZ/pgfplots plus large que la ligne est réduite à la largeur disponible
\BeforeBeginEnvironment{tikzpicture}{\begin{adjustbox}{max width=\linewidth}}
\AfterEndEnvironment{tikzpicture}{\end{adjustbox}}
% légendes pgfplots lisibles par-dessus les courbes
\pgfplotsset{every axis/.append style={legend style={fill=white,fill opacity=0.92,text opacity=1,draw=black!15,inner sep=3pt}}}
% étiquettes d'arêtes (syntaxe edge["..."]) avec fond blanc
\tikzset{every edge quotes/.append style={fill=white,inner sep=1.2pt}}
% bulles de cartes mentales : texte à la ligne dans la bulle, bulle un peu plus grande
\tikzset{every concept/.append style={text width=2.0cm,align=center,minimum size=2.3cm,inner sep=2pt}}
%% FIGFIX-END
\begin{document}

\pagegarde{Integration activities / Evaluation / Remediation (Chapter 1)}{Anglais}{1ère A}{Chapter 1 : Family and Social Life}{ENA13}{2 h}{Cette leçon de révision clôt le chapitre « Family and Social Life » : elle synthétise le vocabulaire et la grammaire étudiés, propose des exercices progressifs de type examen et guide chaque élève dans l'auto-évaluation et la remédiation de ses difficultés.
\vspace{4pt}
\begin{multicols}{2}\small
\begin{itemize}
\item À la fin de cette leçon, je sais mobiliser le vocabulaire de la famille et de la vie sociale dans une tâche d'intégration
\item À la fin de cette leçon, je sais réutiliser correctement les points de grammaire du chapitre (present perfect, comparatifs et superlatifs, possessifs)
\item À la fin de cette leçon, je sais comprendre un texte sur les relations familiales et répondre à des questions de type examen
\item À la fin de cette leçon, je sais rédiger un paragraphe cohérent sur ma famille ou une cérémonie sociale
\item À la fin de cette leçon, je sais m'auto-évaluer à l'aide d'une grille et identifier mes points faibles
\item À la fin de cette leçon, je sais remédier à mes erreurs grâce à des exercices ciblés
\end{itemize}\end{multicols}}

\begin{sommaire}
\begin{multicols}{2}
\begin{enumerate}
\item Warm-up : diagnostic rapide
\item Vocabulary synthesis
\item Grammar synthesis
\item Reading : integration text
\item Progressive exam-type practice
\item Self-evaluation and remediation
\item Activité d'intégration
\item Méthodes et exercices résolus
\item Exercices
\item Corrigés des exercices
\item Fiche bilan
\end{enumerate}
\end{multicols}
\end{sommaire}

\begin{objectifs}
\begin{itemize}
\item À la fin de cette leçon, je sais mobiliser le vocabulaire de la famille et de la vie sociale dans une tâche d'intégration
\item À la fin de cette leçon, je sais réutiliser correctement les points de grammaire du chapitre (present perfect, comparatifs et superlatifs, possessifs)
\item À la fin de cette leçon, je sais comprendre un texte sur les relations familiales et répondre à des questions de type examen
\item À la fin de cette leçon, je sais rédiger un paragraphe cohérent sur ma famille ou une cérémonie sociale
\item À la fin de cette leçon, je sais m'auto-évaluer à l'aide d'une grille et identifier mes points faibles
\item À la fin de cette leçon, je sais remédier à mes erreurs grâce à des exercices ciblés
\item À la fin de cette leçon, je sais gérer mon temps et mes stratégies face à une épreuve de type Bac
\end{itemize}
\end{objectifs}

\begin{prerequis}
\begin{itemize}
\item Vocabulaire de la famille et des liens de parenté (Seconde et début du chapitre 1)
\item Present perfect simple et ses marqueurs (already, just, never, since, for)
\item Comparatifs et superlatifs des adjectifs
\item Adjectifs et pronoms possessifs
\item Connecteurs simples de la narration et de l'opinion (first, then, however, in my opinion)
\end{itemize}
\end{prerequis}

\newpage

\coursec{Warm-up : diagnostic rapide}

\begin{textebox}[Manka's challenge]
The end-of-term exam is in two weeks. Manka, a Form 5 student in Bamenda, opens her copybook and realises she has studied many things about family and social life: describing relatives, talking about marriage customs, using the present perfect and comparatives... but can she use all of them together in one real task? Today, like Manka, you will bring everything from Chapter 1 together, test yourself, and fix your weak points before the big day.
\end{textebox}

\textbf{Glossaire du texte.} \emph{copybook} (nom) : cahier ; \emph{to realise} (verbe) : se rendre compte ; \emph{customs} (nom pluriel) : coutumes ; \emph{weak points} (nom pluriel) : points faibles ; \emph{the big day} (expression) : le grand jour (ici, le jour de l'examen).

\textbf{Problématique de la séance.} Manka a appris beaucoup de choses séparément : du vocabulaire, des règles de grammaire, des expressions. Mais un examen ne teste pas les notions une par une : il demande de \cle{mobiliser tout ensemble}. La question qui guide toute cette leçon est donc : \emph{How can I check what I really know, and how can I repair what is broken?} (Comment vérifier ce que je sais vraiment, et comment réparer ce qui est cassé ?)

\courssub{Le contrat de la séance : réviser, s'évaluer, remédier}

Avant de commencer, fixons les règles du jeu. Cette séance de révision suit trois étapes, comme un mécanicien qui prépare une voiture avant un long voyage :

\begin{enumerate}
\item \cle{Réviser} (\eng{revise}) : rassembler en un seul regard tout le chapitre 1.
\item \cle{S'évaluer} (\eng{evaluate yourself}) : se tester \textbf{sans regarder le cours}, pour découvrir honnêtement ses forces et ses faiblesses.
\item \cle{Remédier} (\eng{remediate}) : retravailler \textbf{uniquement} ce qui a échoué, jusqu'à ce que cela devienne solide.
\end{enumerate}

\begin{methode}[Réviser efficacement : la méthode en trois temps]
Beaucoup d'élèves « révisent » en relisant leur cahier pendant des heures. C'est la méthode la moins efficace : relire donne une \emph{illusion} de connaissance. Voici la méthode des bons élèves :
\begin{enumerate}
\item \textbf{Se tester d'abord, sans le cours.} Ferme ton cahier et essaie de répondre à une question ou de faire un exercice. C'est l'effort de rappel qui grave la notion dans la mémoire.
\item \textbf{Vérifier ensuite.} Ouvre le cours et compare : qu'est-ce qui était juste ? Qu'est-ce qui était faux ou absent ?
\item \textbf{Retravailler seulement ce qui a échoué.} Ne perds pas de temps à relire ce que tu maîtrises déjà. Concentre toute ton énergie sur tes points faibles, puis re-teste-toi le lendemain.
\end{enumerate}
Toute cette leçon applique cette méthode : le warm-up est ton premier test « à cahier fermé ».
\end{methode}

\courssub{Quiz oral de diagnostic : cinq questions}

Fermez vos cahiers. L'enseignant pose cinq questions ; chaque élève répond à voix haute ou sur une ardoise. L'objectif n'est pas d'avoir tout juste : c'est de \cle{découvrir} ce que vous savez encore.

\begin{enumerate}
\item \eng{What is the English word for “beau-père” (le père de ton mari ou de ta femme)?}
\item \eng{Complete: “I have known my best friend \ldots\ 2019.” (since ou for ?)}
\item \eng{What is the comparative of “good”?}
\item \eng{Name two traditional ceremonies in your community.}
\item \eng{Correct this sentence: “My parents-in-laws lives in Douala.”}
\end{enumerate}

\begin{exoresolu}[Quiz de diagnostic : corrigé commenté]
Réponds aux cinq questions du quiz sans regarder ton cours, puis vérifie.
\tcblower
\begin{enumerate}
\item \eng{father-in-law}. Pluriel : \eng{fathers-in-law}. Les mots composés de la belle-famille prennent \eng{-in-law} : \eng{mother-in-law} (belle-mère), \eng{brother-in-law} (beau-frère), \eng{sister-in-law} (belle-sœur).
\item \eng{since 2019}. On utilise \cle{since} + un point de départ (\eng{since 2019, since Monday, since I was ten}) et \cle{for} + une durée (\eng{for five years, for a long time}).
\item \eng{better}. Le comparatif de \eng{good} est irrégulier : \eng{good → better → the best}. Jamais \emph{gooder} ni \emph{more good}.
\item Réponses possibles : \eng{a traditional wedding, a funeral ceremony, a birth celebration, a traditional rite, a bride-price ceremony}.
\item \eng{My parents-in-law live in Douala.} Deux erreurs : le pluriel se met sur \eng{parents} (\eng{parents-in-law}, pas \emph{parents-in-laws}), et le sujet est pluriel, donc \eng{live} sans \eng{-s}.
\end{enumerate}
\end{exoresolu}

\courssub{Brainstorming : la carte mentale du chapitre}

Maintenant, ouvrons la mémoire collective. Question au tableau : \eng{What do you remember about Chapter 1?} Chaque élève propose un mot, une règle ou un exemple, et la classe construit ensemble la carte mentale suivante. Le chapitre 1 repose sur \cle{trois piliers} : les membres de la famille (\eng{family members}), les événements sociaux (\eng{social events} : mariages, funérailles, rites traditionnels) et les outils de grammaire (\eng{grammar tools}) qui permettent d'en parler.

\begin{popfigure}
\begin{center}
\begin{tikzpicture}[mindmap, grow cyclic,
  every node/.style={concept, text=white, align=center},
  root concept/.append style={fill=popblue, font=\bfseries},
  level 1/.append style={level distance=3.2cm, sibling angle=120},
  level 2/.append style={level distance=2.4cm, sibling angle=45, font=\small}]
\node[root concept]{Family and Social Life}
  child[concept color=poporange]{ node {Family members}
    child { node[fill=poporangeL, text=popdark] {parents, uncle, aunt, cousin} }
    child { node[fill=poporangeL, text=popdark] {in-laws, step-family} }
  }
  child[concept color=popgreen]{ node {Social events}
    child { node[fill=popgreenL, text=popdark] {weddings, funerals} }
    child { node[fill=popgreenL, text=popdark] {traditional rites, feelings} }
  }
  child[concept color=poppurple]{ node {Grammar tools}
    child { node[fill=poppurpleL, text=popdark] {present perfect} }
    child { node[fill=poppurpleL, text=popdark] {comparatives and superlatives} }
  };
\end{tikzpicture}
\end{center}
\legende{Carte mentale de synthèse : les trois piliers du chapitre 1.}
\end{popfigure}

\courssub{Rappel express des trois piliers}

\textbf{Pilier 1 : family members.} Le vocabulaire de la famille se divise en trois cercles : la famille proche (\eng{nuclear family} : \eng{father, mother, brother, sister, husband, wife, son, daughter}), la famille élargie (\eng{extended family} : \eng{grandfather, grandmother, uncle, aunt, nephew, niece, cousin}) et la famille par alliance (\eng{in-laws}) ou recomposée (\eng{stepfather, stepmother, half-brother}).

\textbf{Pilier 2 : social events.} Pour décrire la vie sociale : \eng{a wedding} (un mariage), \eng{a funeral} (des funérailles), \eng{a traditional rite} (un rite traditionnel), \eng{a bride price / dowry} (la dot), \eng{a celebration} (une fête), \eng{to mourn} (porter le deuil), \eng{to get married} (se marier), \eng{to celebrate} (célébrer).

\textbf{Pilier 3 : grammar tools.} Deux outils dominent le chapitre :
\begin{itemize}
\item Le \cle{present perfect} : \eng{have/has + participe passé}, pour relier le passé au présent. \eng{She has lived in Buea since 2015.} « Elle vit à Buea depuis 2015. »
\item Les \cle{comparatifs et superlatifs} : \eng{taller than}, \eng{more beautiful than}, \eng{the tallest}, \eng{the most beautiful}, et les irréguliers \eng{good → better → the best}, \eng{bad → worse → the worst}.
\end{itemize}

\begin{center}
\begin{tabularx}{\linewidth}{|X|X|X|}
\hline
\rowcolor{popblueL} Français & English & Remarque \\
\hline
les parents (père et mère) & parents & toujours au pluriel dans ce sens \\
\hline
les parents (famille élargie) & relatives / relations & \eng{my relatives in the village} \\
\hline
le beau-père / la belle-mère & father-in-law / mother-in-law & pluriel : \eng{fathers-in-law} \\
\hline
se marier avec quelqu'un & to marry somebody / to get married to somebody & jamais \emph{to marry with} \\
\hline
la dot & the bride price / dowry & coutume très présente au Cameroun \\
\hline
porter le deuil & to mourn / to be in mourning & \eng{The family is mourning.} \\
\hline
depuis (point de départ) & since & \eng{since 2020} \\
\hline
depuis (durée) & for & \eng{for three years} \\
\hline
\end{tabularx}
\end{center}

\begin{attention}[« Parents » et « relatives » : le faux ami du chapitre]
En français, « mes parents » peut désigner le père et la mère \emph{ou} toute la famille. En anglais, c'est impossible :
\begin{itemize}
\item \eng{parents} = uniquement le père et la mère. \eng{My parents live in Yaoundé.} « Mes parents (père et mère) habitent à Yaoundé. »
\item \eng{relatives} ou \eng{relations} = toute la famille élargie. \eng{All my relatives came to the wedding.} « Tous mes parents (oncles, tantes, cousins...) sont venus au mariage. »
\end{itemize}
Dire \eng{my parents} en pensant à toute la famille est une erreur classique du francophone : l'anglophone comprendra seulement « mon père et ma mère ».
\end{attention}

\begin{attention}[Deux autres pièges à revérifier avant l'examen]
\begin{itemize}
\item \eng{to marry} se construit \textbf{sans préposition} : \eng{She married a teacher.} « Elle a épousé un enseignant. » Avec \eng{get married}, on ajoute \eng{to} : \eng{She got married to a teacher.} Jamais \emph{married with}.
\item Au present perfect, attention aux participes passés irréguliers : \eng{be → been}, \eng{go → gone}, \eng{see → seen}, \eng{write → written}, \eng{take → taken}. \eng{I have never seen a traditional wedding.} « Je n'ai jamais vu de mariage traditionnel. »
\end{itemize}
\end{attention}

\begin{savaistu}[« My people » : l'anglais réel du Cameroun anglophone]
Dans les régions du Nord-Ouest et du Sud-Ouest, quand un élève dit \eng{“I am going to see my people”}, il ne parle pas seulement de son père et de sa mère : \eng{my people} désigne toute la famille élargie, le village, parfois même la communauté. De même, \eng{“my father”} peut désigner un oncle paternel, et \eng{“my brother”} un cousin. Cet usage reflète la conception africaine de la famille, bien plus large que la famille nucléaire européenne. L'anglais réel dépasse le manuel : une langue vit et s'adapte à la culture de ceux qui la parlent.
\end{savaistu}

\courssub{Ta fiche personnelle de diagnostic}

Dernière étape du warm-up : chacun fait le point \emph{pour lui-même}. Sur une petite fiche (ou en marge de ton cahier), recopie ce tableau et complète-le honnêtement :

\begin{center}
\begin{tabularx}{\linewidth}{|X|X|}
\hline
\rowcolor{popblueL} My two weakest points & Why I am not sure \\
\hline
1. \ldots & \ldots \\
\hline
2. \ldots & \ldots \\
\hline
\end{tabularx}
\end{center}

Exemple de fiche remplie par Manka : \eng{1. since / for — I always mix them up. 2. The plural of mother-in-law.} « 1. since / for — je les confonds toujours. 2. Le pluriel de mother-in-law. »

Garde précieusement cette fiche : c'est ta \cle{feuille de route} pour toute la séance. À la fin de la leçon, tu la reprendras pour vérifier que tes deux points faibles ont été réparés. C'est exactement cela, remédier : viser juste, travailler peu, mais travailler là où c'est nécessaire.

\begin{experience}[Diagnostic express en binômes]
Par groupes de deux, en cinq minutes :
\begin{enumerate}
\item Élève A pose les cinq questions du quiz à élève B, qui répond sans cahier ; puis on inverse les rôles.
\item Chacun annonce à son partenaire ses deux points faibles en anglais : \eng{“I am not sure about \ldots”} ou \eng{“I always forget \ldots”}.
\item Le partenaire donne un conseil : \eng{“You should practise \ldots”} ou \eng{“Look at the table again.”}
\end{enumerate}
L'enseignant circule et note les difficultés les plus fréquentes de la classe pour orienter la remédiation de la fin de séance.
\end{experience}

\begin{aretenir}[L'essentiel du warm-up]
\begin{itemize}
\item Réviser, ce n'est pas relire : c'est \cle{se tester}, vérifier, puis retravailler ses points faibles.
\item Le chapitre 1 repose sur trois piliers : \eng{family members}, \eng{social events}, \eng{grammar tools} (present perfect, comparatifs et superlatifs).
\item Trois pièges à éliminer : \eng{parents} contre \eng{relatives} ; \eng{since} + point de départ contre \eng{for} + durée ; \eng{to marry somebody} sans préposition.
\end{itemize}
\end{aretenir}

\coursec{Vocabulary synthesis}

Le diagnostic rapide du warm-up vous a permis de repérer vos points forts et vos lacunes. Nous allons maintenant consolider méthodiquement tout le vocabulaire du chapitre \eng{Family and Social Life}. L'objectif n'est pas de réciter des listes, mais de savoir \cle{employer} chaque mot dans une phrase correcte, avec la bonne préposition, la bonne collocation et le bon registre. Nous procéderons en trois temps : observation d'un texte, synthèse des champs lexicaux, puis entraînement ciblé sur les pièges des francophones.

\courssub{Observation : un texte pour réactiver le lexique}

Lisez d'abord ce texte. Soulignez mentalement tous les mots qui appartiennent au thème de la famille et de la vie sociale.

\begin{textebox}[Family and Social Life in Cameroon — reading text]
In many Cameroonian communities, the extended family plays a central role: uncles, aunts and grandparents all help to raise the children. When a young man wants to marry, his relatives meet the girl's family to discuss the engagement and to negotiate the bride price. Last December, my uncle paid the bride price for his son's wedding, and the whole village attended the ceremony. The chief and the elders blessed the couple according to the traditional rites of the community.

Social life does not stop at happy events. When someone dies, neighbours and relatives gather for the wake, and hundreds of people attend the funeral to mourn the dead and to support the grieving family. A few weeks later, a naming ceremony may be organised for a newborn baby, and once again the community comes together.

Relations inside a family are not always easy. Brothers and sisters sometimes quarrel, and a stepmother may find it difficult to get on with her stepchildren. But in most families, people respect their elders, take care of the sick and rely on one another in difficult times. Even when young people move to Yaoundé or Douala for work, they keep in touch with their parents and send money home. An orphan, for example, is rarely left alone: a guardian, usually an uncle or an aunt, takes care of the child. This solidarity is the strength of the extended family.
\end{textebox}

\textbf{Glossaire.} \emph{to raise} (v.) : élever ; \emph{bride price} (n.) : dot ; \emph{wake} (n.) : veillée funèbre ; \emph{to mourn} (v.) : pleurer (un défunt), être en deuil ; \emph{grieving} (adj.) : endeuillé ; \emph{newborn} (adj./n.) : nouveau-né ; \emph{stepchildren} (n.) : beaux-enfants ; \emph{guardian} (n.) : tuteur, tutrice.

\textbf{Questions de compréhension.}
\begin{enumerate}
\item Who helps to raise children in many Cameroonian communities?
\item What did the narrator's uncle pay last December, and why?
\item Why do people attend a funeral?
\item What happens at a naming ceremony?
\item Who usually takes care of an orphan?
\end{enumerate}

\courssub{Champ lexical 1 : the family}

\begin{definition}[Extended family]
\eng{Extended family} : \emph{all the people in a family beyond parents and children (grandparents, uncles, aunts, cousins)} — la famille élargie, si importante en Afrique. On l'oppose à la \eng{nuclear family} (famille nucléaire : père, mère, enfants).
\end{definition}

\begin{center}
\begin{tabularx}{\linewidth}{|X|X|X|}\hline
\rowcolor{popblueL} English & Français & Remarque \\ \hline
nuclear family & famille nucléaire & parents + children \\ \hline
extended family & famille élargie & grandparents, uncles, aunts, cousins \\ \hline
single parent & parent isolé (seul) & a single mother / a single father \\ \hline
stepmother / stepfather & belle-mère / beau-père (remariage) & préfixe \eng{step-} = lien par remariage \\ \hline
sibling & frère ou sœur & mot neutre : \eng{How many siblings do you have?} \\ \hline
spouse & époux / épouse & registre soutenu, administratif \\ \hline
in-laws & belle-famille & \eng{my mother-in-law} = ma belle-mère \\ \hline
orphan & orphelin(e) & an orphan of war \\ \hline
guardian & tuteur, tutrice & personne qui s'occupe d'un orphelin \\ \hline
\end{tabularx}
\end{center}

\begin{exemplebox}[Exemples traduits]
\eng{My uncle paid the bride price last December.} « Mon oncle a payé la dot en décembre dernier. »

\eng{She gets on well with her in-laws.} « Elle s'entend bien avec sa belle-famille. »

\eng{After the accident, the orphans were raised by a guardian.} « Après l'accident, les orphelins ont été élevés par un tuteur. »

\eng{He has three siblings: two brothers and a sister.} « Il a trois frères et sœurs : deux frères et une sœur. »
\end{exemplebox}

\courssub{Champ lexical 2 : social events and community}

\begin{center}
\begin{tabularx}{\linewidth}{|X|X|X|}\hline
\rowcolor{popblueL} English & Français & Remarque \\ \hline
wedding & mariage (cérémonie) & \eng{to attend a wedding} \\ \hline
engagement & fiançailles & \eng{to get engaged} = se fiancer \\ \hline
dowry / bride price & dot & \eng{to pay the bride price} \\ \hline
funeral & funérailles, enterrement & TOUJOURS singulier pour une cérémonie \\ \hline
wake & veillée funèbre & avant l'enterrement \\ \hline
naming ceremony & cérémonie de nom (baptême traditionnel) & pour un nouveau-né \\ \hline
traditional rite & rite traditionnel & \eng{to perform a rite} \\ \hline
chief & chef (de village) & \eng{the chief and the elders} \\ \hline
elder & ancien, aîné & respect des anciens \\ \hline
community & communauté & \eng{the whole community} \\ \hline
\end{tabularx}
\end{center}

\begin{attention}[Funeral : singulier et verbe correct]
\eng{Funeral} est toujours \cle{singulier} quand on parle d'une cérémonie : \eng{The funeral took place on Saturday.} On dit \eng{to attend a funeral}, JAMAIS \eng{to assist a funeral}. Le verbe anglais \eng{to assist} signifie « aider » (\eng{The nurse assisted the doctor} = l'infirmière a aidé le médecin). Pour « assister à » une cérémonie, utilisez toujours \eng{to attend} ou \eng{to be present at}.
\end{attention}

\courssub{Champ lexical 3 : feelings and relationships}

Ces verbes et expressions décrivent les relations humaines. Retenez chaque verbe \cle{avec sa préposition}.

\begin{center}
\begin{tabularx}{\linewidth}{|X|X|X|}\hline
\rowcolor{popblueL} English & Français & Exemple \\ \hline
to get on with & s'entendre avec & She gets on with her stepfather. \\ \hline
to quarrel (with) & se disputer (avec) & The brothers quarrel about money. \\ \hline
to support & soutenir & Friends support the grieving family. \\ \hline
to respect & respecter & Children respect their elders. \\ \hline
to be close to & être proche de & I am very close to my grandmother. \\ \hline
to take care of & s'occuper de & She takes care of her sick mother. \\ \hline
to rely on & compter sur & You can rely on your family. \\ \hline
\end{tabularx}
\end{center}

\begin{propriete}[Collocations clés à mémoriser comme des blocs]
Une \cle{collocation} est une combinaison de mots qui va naturellement ensemble en anglais ; on ne la traduit pas mot à mot.
\begin{itemize}
\item \eng{to raise children} = élever des enfants (jamais \eng{grow children})
\item \eng{to pay the bride price} = payer la dot
\item \eng{to attend a ceremony} = assister à une cérémonie
\item \eng{to mourn the dead} = pleurer les morts, honorer le deuil
\item \eng{to keep in touch (with)} = rester en contact (avec)
\end{itemize}
\end{propriete}

\begin{popfigure}
\begin{center}
\begin{tabular}{|>{\centering}p{0.3\linewidth}|>{\centering}p{0.3\linewidth}|>{\centering\arraybackslash}p{0.3\linewidth}|}
\hline
\rowcolor{popblueL} \textbf{\textcolor{popdark}{Family}} & \textbf{\textcolor{popdark}{Social events}} & \textbf{\textcolor{popdark}{Feelings}} \\
\hline
nuclear family & wedding & to get on with \\
\hline
sibling & engagement & to quarrel \\
\hline
spouse & funeral & to support \\
\hline
in-laws & naming ceremony & to be close to \\
\hline
orphan & traditional rite & to rely on \\
\hline
\end{tabular}
\end{center}
\legende{Synthèse visuelle : trois champs lexicaux du chapitre, cinq mots-clés par colonne.}
\end{popfigure}

\courssub{Word families : construire le vocabulaire}

Un même radical donne plusieurs mots de natures différentes. Connaître la famille d'un mot, c'est multiplier son vocabulaire.

\begin{center}
\begin{tabular}{|l|l|l|}\hline
\rowcolor{popblueL} Verbe & Nom & Adjectif / Autre \\ \hline
to marry (épouser) & marriage (mariage) & married (marié) \\ \hline
to die (mourir) & death (mort, décès) & dead (mort) \\ \hline
to mourn (pleurer un mort) & mourning (deuil) & mournful (triste, en deuil) \\ \hline
to relate (raconter, relier) \textbackslash{} be related (être parent) & relation / relationship (relation) & relative (parent, membre de la famille) \\ \hline
\end{tabular}
\end{center}

\begin{exemplebox}[La même famille dans trois phrases]
\eng{They will marry in June. Their marriage will take place in Buea. They are happily married.} « Ils se marieront en juin. Leur mariage aura lieu à Buea. Ils sont heureux en ménage. »

\eng{Her grandfather died last year. His death saddened the village. The family still mourns the dead.} « Son grand-père est mort l'an dernier. Son décès a attristé le village. La famille pleure encore le défunt. »

\eng{All my relatives came to the wedding. Family relations are important here.} « Tous mes parents (élargis) sont venus au mariage. Les relations familiales sont importantes ici. »
\end{exemplebox}

\courssub{Faux amis : les pièges des francophones}

\begin{attention}[Trois faux amis à éliminer définitivement]
\begin{itemize}
\item \eng{to assist} $\neq$ « assister à ». \eng{to assist} = aider. « Assister à une cérémonie » = \eng{to attend a ceremony}.
\item \eng{sensible} $\neq$ « sensible » (émotif). \eng{sensible} = raisonnable, sensé. « Sensible » (émotif) = \eng{sensitive}. \eng{She is very sensitive} = elle est très sensible.
\item \eng{parents} $\neq$ « parents » au sens large (famille élargie). \eng{parents} = le père et la mère uniquement. Les « parents éloignés » = \eng{relatives} ou \eng{relations}. \eng{My relatives live in Bamenda} = mes parents (famille élargie) vivent à Bamenda.
\end{itemize}
\end{attention}

\begin{exoresolu}[Exercice de synthèse : complétez avec le mot correct]
Énoncé. Complétez chaque phrase avec un mot du chapitre : \eng{guardian — bride price — funeral — siblings — sensitive — rely on}.
\begin{enumerate}
\item My uncle negotiated the \dots\ before the wedding.
\item Hundreds of people attended the \dots\ to mourn the old chief.
\item After her parents' death, her aunt became her \dots.
\item I have two \dots : one brother and one sister.
\item Don't cry so easily, you are too \dots !
\item In difficult times, you can always \dots\ your family.
\end{enumerate}
\tcblower
Solution détaillée.
\begin{enumerate}
\item \eng{bride price} — la dot se négocie avant le mariage (\eng{to negotiate / to pay the bride price}).
\item \eng{funeral} — toujours au singulier pour une cérémonie, avec le verbe \eng{to attend}.
\item \eng{guardian} — le tuteur s'occupe d'un orphelin (\eng{to take care of an orphan}).
\item \eng{siblings} — mot neutre qui regroupe frères et sœurs.
\item \eng{sensitive} — faux ami : « sensible » (émotif) se dit \eng{sensitive}, pas \eng{sensible}.
\item \eng{rely on} — « compter sur » exige la préposition \eng{on}.
\end{enumerate}
\end{exoresolu}

\begin{methode}[Comment mémoriser durablement ce lexique]
\begin{enumerate}
\item Apprenez chaque mot \cle{dans une phrase complète}, jamais isolément.
\item Apprenez le mot \cle{avec sa collocation} : \eng{to attend a funeral}, \eng{to raise children}.
\item Regroupez par \cle{famille de mots} : \eng{marry / marriage / married}.
\item Testez-vous en traduisant du français vers l'anglais, puis vérifiez la préposition.
\item Méfiez-vous des mots qui ressemblent au français : ils cachent souvent un faux ami.
\end{enumerate}
\end{methode}

\begin{exercice}[Pour vous entraîner]
Traduisez en anglais : 1) Ma belle-mère s'entend bien avec mes parents. 2) Toute la communauté a assisté au mariage. 3) Les anciens respectent les rites traditionnels. 4) Il reste en contact avec sa famille élargie. 5) Elle est raisonnable mais très sensible.
\end{exercice}

\begin{corrige}[Exercice « Pour vous entraîner »]
1) \eng{My mother-in-law gets on well with my parents.} 2) \eng{The whole community attended the wedding.} 3) \eng{The elders respect traditional rites.} 4) \eng{He keeps in touch with his extended family.} 5) \eng{She is sensible but very sensitive.}
\end{corrige}

Votre vocabulaire est désormais consolidé. La section suivante, \eng{Grammar synthesis}, reprendra de la même manière les points de grammaire du chapitre.

\coursec{Grammar synthesis}

Après avoir révisé le vocabulaire de la famille et de la vie sociale (\emph{Vocabulary synthesis}), il est temps de consolider les outils grammaticaux qui permettent de parler de nos expériences familiales, de comparer les membres d'une famille et d'exprimer la possession. Ce chapitre a mobilisé quatre grands points : le \cle{present perfect}, les \cle{comparatifs} et \cle{superlatifs}, et les \cle{possessifs}. Reprenons-les méthodiquement, en partant d'observations concrètes.

\courssub{The present perfect : an experience connected to the present}

\textbf{Observation}\par

Lisons d'abord ce court texte, puis observons les verbes en gras.

\begin{textebox}[A family reunion in Buea]
My family has always been very close. We have lived in Buea since 2010, and I have never regretted leaving the big city. My elder sister has just got married, and the whole family has already met her husband, a kind teacher from Bamenda. My parents have known each other for thirty-five years, and they have never had a serious argument. Have you ever attended a traditional wedding in the South-West Region? I have been to three this year, and each one has taught me something new about our customs. Last Saturday, my uncle arrived from Douala with gifts for everyone. He has not visited us yet this month, so the children were very excited.
\end{textebox}

\begin{definition}[Glossaire]
\begin{itemize}
\item \eng{elder sister} (n.) : sœur aînée
\item \eng{to regret} (v.) : regretter
\item \eng{an argument} (n.) : une dispute
\item \eng{to attend} (v.) : assister à
\item \eng{customs} (n.) : coutumes, traditions
\item \eng{gifts} (n.) : cadeaux
\end{itemize}
\end{definition}

Remarquez : toutes les actions en gras sont soit des \cle{expériences} sans date précise, soit des situations qui \cle{ont commencé dans le passé et continuent aujourd'hui}. C'est exactement le domaine du present perfect.

\begin{propriete}[Forme du present perfect]
Le present perfect se construit avec l'auxiliaire \cle{have} (ou \cle{has} à la 3\up{e} personne du singulier) suivi du \cle{participe passé} du verbe.
\begin{center}
\begin{tabular}{|l|l|l|}
\hline
\rowcolor{popblueL} Type & Forme & Exemple \\
\hline
Affirmative & sujet + have/has + participe passé & \eng{She has just got married.} \\
\hline
Négative & sujet + have/has + not + participe passé & \eng{He has not visited us yet.} \\
\hline
Interrogative & Have/Has + sujet + participe passé ? & \eng{Have you ever met him?} \\
\hline
\end{tabular}
\end{center}
Attention aux verbes irréguliers : \eng{go -- went -- gone}, \eng{see -- saw -- seen}, \eng{be -- was/were -- been}, \eng{know -- knew -- known}, \eng{meet -- met -- met}.
\end{propriete}

\begin{propriete}[Emplois du present perfect]
\begin{enumerate}
\item \cle{Expérience de vie} sans date précise : \eng{I have visited Buea.} (J'ai visité Buéa — à un moment ou à un autre.)
\item \cle{Action récente} dont le résultat est visible maintenant : \eng{She has just got married.} (Elle vient de se marier.)
\item \cle{Situation commencée dans le passé et qui continue} : \eng{They have known each other for ten years.} (Ils se connaissent depuis dix ans.)
\end{enumerate}
\end{propriete}

\begin{exemplebox}[Exemples traduits]
\begin{itemize}
\item \eng{She has just got married.} = Elle vient de se marier.
\item \eng{They have known each other for ten years.} = Ils se connaissent depuis dix ans.
\item \eng{We have lived in Buea since 2010.} = Nous vivons à Buéa depuis 2010.
\item \eng{Have you ever eaten ndolé?} = As-tu déjà mangé du ndolé ?
\end{itemize}
\end{exemplebox}

\courssub{The markers of the present perfect}

Certains petits mots signalent presque toujours le present perfect. Apprenez leur \cle{position} dans la phrase.

\begin{center}
\begin{tabularx}{\linewidth}{|l|X|X|}
\hline
\rowcolor{popblueL} Marqueur & Position & Exemple \\
\hline
\eng{just} (venir de) & entre have/has et le participe & \eng{She has just arrived.} \\
\hline
\eng{already} (déjà) & entre have/has et le participe & \eng{We have already met him.} \\
\hline
\eng{yet} (encore/déjà) & à la fin, en négatif et interrogatif & \eng{He has not come yet.} \\
\hline
\eng{ever} (déjà, jamais) & entre have/has et le participe, en question & \eng{Have you ever been to Bamenda?} \\
\hline
\eng{never} (jamais) & entre have/has et le participe & \eng{I have never lied to my parents.} \\
\hline
\eng{since} (depuis, point de départ) & devant une date ou un moment & \eng{since 2010, since Monday} \\
\hline
\eng{for} (depuis, durée) & devant une durée & \eng{for ten years, for a long time} \\
\hline
\end{tabularx}
\end{center}

\begin{attention}[Le piège du « depuis » français]
En français, on dit « Je vis ici \textbf{depuis} 2010 » avec un \emph{présent}. L'anglais refuse cette construction : il exige le present perfect. Ne dites jamais \eng{I am living here since 2010} ; dites \eng{I have lived here since 2010}. De même : « Ils se connaissent depuis dix ans » = \eng{They have known each other for ten years} (et non \eng{they know each other since ten years}).
\end{attention}

\courssub{Present perfect or past simple ?}

C'est le choix le plus délicat pour un francophone, car le français utilise souvent le passé composé dans les deux cas.

\begin{propriete}[Règle d'or]
Le present perfect ne s'emploie \cle{JAMAIS} avec un repère de temps passé et terminé (\eng{yesterday, last year, in 2010, two days ago, when I was a child}). Si une telle date apparaît, le \cle{prétérit} est obligatoire.
\begin{itemize}
\item \eng{I have visited Buea.} (expérience, sans date) $\rightarrow$ present perfect
\item \eng{I visited Buea in 2023.} (date précise, terminée) $\rightarrow$ prétérit
\end{itemize}
\end{propriete}

\begin{popfigure}
\begin{tikzpicture}[scale=1]
\draw[-{Stealth}, thick, popdark] (-0.5,0) -- (10.5,0);
\node[below, popdark] at (10.3,0) {\small time};
\node[circle, fill=popgreen, inner sep=2.5pt, label=below:{\small now}] (now) at (9,0) {};
\draw[-{Stealth}, very thick, popblue] (2,0.55) .. controls (5,1.1) .. (8.8,0.15);
\node[popblue, align=center] at (5,1.5) {\small present perfect : past $\rightarrow$ now};
\node[circle, fill=poporange, inner sep=2.5pt, label=below:{\small in 2023}] at (3,-0.9) {};
\node[poporange, align=center] at (3,-1.7) {\small past simple : finished, dated};
\end{tikzpicture}
\legende{Le prétérit est un point fermé et daté dans le passé ; le present perfect relie le passé au présent.}
\end{popfigure}

\begin{methode}[Choisir entre prétérit et present perfect]
\begin{enumerate}
\item Repérez le complément de temps dans la phrase.
\item Posez-vous la question : « Y a-t-il une date passée précise et terminée ? » (\eng{yesterday, last week, in 2015, ago}).
\item Si \textbf{oui} $\rightarrow$ utilisez le \cle{prétérit} : \eng{My cousin got married last year.}
\item Si \textbf{non} (expérience, résultat présent, \eng{since/for/just/already/yet/ever/never}) $\rightarrow$ utilisez le \cle{present perfect} : \eng{My cousin has just got married.}
\end{enumerate}
\end{methode}

\begin{exoresolu}[Past simple or present perfect ?]
Complétez avec le verbe au temps correct : (a) \eng{I \ldots{} (never / see) a lion.} (b) \eng{We \ldots{} (move) to Yaoundé in 2019.} (c) \eng{She \ldots{} (live) with her aunt since 2020.} (d) \eng{They \ldots{} (visit) their grandparents last Sunday.}
\tcblower
\textbf{Solution détaillée.} (a) \eng{I have never seen a lion.} — expérience de vie sans date, marqueur \eng{never} $\rightarrow$ present perfect. (b) \eng{We moved to Yaoundé in 2019.} — date passée précise \eng{in 2019} $\rightarrow$ prétérit obligatoire. (c) \eng{She has lived with her aunt since 2020.} — situation qui continue, marqueur \eng{since} $\rightarrow$ present perfect (le français « elle vit chez sa tante depuis 2020 » est un piège). (d) \eng{They visited their grandparents last Sunday.} — repère terminé \eng{last Sunday} $\rightarrow$ prétérit.
\end{exoresolu}

\courssub{Comparatives and superlatives}

\textbf{Observation}\par

\eng{My brother is taller than me. My grandmother is the wisest person in the family. Douala is more crowded than Buea.} Observons : les adjectifs courts prennent \eng{-er / -est}, les adjectifs longs prennent \eng{more / the most}.

\begin{propriete}[Formation des comparatifs et superlatifs]
\begin{center}
\begin{tabular}{|l|l|l|l|}
\hline
\rowcolor{popblueL} Adjectif & Comparatif & Superlatif & Exemple \\
\hline
Court (1 syllabe) & adj + \eng{-er than} & \eng{the} + adj + \eng{-est} & \eng{taller than / the tallest} \\
\hline
Court en -e & adj + \eng{-r} & \eng{the} + adj + \eng{-st} & \eng{nicer / the nicest} \\
\hline
Consonne + y & y $\rightarrow$ \eng{-ier} & \eng{the} + \eng{-iest} & \eng{happier / the happiest} \\
\hline
CVC (big, hot) & consonne doublée & consonne doublée & \eng{bigger / the biggest} \\
\hline
Long (2+ syllabes) & \eng{more} + adj + \eng{than} & \eng{the most} + adj & \eng{more patient / the most patient} \\
\hline
\end{tabular}
\end{center}
\end{propriete}

\begin{propriete}[Comparatifs et superlatifs irréguliers]
\begin{center}
\begin{tabular}{|l|l|l|}
\hline
\rowcolor{popblueL} Adjectif & Comparatif & Superlatif \\
\hline
\eng{good} (bon) & \eng{better} & \eng{the best} \\
\hline
\eng{bad} (mauvais) & \eng{worse} & \eng{the worst} \\
\hline
\eng{far} (loin) & \eng{further} & \eng{the furthest} \\
\hline
\end{tabular}
\end{center}
\end{propriete}

\begin{exemplebox}[Exemples traduits]
\begin{itemize}
\item \eng{My grandmother is the wisest person in the family.} = Ma grand-mère est la personne la plus sage de la famille.
\item \eng{My younger sister is more talkative than my brother.} = Ma sœur cadette est plus bavarde que mon frère.
\item \eng{This is the best reunion we have ever had.} = C'est la meilleure réunion de famille que nous ayons jamais eue.
\item \eng{Bamenda is further from Douala than Buea is.} = Bamenda est plus loin de Douala que Buéa.
\end{itemize}
\end{exemplebox}

\begin{attention}[Erreurs fréquentes des francophones]
\begin{itemize}
\item Ne doublez jamais la marque : pas de \eng{more taller} ni de \eng{the most biggest}. Une seule marque de comparaison.
\item N'oubliez pas \eng{than} après le comparatif : \eng{My brother is taller than me} (et non \eng{taller that me}).
\item Le superlatif exige toujours \eng{the} : \eng{the best}, jamais \eng{a best}.
\end{itemize}
\end{attention}

\courssub{Possessives : adjectives, pronouns and the genitive}

\textbf{Observation}\par

\eng{This is my book. It is mine. Her brother lives in Douala; hers is a teacher. Mary's children are polite.} Distinguons trois outils : l'\cle{adjectif possessif} (devant un nom), le \cle{pronom possessif} (à la place du nom) et le \cle{génitif} \eng{'s}.

\begin{popfigure}
\begin{tikzpicture}
\node[draw=popblue, fill=popblueL, rounded corners, align=center, font=\small\bfseries] (titre) at (4.5,3.2) {Possessives};
\node[draw=popgreen, fill=popgreenL, rounded corners, align=center, font=\small] (adj) at (0.8,1.6) {Adjective\\+ noun\\\eng{my book}};
\node[draw=poporange, fill=poporangeL, rounded corners, align=center, font=\small] (pro) at (8.2,1.6) {Pronoun\\alone\\\eng{mine}};
\draw[-{Stealth}, thick, popmuted] (titre) -- (adj);
\draw[-{Stealth}, thick, popmuted] (titre) -- (pro);
\node[align=center, font=\small] at (0.8,-0.3) {\eng{my -- mine}\\\eng{your -- yours}\\\eng{his -- his}\\\eng{her -- hers}};
\node[align=center, font=\small] at (8.2,-0.3) {\eng{its -- (rare)}\\\eng{our -- ours}\\\eng{their -- theirs}};
\end{tikzpicture}
\legende{L'adjectif possessif accompagne un nom ; le pronom possessif le remplace.}
\end{popfigure}

\begin{propriete}[Tableau des possessifs]
\begin{center}
\begin{tabular}{|l|l|l|l|}
\hline
\rowcolor{popblueL} Adjectif possessif & Pronom possessif & Français (adj.) & Français (pronom) \\
\hline
\eng{my} & \eng{mine} & mon, ma, mes & le mien, la mienne \\
\hline
\eng{your} & \eng{yours} & ton, ta, tes, votre & le tien, le vôtre \\
\hline
\eng{his} & \eng{his} & son, sa, ses (à lui) & le sien (à lui) \\
\hline
\eng{her} & \eng{hers} & son, sa, ses (à elle) & le sien (à elle) \\
\hline
\eng{its} & --- & son, sa, ses (chose) & --- \\
\hline
\eng{our} & \eng{ours} & notre, nos & le nôtre \\
\hline
\eng{their} & \eng{theirs} & leur, leurs & le leur \\
\hline
\end{tabular}
\end{center}
L'adjectif possessif anglais s'accorde avec le \cle{possesseur}, pas avec l'objet possédé : \eng{Mary loves her father} (Mary aime son père) — \eng{her} parce que Mary est une femme.
\end{propriete}

\begin{propriete}[Le génitif 's]
Pour exprimer la possession avec un nom de personne, on ajoute \eng{'s} au possesseur : \eng{my father's house} (la maison de mon père), \eng{the children's toys} (les jouets des enfants — pluriel irrégulier sans -s), \eng{my parents' village} (le village de mes parents — pluriel en -s : apostrophe seule).
\end{propriete}

\begin{attention}[Confusions à éliminer absolument]
\begin{itemize}
\item \eng{its} (possessif : son, sa) $\neq$ \eng{it's} (= \eng{it is}) : \eng{The dog wags its tail.} / \eng{It's raining.}
\item \eng{their} (leur) $\neq$ \eng{there} (là-bas) $\neq$ \eng{they're} (= \eng{they are}) : \eng{Their house is over there, and they're happy.}
\item Le possessif anglais ne prend jamais d'article : \eng{my mother}, jamais \eng{the my mother}.
\end{itemize}
\end{attention}

\begin{exoresolu}[Choose the correct possessive]
Complétez : (a) \eng{This pen is not yours; it is \ldots{} (I).} (b) \eng{Mary has lost \ldots{} keys.} (c) \eng{The cat is drinking \ldots{} milk.} (d) \eng{We are proud of \ldots{} school.}
\tcblower
\textbf{Solution détaillée.} (a) \eng{mine} — pronom possessif, car il remplace \eng{my pen}. (b) \eng{her} — adjectif possessif devant le nom \eng{keys}, accordé avec Mary (une femme). (c) \eng{its} — possessif de la chose, sans apostrophe. (d) \eng{our} — adjectif possessif de la 1\up{re} personne du pluriel devant \eng{school}.
\end{exoresolu}

\courssub{French / English : a final comparison}

\begin{center}
\begin{tabularx}{\linewidth}{|X|X|X|}
\hline
\rowcolor{popblueL} Français & English & Remarque \\
\hline
Je vis ici depuis 2010. & \eng{I have lived here since 2010.} & « depuis » + présent $\rightarrow$ present perfect \\
\hline
Il a plu hier. & \eng{It rained yesterday.} & date passée $\rightarrow$ prétérit \\
\hline
Mon frère est plus grand que moi. & \eng{My brother is taller than me.} & adjectif court : \eng{-er than} \\
\hline
C'est la meilleure élève. & \eng{She is the best student.} & irrégulier : \eng{good/better/best} \\
\hline
Ce livre est le mien. & \eng{This book is mine.} & pronom possessif, sans nom après \\
\hline
La maison de mon père. & \eng{My father's house.} & génitif \eng{'s} \\
\hline
\end{tabularx}
\end{center}

\begin{aretenir}[L'essentiel de la synthèse grammaticale]
\begin{itemize}
\item Present perfect = \eng{have/has} + participe passé ; expérience, résultat présent, durée avec \eng{since/for} ; jamais avec une date passée terminée.
\item Prétérit = action datée et terminée (\eng{yesterday, last year, in 2010, ago}).
\item Comparatif : \eng{-er than} / \eng{more ... than} ; superlatif : \eng{the -est} / \eng{the most} ; irréguliers : \eng{better/best, worse/worst, further/furthest}.
\item Adjectif possessif + nom (\eng{my book}) ; pronom possessif seul (\eng{mine}) ; génitif \eng{'s} pour les personnes.
\item Pièges : \eng{its/it's}, \eng{their/there/they're}, et le « depuis » français qui exige le present perfect.
\end{itemize}
\end{aretenir}

Ces quatre piliers grammaticaux étant désormais consolidés, vous êtes prêts à les mobiliser dans le texte d'intégration de la section suivante (\emph{Reading : integration text}), où ils apparaîtront en contexte authentique.

\coursec{Reading : integration text}

Après avoir synthétisé le vocabulaire et la grammaire du chapitre \emph{Family and Social Life}, il est temps de les remobiliser dans une situation d'intégration complète : la compréhension d'un texte authentique. C'est exactement ce que l'on te demandera à l'examen : lire un texte, repérer l'idée générale, vérifier des détails, comprendre des mots en contexte, puis réutiliser le texte pour produire. Nous allons travailler sur un récit de réunion de famille à Buea, qui rassemble le lexique de la famille et les points de grammaire du chapitre (comparatifs et superlatifs, present perfect, past perfect, pronoms relatifs).

\courssub{La méthode de lecture en trois étapes}

Avant de lire le texte, rappelons la démarche professionnelle du lecteur efficace. On ne lit jamais un texte d'examen mot à mot du début à la fin : on procède en trois passages.

\begin{methode}[Lire efficacement : skimming, scanning, checking]
\begin{enumerate}
\item \textbf{Skimming} (lecture rapide) : lis le texte une première fois rapidement, sans t'arrêter sur les mots inconnus, pour saisir l'\cle{idée générale}. C'est cette étape qui permet de choisir le meilleur titre.
\item \textbf{Scanning} (lecture ciblée) : relis en cherchant des \cle{informations précises} (noms, dates, lieux, raisons) pour répondre aux questions détaillées. Souligne dans le texte la phrase qui justifie chaque réponse.
\item \textbf{Checking} (vérification) : relis tes réponses et compare-les au texte. Vérifie l'orthographe des mots recopiés et la logique de tes justifications.
\end{enumerate}
\end{methode}

\begin{popfigure}
\begin{tikzpicture}[node distance=1.1cm, every node/.style={font=\small}]
\node[draw=popblue, fill=popblueL, rounded corners, minimum width=3.4cm, minimum height=1.2cm, align=center] (a) {\textbf{1. Skimming}\\ idée générale};
\node[draw=poporange, fill=poporangeL, rounded corners, minimum width=3.4cm, minimum height=1.2cm, align=center, right=of a] (b) {\textbf{2. Scanning}\\ détails précis};
\node[draw=popgreen, fill=popgreenL, rounded corners, minimum width=3.4cm, minimum height=1.2cm, align=center, right=of b] (c) {\textbf{3. Checking}\\ vérification};
\draw[-{Stealth}, very thick, popblue] (a) -- (b);
\draw[-{Stealth}, very thick, poporange] (b) -- (c);
\draw[-{Stealth}, thick, popmuted, dashed] (c.south) to[bend left=25] node[below, font=\footnotesize\itshape]{relecture du texte si doute} (a.south);
\end{tikzpicture}
\legende{La méthode de lecture en trois étapes : du général vers le particulier, puis vérification.}
\end{popfigure}

\courssub{Texte d'intégration : A Family Reunion in Buea}

\begin{textebox}[A Family Reunion in Buea]
Last Saturday, the Ngong family gathered in Buea for their annual reunion. Relatives came from Douala, Bamenda and even Lagos. More than sixty people filled the large courtyard of the family house, and the women prepared traditional dishes such as eru, ndolé and fried plantains.

Grandfather Ngong, the eldest member of the family, has always been the wisest of us all. During the meal, he recalled how he had paid the bride price for his wife fifty years earlier. “In those days,” he explained, “a dowry was a sign of respect between two families, not a business.” The younger cousins, who had never attended a traditional wedding, listened carefully. Some of them had travelled from far away and had not seen their uncles and aunts for years.

After lunch, the family discussed important matters: the school fees of the youngest children, the health of Grandmother Ngong, and a piece of land in Molyko. The in-laws were also present and everyone got on well with each other. Everyone spoke freely, and the meeting ended peacefully. “Family bonds are more precious than money,” Grandfather said. “When we stay united, no problem is too big.”

Everyone agreed that it was the most moving reunion they had ever had. Before leaving, the relatives promised to gather again the following year, and the children exchanged phone numbers with cousins they had just met for the first time.
\end{textebox}

\begin{definition}[Glossaire du texte]
\begin{center}
\begin{tabularx}{\linewidth}{|l|l|X|}
\hline
\rowcolor{popblueL} \textbf{English} & \textbf{Français} & \textbf{Exemple du texte} \\
\hline
a reunion & une réunion (de famille) & their annual reunion \\
\hline
to gather & se rassembler & the Ngong family gathered in Buea \\
\hline
dishes & des plats (cuisine) & traditional dishes such as eru \\
\hline
to recall & se souvenir, rappeler & he recalled how he had paid the bride price \\
\hline
a bond & un lien & family bonds are precious \\
\hline
a dowry / the bride price & la dot & he had paid the bride price \\
\hline
in-laws & beaux-parents, belle-famille & The in-laws were also present \\
\hline
to get on with & s'entendre avec & everyone got on well with each other \\
\hline
\end{tabularx}
\end{center}
\end{definition}

\begin{attention}[Faux amis et pièges de vocabulaire]
\begin{itemize}
\item \eng{dishes} signifie « plats » (préparations culinaires) et non « des disques ». \eng{The dishes were delicious.} « Les plats étaient délicieux. »
\item \eng{to recall} = se souvenir ; ce n'est pas « rappeler au téléphone ». On dit \eng{to call back} pour rappeler quelqu'un.
\item \eng{in-laws} désigne la belle-famille (beaux-parents, beaux-frères, belles-sœurs) : \eng{my mother-in-law} = ma belle-mère. Ne confonds pas avec \eng{outlaws} (hors-la-loi) !
\item \eng{actually} = en fait, réellement (et non « actuellement », qui se dit \eng{currently}).
\end{itemize}
\end{attention}

\courssub{Compréhension globale : choisir le meilleur titre}

\begin{exercice}[Global comprehension]
Choose the best title for the text. Tick only one answer.
\begin{enumerate}
\item[a)] A Traditional Wedding in Buea
\item[b)] A Family Reunion in Buea
\item[c)] The History of the Bride Price in Cameroon
\end{enumerate}
\end{exercice}

\begin{corrige}[Compréhension globale]
La bonne réponse est \textbf{b) A Family Reunion in Buea}. Le texte entier décrit une réunion de famille (arrivée des proches, repas, discussion, départ). Le titre a) est faux : on parle d'un mariage traditionnel, mais seulement comme souvenir du grand-père, ce n'est pas le sujet. Le titre c) est un détail du texte (la dot), pas l'idée générale. À l'examen, méfie-toi des titres qui reprennent un simple détail du texte : c'est le piège classique du skimming.
\end{corrige}

\courssub{Compréhension détaillée : vrai ou faux justifié}

\begin{methode}[Répondre aux questions True / False]
\begin{enumerate}
\item Lis l'affirmation et repère les mots-clés (nom, date, lieu, verbe).
\item Retrouve dans le texte la phrase qui parle de cette information (scanning).
\item Compare : l'affirmation dit-elle exactement la même chose que le texte ?
\item Écris \textbf{True} ou \textbf{False}, puis \cle{recopie la phrase du texte} qui justifie ta réponse, entre guillemets. Sans justification, la réponse est souvent \textbf{non notée}, même si elle est correcte.
\end{enumerate}
\end{methode}

\begin{exoresolu}[True or False — exercice résolu]
\textbf{Énoncé.} Read the text and say if the following statements are True or False. Justify each answer with a quotation from the text.
\begin{enumerate}
\item The reunion took place in Douala.
\item Grandfather Ngong paid the bride price fifty years ago.
\item The younger cousins have already attended many traditional weddings.
\item The family meeting ended in a quarrel.
\end{enumerate}
\tcblower
\textbf{Solution détaillée.}
\begin{enumerate}
\item \textbf{False.} “The Ngong family gathered in \emph{Buea} for their annual reunion.” Le texte dit Buea, pas Douala ; les proches viennent de Douala, mais la réunion a lieu à Buea.
\item \textbf{True.} “He recalled how he had paid the bride price for his wife fifty years earlier.” \emph{Fifty years earlier} = fifty years ago.
\item \textbf{False.} “The younger cousins, who had \emph{never} attended a traditional wedding, listened carefully.” Le mot \eng{never} contredit l'affirmation.
\item \textbf{False.} “Everyone spoke freely, and the meeting ended \emph{peacefully}.” Le texte dit le contraire d'une querelle (\eng{peacefully} = pacifiquement).
\end{enumerate}
Remarque la méthode : pour chaque réponse, un seul mot du texte (\emph{Buea}, \emph{never}, \emph{peacefully}) suffit à trancher. C'est ce mot qu'il faut souligner et citer.
\end{exoresolu}

\courssub{QCM : comprendre les mots en contexte}

À l'examen, tu ne connaîtras pas toujours tous les mots du texte. Ce n'est pas grave : le sens se devine grâce au \cle{contexte}, c'est-à-dire les phrases qui entourent le mot.

\begin{exercice}[Words in context — multiple choice]
Choose the correct meaning of the underlined words as they are used in the text.
\begin{enumerate}
\item \textbf{dowry} (paragraph 2) means:
\begin{enumerate}
\item[a)] money or goods given by the groom's family to the bride's family
\item[b)] a gift offered to the church
\item[c)] a contract signed at the town hall
\end{enumerate}
\item \textbf{in-laws} means:
\begin{enumerate}
\item[a)] your husband's or wife's family
\item[b)] your closest friends
\item[c)] your work colleagues
\end{enumerate}
\item \textbf{to get on with} someone means:
\begin{enumerate}
\item[a)] to argue with someone
\item[b)] to have a good relationship with someone
\item[c)] to travel with someone
\end{enumerate}
\end{enumerate}
\end{exercice}

\begin{corrige}[QCM vocabulaire en contexte]
\begin{enumerate}
\item \textbf{a)} — le contexte le prouve : “he had paid the bride price for his wife” et “a dowry was a sign of respect between two families”. La dot passe de la famille du marié vers celle de la mariée.
\item \textbf{a)} — \eng{in-laws} = la belle-famille : \eng{my father-in-law} (mon beau-père), \eng{my sister-in-law} (ma belle-sœur).
\item \textbf{b)} — \eng{to get on well with somebody} = bien s'entendre avec quelqu'un : \eng{I get on well with my cousins.} « Je m'entends bien avec mes cousins. »
\end{enumerate}
\end{corrige}

\courssub{Questions ouvertes : who, when, why}

\begin{exoresolu}[Short answers — exercice résolu]
\textbf{Énoncé.} Answer the following questions in complete sentences.
\begin{enumerate}
\item \textbf{Who} is the wisest member of the Ngong family?
\item \textbf{When} did the family gather in Buea?
\item \textbf{Why} was this reunion special for the younger cousins?
\end{enumerate}
\tcblower
\textbf{Solution détaillée.}
\begin{enumerate}
\item Grandfather Ngong is the wisest member of the family. (Le texte : “Grandfather Ngong... has always been the wisest of us all.”)
\item The family gathered in Buea last Saturday. (Première phrase du texte.)
\item Because they had never attended a traditional wedding and they met some cousins for the first time. (Paragraphes 2 et 4.)
\end{enumerate}
Règle d'or : réponds toujours par une \cle{phrase complète} (sujet + verbe), jamais par un seul mot. “Grandfather Ngong” seul ne vaut souvent que la moitié des points.
\end{exoresolu}

\begin{attention}[Erreurs fréquentes des francophones dans les réponses écrites]
\begin{itemize}
\item Oublier la majuscule au pronom \eng{I} et aux noms propres : écris \eng{Grandfather Ngong}, pas \eng{grandfather ngong}.
\item Traduire mot à mot : « il a payé la dot il y a cinquante ans » se dit \eng{He paid the bride price fifty years ago} (jamais \eng{he has paid... ago} : avec \eng{ago}, on utilise le prétérit, pas le present perfect).
\item Recopier la question au lieu de répondre : à \eng{Why was the reunion special?}, ne réponds pas \eng{Because the reunion was special} mais donne la raison précise trouvée dans le texte.
\end{itemize}
\end{attention}

\courssub{Grammaire en contexte : ce que le texte révise}

Le texte d'intégration n'est pas choisi au hasard : il contient toutes les structures du chapitre. Observons-les.

\begin{center}
\begin{tabularx}{\linewidth}{|X|X|l|}
\hline
\rowcolor{popblueL} \textbf{Phrase du texte} & \textbf{Point de grammaire} & \textbf{Rappel} \\
\hline
the wisest of us all & superlatif & the + adj. + -est \\
\hline
more precious than money & comparatif & more + adj. long + than \\
\hline
he had paid the bride price & past perfect & had + participe passé \\
\hline
who had never attended & pronom relatif + past perfect & who + had + p. passé \\
\hline
the most moving reunion they had ever had & superlatif + past perfect & the most + adj. \\
\hline
\end{tabularx}
\end{center}

\begin{exemplebox}[Exemples traduits]
\eng{Family bonds are more precious than money.} « Les liens familiaux sont plus précieux que l'argent. »

\eng{He recalled how he had paid the bride price fifty years earlier.} « Il se souvint comment il avait payé la dot cinquante ans plus tôt. » (Le past perfect \eng{had paid} marque une action antérieure au moment du récit.)

\eng{It was the most moving reunion they had ever had.} « C'était la réunion la plus émouvante qu'ils aient jamais eue. »
\end{exemplebox}

\courssub{Production guidée : ta propre réunion de famille}

\begin{experience}[Guided writing — My family reunion]
Rédige \textbf{cinq phrases} sur une réunion de famille que tu as vécue (ou imaginée), en réutilisant au moins \textbf{cinq mots du texte} : \eng{reunion}, \eng{to gather}, \eng{dishes}, \eng{to recall}, \eng{bond}, \eng{relatives}, \eng{eldest}, \eng{peacefully}.

Modèle de départ : \eng{Last year, my family gathered in Yaoundé for a reunion. My grandmother recalled...}
\end{experience}

\begin{exemplebox}[Exemple de production attendue]
\eng{Last December, my relatives gathered in Douala for our annual reunion. My mother prepared delicious dishes like ndolé and miondo. My grandfather, the eldest of the family, recalled his childhood in the village. I get on very well with my cousins, and we played football together. Our family bond is stronger than anything else.}
\end{exemplebox}

\begin{savaistu}[Les « family meetings » au Cameroun]
Les réunions de famille (\eng{family meetings}) sont une institution sociale majeure au Cameroun : on y règle les dots (\eng{dowries}), les successions (\eng{inheritances}), les conflits et les projets communs. Dans les régions anglophones du Nord-Ouest et du Sud-Ouest, comme à Buea ou Bamenda, ces réunions rassemblent souvent des proches venus de tout le pays et même de la diaspora. En anglais, on parle aussi de \eng{family council} pour désigner le conseil de famille qui prend les décisions importantes.
\end{savaistu}

\begin{aretenir}[L'essentiel de la section]
\begin{itemize}
\item Lis en trois étapes : \textbf{skimming} (idée générale), \textbf{scanning} (détails), \textbf{checking} (vérification).
\item Pour les questions vrai/faux, \textbf{cite toujours la phrase du texte} qui justifie ta réponse.
\item Devine le sens des mots inconnus grâce au \textbf{contexte}, pas au dictionnaire mental mot à mot.
\item Réponds aux questions ouvertes par des \textbf{phrases complètes}.
\item Réutilise le vocabulaire et la grammaire du texte dans ta propre production écrite.
\end{itemize}
\end{aretenir}

\coursec{Progressive exam-type practice}

Après la lecture du texte d'intégration, nous avons réactivé le vocabulaire de la vie familiale et sociale et vérifié notre compréhension globale d'un document authentique. Il est temps maintenant de passer à l'étape décisive de tout chapitre : \cle{l'entraînement de type examen}. Dans cette section, vous allez travailler \emph{comme le jour de l'épreuve} : quatre niveaux de difficulté croissante, un barème affiché, une auto-correction, puis un bilan personnel. L'objectif n'est pas seulement de « faire des exercices », mais d'apprendre à \cle{s'auto-évaluer} et à construire son propre plan de \cle{remédiation}.

\courssub{How to manage your time in an exam}

Avant de commencer les exercices, rappelons une évidence trop souvent oubliée : à l'examen, une grande partie des points perdus ne le sont pas par ignorance, mais par mauvaise gestion du temps et négligence de la présentation.

\begin{methode}[Managing your time during an exam]
\begin{enumerate}
\item \textbf{Read the whole paper first.} Lisez tout le sujet avant d'écrire quoi que ce soit : repérez les questions faciles, les questions difficiles et le nombre de points de chaque partie.
\item \textbf{Start with the easy questions.} Traitez d'abord ce que vous maîtrisez : cela sécurise des points et vous met en confiance.
\item \textbf{Respect the time per section.} Si une partie vaut 5 points sur 20, ne lui consacrez pas la moitié de l'heure.
\item \textbf{Keep the last 10 minutes to proofread.} Gardez toujours dix minutes pour relire : majuscules, points finaux, accords, orthographe des mots du chapitre.
\item \textbf{Never leave a blank.} Une réponse imparfaite rapporte toujours plus qu'une case vide.
\end{enumerate}
\end{methode}

\begin{propriete}[Capital letters and full stops]
À l'écrit, toute phrase anglaise commence par une \cle{majuscule} et se termine par un \cle{point} (full stop), un point d'interrogation ou un point d'exclamation. Le pronom personnel \eng{I} s'écrit \emph{toujours} avec une majuscule, même au milieu de la phrase. Les noms propres (villes, pays, personnes), les jours et les mois prennent aussi la majuscule. Les copies perdent des points sur ces détails : un texte sans ponctuation est sanctionné, même si les idées sont bonnes.

\eng{my uncle and i went to buea last saturday} \textrightarrow{} incorrect.

\eng{My uncle and I went to Buea last Saturday.} \textrightarrow{} correct. « Mon oncle et moi sommes allés à Buea samedi dernier. »
\end{propriete}

\courssub{Level 1 — Reactivation : vocabulary in context (5 points)}

Consigne en anglais, comme à l'examen : \eng{Complete each sentence with the correct word from the word bank.}

\begin{center}
\begin{tabularx}{\linewidth}{|X|X|X|}\hline
\rowcolor{popblueL} English word & French & Example \\ \hline
relative & parent (membre de la famille) & All my \eng{relatives} came to the wedding. \\ \hline
dowry & dot & The \eng{dowry} was discussed by both families. \\ \hline
ceremony & cérémonie & The \eng{ceremony} lasted three hours. \\ \hline
neighbour & voisin & Our \eng{neighbours} helped us prepare the feast. \\ \hline
to get married & se marier & My cousin will \eng{get married} in December. \\ \hline
\end{tabularx}
\end{center}

\begin{exercice}[Level 1 — Word bank]
Word bank : \eng{relatives — dowry — ceremony — neighbour — get married}

\begin{enumerate}
\item My elder sister will \ldots\ldots\ldots{} next month in Douala.
\item During the traditional \ldots\ldots\ldots{}, the two families exchange gifts.
\item In my village, the \ldots\ldots\ldots{} is negotiated before the wedding.
\item Our \ldots\ldots\ldots{} Mr Etoundi always greets us in the morning.
\item All my \ldots\ldots\ldots{} from Bamenda came for the funeral.
\end{enumerate}
Barème : 1 point par réponse correcte, orthographe exacte exigée.
\end{exercice}

\begin{corrige}[Level 1]
\begin{enumerate}
\item \eng{get married} (attention : pas de « married » seul ; on dit \eng{to get married}).
\item \eng{ceremony} (un seul « r », un seul « m »).
\item \eng{dowry} (ne pas confondre avec \eng{diary} = journal intime).
\item \eng{neighbour} (orthographe britannique ; \eng{neighbor} en américain est accepté).
\item \eng{relatives} (au pluriel, car \eng{all}).
\end{enumerate}
\end{corrige}

\courssub{Level 2 — Guided grammar : past simple or present perfect (5 points)}

Rappel de la règle du chapitre : on utilise le \cle{past simple} quand l'action est terminée et datée (\eng{yesterday, last year, in 2019, two days ago}) ; on utilise le \cle{present perfect} quand l'action a un lien avec le présent ou qu'aucune date précise n'est donnée (\eng{ever, never, already, just, since, for}).

\begin{center}
\begin{tabularx}{\linewidth}{|X|X|X|}\hline
\rowcolor{popblueL} Past simple & Present perfect & Remarque \\ \hline
\eng{I visited Buea last year.} & \eng{I have visited Buea twice.} & date précise \textrightarrow{} past simple \\ \hline
\eng{She got married in 2020.} & \eng{She has just got married.} & \eng{just} \textrightarrow{} present perfect \\ \hline
\eng{They left two hours ago.} & \eng{They have already left.} & \eng{ago} \textrightarrow{} past simple \\ \hline
\end{tabularx}
\end{center}

\begin{exercice}[Level 2 — Conjugate the verbs]
Put the verbs in brackets into the past simple or the present perfect.

\begin{enumerate}
\item My grandparents (to celebrate) \ldots\ldots\ldots{} their golden wedding last Saturday.
\item I (never / to attend) \ldots\ldots\ldots{} a traditional engagement party.
\item We (to live) \ldots\ldots\ldots{} in Yaoundé since 2015.
\item The bride (to wear) \ldots\ldots\ldots{} a beautiful dress at the ceremony two days ago.
\item (you / ever / to meet) \ldots\ldots\ldots{} my uncle from Bamenda?
\end{enumerate}
Barème : 1 point par verbe correctement conjugué.
\end{exercice}

\begin{corrige}[Level 2]
\begin{enumerate}
\item \eng{celebrated} — past simple (\eng{last Saturday} : date précise).
\item \eng{have never attended} — present perfect (\eng{never} ; participe passé de \eng{attend} : \eng{attended}).
\item \eng{have lived} — present perfect (\eng{since 2015} : durée jusqu'à aujourd'hui).
\item \eng{wore} — past simple ; verbe irrégulier \eng{wear — wore — worn}.
\item \eng{Have you ever met} — present perfect ; verbe irrégulier \eng{meet — met — met}.
\end{enumerate}
\end{corrige}

\begin{attention}[Piège des francophones : « since » et « for »]
Ne traduisez jamais « depuis 2015 » par un présent ou un passé composé maladroit : \eng{I live here since 2015} est faux. Avec \eng{since} (point de départ) et \eng{for} (durée), l'anglais exige le \cle{present perfect} : \eng{I have lived here since 2015.} (« J'habite ici depuis 2015. ») ; \eng{I have known her for ten years.} (« Je la connais depuis dix ans. ») De même, « il y a deux jours » se dit \eng{two days ago} et appelle le past simple.
\end{attention}

\courssub{Level 3 — Transformation : comparatives, superlatives and possessives (4 points)}

Rappel expressif : le \cle{comparatif} oppose deux éléments (\eng{higher than}), le \cle{superlatif} désigne l'élément extrême d'un groupe (\eng{the highest}). Les adjectifs courts (une syllabe) prennent \eng{-er / the -est} ; les adjectifs longs (deux syllabes et plus) prennent \eng{more / the most}. Les \cle{adjectifs possessifs} (\eng{my, your, his, her, our, their}) s'emploient \emph{devant un nom} ; les \cle{pronoms possessifs} (\eng{mine, yours, his, hers, ours, theirs}) le \emph{remplacent}.

\begin{exemplebox}[Exemples traduits à observer]
\eng{The dowry was higher than I expected.} « La dot était plus élevée que je ne l'espérais. »

\eng{This is the most beautiful ceremony I have ever attended.} « C'est la plus belle cérémonie à laquelle j'aie jamais assisté. »

\eng{This is my dress. It is mine.} « C'est ma robe. C'est la mienne. »
\end{exemplebox}

\begin{exoresolu}[Transformation guidée]
Énoncé : transformez \eng{My village is more peaceful than the town.} en superlatif, puis remplacez l'adjectif possessif de \eng{This is my bicycle.} par un pronom possessif.

\tcblower
Solution détaillée :
\begin{enumerate}
\item Superlatif : l'adjectif \eng{peaceful} est long (deux syllabes), donc \eng{the most peaceful}. On précise le groupe de référence : \eng{My village is the most peaceful place in the region.} (« Mon village est l'endroit le plus paisible de la région. »)
\item Pronom possessif : \eng{my bicycle} devient \eng{mine} : \eng{This bicycle is mine.} (« Ce vélo est le mien. ») Le pronom possessif ne s'écrit jamais avec apostrophe : \eng{mine}, pas \eng{mine's}.
\end{enumerate}
\end{exoresolu}

\begin{exercice}[Level 3 — Transform the sentences]
\begin{enumerate}
\item \eng{The bride's dress is more expensive than her sister's.} \textrightarrow{} superlatif (\eng{in the family}).
\item \eng{This is my notebook.} \textrightarrow{} remplacez par un pronom possessif.
\item \eng{Their house is bigger than ours.} \textrightarrow{} superlatif (\eng{in the neighbourhood}).
\item \eng{Is this your pen?} \textrightarrow{} remplacez par un pronom possessif.
\end{enumerate}
Barème : 1 point par transformation correcte.
\end{exercice}

\begin{corrige}[Level 3]
\begin{enumerate}
\item \eng{The bride's dress is the most expensive in the family.}
\item \eng{This notebook is mine.}
\item \eng{Their house is the biggest in the neighbourhood.} (adjectif court : \eng{big \textrightarrow{} the biggest}, avec redoublement du « g »).
\item \eng{Is this pen yours?}
\end{enumerate}
\end{corrige}

\courssub{Level 4 — Bac-type mini-test (20 minutes, 6 points)}

Voici maintenant la mini-épreuve complète, à réaliser en conditions réelles : vingt minutes, sans dictionnaire, sans aide. Elle comporte trois parties : compréhension écrite (2 points), grammaire (2 points), expression écrite (2 points).

\begin{textebox}[Reading text : A family reunion in Bafoussam]
Last December, the Nguemo family organised a big reunion in Bafoussam, in the West Region of Cameroon. More than eighty relatives travelled from Douala, Yaoundé, Bamenda and even from abroad to attend the event. The purpose of the gathering was to celebrate the eightieth birthday of Grandmother Mado, the oldest member of the family.

The celebration started early in the morning with a traditional ceremony. The elders blessed the grandmother and gave her gifts: a piece of traditional cloth, a calabash and a walking stick. Then the younger relatives sang songs and performed dances. Everybody wore colourful clothes, and the compound was decorated with flowers and palm branches.

In the afternoon, the family shared an enormous meal: eru, ndolé, fried plantains and grilled fish. During the meal, the eldest uncle made a speech. He reminded everybody that family solidarity is the greatest wealth. “When one member of the family suffers, the whole family suffers,” he said. “And when one member succeeds, the whole family celebrates.”

Before leaving, each relative promised to contribute money to repair the family house in the village. Grandmother Mado was so happy that she cried. “This is the most beautiful day of my life,” she said. “I have never seen the whole family together like this.”

The reunion ended late at night, but nobody wanted to go home. Such gatherings are very important in Cameroonian social life: they strengthen family ties and keep traditions alive for the younger generation.
\end{textebox}

\textbf{Glossary :} \eng{gathering} (n.) : rassemblement ; \eng{elders} (n.) : les aînés ; \eng{to bless} (v.) : bénir ; \eng{walking stick} (n.) : canne ; \eng{compound} (n.) : concession, cour ; \eng{wealth} (n.) : richesse ; \eng{to strengthen} (v.) : renforcer ; \eng{ties} (n.) : liens.

\begin{exercice}[Level 4 — Mini-test, 20 minutes]
\textbf{A. Reading comprehension (2 points).} Answer in complete sentences.
\begin{enumerate}
\item Why did the Nguemo family organise a reunion?
\item What did the eldest uncle say about family solidarity?
\end{enumerate}

\textbf{B. Grammar (2 points).} Choose or correct.
\begin{enumerate}
\item The reunion (took / has taken) place last December.
\item Grandmother Mado is (the older / the oldest) member of the family.
\item “This is (mine / my) happiest day,” she said. \textrightarrow{} choose the correct possessive.
\item Put into the present perfect : \eng{I (never / to see) such a big gathering.}
\end{enumerate}

\textbf{C. Writing (2 points).} Write a short paragraph of about 60 words on : \eng{An important ceremony in my community.} Describe the ceremony, the people who attend and why it matters to you.
\end{exercice}

\begin{corrige}[Level 4]
\textbf{A.} 1. \eng{They organised it to celebrate the eightieth birthday of Grandmother Mado.} 2. \eng{He said that family solidarity is the greatest wealth: when one member suffers, the whole family suffers, and when one succeeds, everybody celebrates.}

\textbf{B.} 1. \eng{took} (\eng{last December} : past simple). 2. \eng{the oldest}. 3. \eng{mine} (pronom possessif, sans nom derrière). 4. \eng{I have never seen such a big gathering.}

\textbf{C.} Grille : respect de la consigne et du nombre de mots (0,5 pt) ; grammaire et vocabulaire du chapitre (1 pt) ; présentation, majuscules et ponctuation (0,5 pt).

Modèle de production : \eng{An important ceremony in my community is the traditional wedding. Before the wedding, the two families meet to discuss the dowry. On the day of the ceremony, everybody wears colourful clothes and the bride is beautifully dressed. The elders bless the couple, and then we share a big meal with music and dancing. I love this ceremony because it brings the whole community together.}
\end{corrige}

\begin{attention}[Évitez les traductions mot à mot du français]
Dans l'expression écrite, ne calquez jamais le français :
\begin{itemize}
\item « J'ai 16 ans » = \eng{I am sixteen.} — jamais \eng{I have sixteen years.}
\item « J'ai faim » = \eng{I am hungry.} — jamais \eng{I have hunger.}
\item « C'est la première fois que je vois cela » = \eng{It is the first time I have seen this.}
\item « Assister à une cérémonie » = \eng{to attend a ceremony} (pas de préposition : jamais \eng{to attend to a ceremony}).
\end{itemize}
Pensez directement en anglais, avec des structures simples que vous maîtrisez.
\end{attention}

\courssub{Self-evaluation and remediation}

Chaque élève se corrige maintenant avec la grille, compte ses points sur 20 et remplit sa fiche de bilan personnel. Le barème affiché au tableau : Level 1 = 5 pts ; Level 2 = 5 pts ; Level 3 = 4 pts ; Level 4 = 6 pts.

\begin{center}
\begin{tabularx}{\linewidth}{|X|X|X|}\hline
\rowcolor{popblueL} Mes réussites & Mes erreurs & Mon plan de remédiation \\ \hline
Ex. : present perfect bien maîtrisé & Ex. : confusion \eng{mine / my} & Refaire le tableau des possessifs et 5 phrases \\ \hline
\ldots\ldots\ldots & \ldots\ldots\ldots & \ldots\ldots\ldots \\ \hline
\ldots\ldots\ldots & \ldots\ldots\ldots & \ldots\ldots\ldots \\ \hline
\end{tabularx}
\end{center}

La remédiation n'est pas une punition : c'est une méthode. À chaque type d'erreur correspond une réparation précise, comme le montre l'organigramme suivant.

\begin{popfigure}
\begin{tikzpicture}[
node distance=0.55cm and 1.1cm,
err/.style={rectangle, rounded corners, draw=poppink, fill=poppinkL, align=center, minimum width=4.2cm, minimum height=0.9cm, font=\small},
fix/.style={rectangle, rounded corners, draw=popgreen, fill=popgreenL, align=center, minimum width=4.6cm, minimum height=0.9cm, font=\small},
arr/.style={-{Stealth[length=2.5mm]}, thick, popdark}]
\node[err, align=center] (e1){Erreur de\\vocabulaire};
\node[fix, right=of e1, align=center] (f1){Refaire la fiche lexicale\\(word bank, glossaire)};
\node[err, below=of e1, align=center] (e2){Erreur de\\grammaire};
\node[fix, right=of e2, align=center] (f2){Refaire les exercices\\de la règle concernée};
\node[err, below=of e2, align=center] (e3){Erreur de\\méthode};
\node[fix, right=of e3, align=center] (f3){Refaire la méthode\\de lecture du texte};
\draw[arr] (e1) -- (f1);
\draw[arr] (e2) -- (f2);
\draw[arr] (e3) -- (f3);
\end{tikzpicture}
\legende{Organigramme de remédiation : à chaque erreur, sa réparation.}
\end{popfigure}

\begin{aretenir}[Ce qu'il faut retenir de cette séance]
\begin{itemize}
\item À l'examen : lire tout le sujet, commencer par le facile, garder dix minutes pour relire.
\item Toute phrase commence par une majuscule et finit par un point ; \eng{I} est toujours majuscule.
\item Date précise \textrightarrow{} past simple ; \eng{ever, never, just, since, for} \textrightarrow{} present perfect.
\item Adjectif possessif + nom ; pronom possessif seul : \eng{my book / the book is mine}.
\item S'auto-évaluer, identifier ses erreurs et appliquer le bon plan de remédiation : voilà la clé de la progression.
\end{itemize}
\end{aretenir}

\coursec{Self-evaluation and remediation}

Après la mini-épreuve de type examen que vous venez de passer, il serait tentant de tourner la page et de passer au chapitre suivant. Erreur ! Les meilleurs élèves ne sont pas ceux qui font le plus d'exercices, mais ceux qui savent \cle{s'auto-évaluer} honnêtement et \cle{remédier} à leurs difficultés de façon ciblée. Cette dernière séquence du chapitre \emph{Family and Social Life} vous apprend à devenir votre propre correcteur : c'est une compétence qui vous servira jusqu'au Baccalauréat et bien au-delà.

\courssub{Why self-evaluation matters}

Observez ces deux réactions d'élèves après la correction de la mini-épreuve :

\begin{exemplebox}[Two students, two attitudes]
\eng{Student A: “I got 12 out of 20. English is too hard. I give up.”}\\
« Élève A : J'ai eu 12 sur 20. L'anglais est trop dur. J'abandonne. »

\eng{Student B: “I got 12 out of 20. I lost 4 marks on the present perfect and 2 marks on vocabulary. Tonight, I will review the present perfect for fifteen minutes and learn ten family words.”}\\
« Élève B : J'ai eu 12 sur 20. J'ai perdu 4 points sur le present perfect et 2 points sur le vocabulaire. Ce soir, je révise le present perfect pendant quinze minutes et j'apprends dix mots de famille. »
\end{exemplebox}

L'élève B fait preuve de \cle{self-evaluation} (auto-évaluation) : il transforme une note globale en \textbf{diagnostic précis}, puis en \textbf{plan d'action}. Remarquez au passage le vocabulaire utile : \eng{to get a mark} (obtenir une note), \eng{to lose marks} (perdre des points), \eng{to give up} (abandonner), \eng{to review} (réviser, orthographe britannique ; on trouve aussi \eng{to revise}). C'est exactement la démarche que nous allons suivre.

\begin{methode}[How to self-evaluate honestly — S'auto-évaluer honnêtement]
\begin{enumerate}
\item \textbf{Base your judgement on what you actually produced}, not on what you think you know. Regardez votre copie, vos réponses écrites, vos productions orales — pas vos impressions. « Je connais le present perfect » ne veut rien dire ; « j'ai écrit \eng{I have visited} correctement trois fois » est un fait.
\item \textbf{Count your errors by category} : grammaire, vocabulaire, compréhension, expression. Classez chaque erreur dans une seule catégorie.
\item \textbf{Rate each skill} sur l'échelle à trois niveaux : \emph{acquired} (acquis), \emph{fragile} (fragile), \emph{to rework} (à retravailler).
\item \textbf{Choose ONE priority} : la difficulté qui vous a coûté le plus de points.
\item \textbf{Set a short, concrete action} : quinze minutes, un point précis, une date.
\end{enumerate}
\end{methode}

\begin{attention}[Le piège de l'illusion de compétence]
Relire son cahier et se dire « oui, je sais ça » n'est pas apprendre : c'est ce que les psychologues appellent l'\emph{illusion de compétence}. Le seul test fiable est la \textbf{production} : cachez la leçon et écrivez cinq phrases au present perfect, ou nommez quinze mots de famille sans regarder votre liste. Si vous bloquez, la notion est fragile — même si elle vous semblait « facile » en la relisant.
\end{attention}

\courssub{The “Can-do” grid : your personal dashboard}

La grille \emph{Can-do} (« ce que je sais faire ») reprend les six compétences clés du chapitre 1. Pour chaque item, cochez \textbf{une seule case}, en vous basant sur vos productions réelles de la mini-épreuve :

\begin{itemize}
\item \textbf{Acquired} (acquis) : je l'ai fait correctement, sans hésitation, au moins deux fois.
\item \textbf{Fragile} (fragile) : j'ai réussi une fois, hésité, ou fait une erreur que j'ai su corriger.
\item \textbf{To rework} (à retravailler) : j'ai échoué ou je n'ai pas osé répondre.
\end{itemize}

\begin{popfigure}
\begin{tikzpicture}[
  cell/.style={draw=popline, fill=popcream, minimum height=0.85cm, align=left, text=popdark, font=\small},
  head/.style={draw=popblue, fill=popblueL, minimum height=0.85cm, align=center, text=popdark, font=\small\bfseries},
  check/.style={draw=popline, fill=white, minimum height=0.85cm, align=center, minimum width=2.2cm}
]
\node[head, minimum width=8.2cm, anchor=north west] (h1) at (0,0) {Skill — I can...};
\node[head, minimum width=2.2cm, anchor=west] (h2) at (h1.east) {Acquired};
\node[head, minimum width=2.2cm, anchor=west] (h3) at (h2.east) {Fragile};
\node[head, minimum width=2.4cm, anchor=west] (h4) at (h3.east) {To rework};

\node[cell, minimum width=8.2cm, anchor=north west] (r1) at (0,-0.85) {name 15 family words (uncle, nephew, in-laws...)};
\node[check, anchor=west] (c1a) at (r1.east) {$\square$};
\node[check, anchor=west] (c1b) at (c1a.east) {$\square$};
\node[check, minimum width=2.4cm, anchor=west] (c1c) at (c1b.east) {$\square$};

\node[cell, minimum width=8.2cm, anchor=north west] (r2) at (r1.south west) {use the present perfect (I have never met...)};
\node[check, anchor=west] (c2a) at (r2.east) {$\square$};
\node[check, anchor=west] (c2b) at (c2a.east) {$\square$};
\node[check, minimum width=2.4cm, anchor=west] (c2c) at (c2b.east) {$\square$};

\node[cell, minimum width=8.2cm, anchor=north west] (r3) at (r2.south west) {compare people (taller than, the most caring...)};
\node[check, anchor=west] (c3a) at (r3.east) {$\square$};
\node[check, anchor=west] (c3b) at (c3a.east) {$\square$};
\node[check, minimum width=2.4cm, anchor=west] (c3c) at (c3b.east) {$\square$};

\node[cell, minimum width=8.2cm, anchor=north west] (r4) at (r3.south west) {understand a text about family life};
\node[check, anchor=west] (c4a) at (r4.east) {$\square$};
\node[check, anchor=west] (c4b) at (c4a.east) {$\square$};
\node[check, minimum width=2.4cm, anchor=west] (c4c) at (c4b.east) {$\square$};

\node[cell, minimum width=8.2cm, anchor=north west] (r5) at (r4.south west) {write a paragraph about my family};
\node[check, anchor=west] (c5a) at (r5.east) {$\square$};
\node[check, anchor=west] (c5b) at (c5a.east) {$\square$};
\node[check, minimum width=2.4cm, anchor=west] (c5c) at (c5b.east) {$\square$};

\node[cell, minimum width=8.2cm, anchor=north west] (r6) at (r5.south west) {speak about a ceremony (wedding, funeral...)};
\node[check, anchor=west] (c6a) at (r6.east) {$\square$};
\node[check, anchor=west] (c6b) at (c6a.east) {$\square$};
\node[check, minimum width=2.4cm, anchor=west] (c6c) at (c6b.east) {$\square$};
\end{tikzpicture}
\legende{Grille d'auto-évaluation « Can-do » du chapitre 1 : cochez une seule case par ligne, en vous fondant sur vos productions réelles.}
\end{popfigure}

\begin{exemplebox}[Reading your grid — Lire sa grille]
\eng{If most of your boxes are “Acquired”, congratulations: you are ready for the exam. Keep practising to stay sharp.}\\
« Si la plupart de vos cases sont “Acquis”, félicitations : vous êtes prêt pour l'examen. Continuez à vous entraîner pour rester en forme. »

\eng{If you have two or more boxes in “To rework”, do not panic: choose the most expensive one — the skill that cost you the most marks — and join the corresponding remediation workshop.}\\
« Si vous avez deux cases ou plus en “À retravailler”, pas de panique : choisissez la plus coûteuse — la compétence qui vous a fait perdre le plus de points — et rejoignez l'atelier de remédiation correspondant. »
\end{exemplebox}

\courssub{Differentiated remediation workshops}

Selon votre diagnostic, rejoignez l'atelier qui correspond à votre priorité. Chaque atelier dure \textbf{quinze minutes} et cible \textbf{une seule} difficulté.

\begin{attention}[Une remédiation efficace est courte et ciblée]
Quinze minutes sur UNE difficulté précise valent mieux qu'une heure de révision vague. « Réviser l'anglais » ne mène nulle part ; « refaire dix phrases au present perfect avec \eng{for} et \eng{since} » transforme réellement votre niveau. Une difficulté à la fois, toujours.
\end{attention}

\textbf{Workshop A — Vocabulary : memory card game.}\par
En binômes, fabriquez douze cartes : d'un côté le mot anglais (\eng{stepmother}, \eng{bride}, \eng{relative}), de l'autre la traduction française. À tour de rôle, tirez une carte et donnez l'équivalent, puis une phrase complète : \eng{My stepmother is very kind to me.} « Ma belle-mère est très gentille avec moi. » Une carte n'est « gagnée » que si le mot est utilisé dans une phrase correcte.

\begin{center}
\begin{tabularx}{\linewidth}{|X|X|X|}
\hline
\rowcolor{popblueL} English & Français & Exemple de phrase \\
\hline
nephew & neveu & My nephew lives in Douala. \\
\hline
mother-in-law & belle-mère (épouse du père ou mère du conjoint) & My mother-in-law cooks very well. \\
\hline
engaged & fiancé & They got engaged in Buea. \\
\hline
widow & veuve & The widow raised four children. \\
\hline
household & foyer, ménage & Our household has six people. \\
\hline
\end{tabularx}
\end{center}

\begin{attention}[Faux amis et pièges de vocabulaire du chapitre]
\begin{itemize}
\item \eng{parents} signifie « les parents » (père et mère) ; « les parents » au sens de « membres de la famille » se dit \eng{relatives} ou \eng{relations}.
\item \eng{stepmother} = belle-mère (nouvelle épouse du père) ; \eng{mother-in-law} = belle-mère (mère du conjoint). Deux réalités différentes, deux mots différents.
\item \eng{sensible} signifie « raisonnable, sensé » ; « sensible » se dit \eng{sensitive}.
\end{itemize}
\end{attention}

\textbf{Workshop B — Grammar : targeted drills.}\par
Refaites dix transformations ciblées sur le point qui vous a fait perdre des points. Exemple pour le present perfect : mettez les verbes entre parenthèses à la bonne forme.

\begin{exoresolu}[Targeted drill : present perfect]
Énoncé — Complete with the present perfect : 1. I \ldots\ (never / meet) my cousin from Bamenda. 2. She \ldots\ (live) in Yaoundé since 2019. 3. They \ldots\ (just / celebrate) their wedding anniversary.
\tcblower
Solution — 1. \eng{I have never met my cousin from Bamenda.} (\eng{never} se place entre l'auxiliaire et le participe passé ; participe passé irrégulier de \eng{meet} : \eng{met}.) 2. \eng{She has lived in Yaoundé since 2019.} (troisième personne du singulier : \eng{has} ; \eng{since} + point de départ.) 3. \eng{They have just celebrated their wedding anniversary.} (\eng{just} se place aussi entre \eng{have} et le participe passé.) Piège évité : ne jamais dire \eng{I have met him yesterday} — avec un moment précis du passé (\eng{yesterday}), l'anglais exige le prétérit simple : \eng{I met him yesterday.}
\end{exoresolu}

\textbf{Workshop C — Expression : guided rewriting.}\par
Reprenez le paragraphe que vous avez écrit pendant la mini-épreuve. Soulignez trois erreurs, réécrivez les trois phrases fautives, puis ajoutez \textbf{un connecteur d'opinion} et \textbf{une comparaison}. Comparez ensuite votre nouvelle version avec celle d'un camarade.

\courssub{Collective correction of the most frequent errors}

Voici les trois erreurs les plus fréquentes relevées dans les copies de la mini-épreuve, corrigées collectivement au tableau :

\begin{center}
\begin{tabularx}{\linewidth}{|X|X|X|}
\hline
\rowcolor{popblueL} Erreur fréquente & Correction & Explication \\
\hline
\eng{I have seen him last week.} & \eng{I saw him last week.} & Moment précis du passé (\eng{last week}, \eng{yesterday}, \eng{in 2020}) = prétérit simple, jamais de present perfect. \\
\hline
\eng{My brother is more tall than me.} & \eng{My brother is taller than me.} & Adjectif court (une syllabe) : comparatif en \emph{-er}, pas \eng{more}. \\
\hline
\eng{The family of my wife} & \eng{My wife's family} & Possession : le génitif en 's s'emploie pour les personnes. \\
\hline
\end{tabularx}
\end{center}

\courssub{Final commitment : your revision contract}

Pour terminer, chaque élève écrit sur une fiche \textbf{deux actions concrètes} de révision avant l'examen. Une bonne action est précise, courte et datée :

\begin{exemplebox}[Model commitments — Modèles d'engagements]
\eng{“Before the exam, I will learn ten family words every evening this week and test myself on Saturday.”}\\
« Avant l'examen, j'apprendrai dix mots de famille chaque soir cette semaine et je me testerai samedi. »

\eng{“I will rewrite my paragraph about my family using two comparatives and one opinion connector, and I will show it to my teacher on Thursday.”}\\
« Je réécrirai mon paragraphe sur ma famille avec deux comparatifs et un connecteur d'opinion, et je le montrerai à mon professeur jeudi. »
\end{exemplebox}

\courssub{Useful extra for the Bac : opinion connectors}

\begin{savaistu}[Ce que les examinateurs du Bac récompensent vraiment]
Au Baccalauréat, les correcteurs récompensent la \textbf{clarté} et la \textbf{correction simple} bien plus que les phrases compliquées pleines d'erreurs. Une phrase courte et juste avec un bon connecteur (\eng{In my opinion, the extended family protects its members.}) rapporte plus de points qu'une longue phrase bancale. Écrivez simple, écrivez juste, reliez vos idées.
\end{savaistu}

Pour l'\emph{essay writing}, mémorisez ces quatre connecteurs d'opinion :

\begin{center}
\begin{tabularx}{\linewidth}{|X|X|X|}
\hline
\rowcolor{popblueL} Connecteur & Fonction & Exemple \\
\hline
in my opinion & donner son avis & In my opinion, respect for elders is essential. \\
\hline
moreover & ajouter un argument & Moreover, grandparents transmit our culture. \\
\hline
however & nuancer, opposer & However, young people also need freedom. \\
\hline
to conclude & conclure & To conclude, family remains the heart of our society. \\
\hline
\end{tabularx}
\end{center}

\begin{attention}[Ponctuation des connecteurs]
Ces quatre connecteurs se placent en tête de phrase et sont \textbf{toujours suivis d'une virgule} : \eng{However, young people also need freedom.} Oublier la virgule est une faute de ponctuation sanctionnée à l'examen.
\end{attention}

\begin{exemplebox}[Model sentence — Phrase modèle]
\eng{In my opinion, the extended family is stronger than the nuclear family in Africa.}\\
« À mon avis, la famille élargie est plus solide que la famille nucléaire en Afrique. »
\end{exemplebox}

\begin{exercice}[Your turn — À vous]
Rédigez cinq phrases sur la famille au Cameroun en utilisant chacun des quatre connecteurs d'opinion au moins une fois, avec au moins deux comparatifs et une phrase au present perfect. Puis auto-évaluez votre production avec la grille Can-do : quelle compétence reste fragile ?
\end{exercice}

\begin{corrige}[Exercice — éléments de réponse]
Production possible : \eng{In my opinion, Cameroonian families are very united. Moreover, the extended family is larger than the nuclear family in most villages. However, life in cities like Douala is changing family habits. I have never attended a traditional wedding in the North, but I would love to. To conclude, family is the most important institution in our country.} Auto-évaluation : l'élève vérifie chaque phrase (connecteur présent et suivi d'une virgule ? comparatif correct : \eng{larger than}, \eng{the most important} ? present perfect bien formé : \eng{have never attended} ?) et coche la ligne correspondante de sa grille.
\end{corrige}

Vous savez désormais non seulement parler de la famille et de la vie sociale en anglais, mais aussi évaluer vos propres progrès et réparer vos faiblesses : c'est la marque des grands apprenants. Gardez votre grille Can-do dans votre cahier — elle vous servira à chaque chapitre de l'année.

\newpage

\coursec{Activité d'intégration}

\coursec{LEÇON ENA13 : Integration Activities / Evaluation / Remediation (Chapter 1)}

\courssub{Warm-up : Diagnostic rapide}

\begin{experience}[Quiz éclair : 5 minutes]
Travaillez seul. Répondez sur une feuille volante que vous remettrez (anonyme) pour que le professeur identifie les points à revoir.
\begin{enumerate}
\item Conjuguez au \cle{present perfect} : \eng{My cousins (arrive) \_\_\_\_\_\_\_\_\_\_ just \_\_\_\_\_\_\_\_\_\_ from Yaoundé.}
\item Formez le comparatif et le superlatif : \eng{noisy} ; \eng{expensive} ; \eng{good}.
\item Traduisez : « Ma famille s'est réunie pour un mariage traditionnel la semaine dernière. »
\item Citez 5 mots du lexique « Famille et vie sociale » (nom, verbe ou adjectif).
\item Corrigez l'erreur : \eng{*I have seen my uncle yesterday.} \eng{*This wedding is more bigger than the last one.}
\end{enumerate}
\end{experience}

\begin{aretenir}[Objectifs de la leçon]
Cette séance de \cle{rémédiation} et d'\cle{intégration} clôt le Chapter 1 \eng{Family and Social Life}. Elle vise à :
\begin{itemize}
\item \textbf{Mobiliser} le vocabulaire (famille, cérémonies, vie communautaire), la grammaire (present perfect, comparatif/superlatif) et les compétences (reading, writing, speaking) dans une tâche unique authentique ;
\item \textbf{S'auto-évaluer} à l'aide d'une grille \cle{Can-do} alignée sur les attendus du Baccalauréat (épreuve de composition et de compréhension) ;
\item \textbf{Remédier} aux erreurs fossilisées (dates avec present perfect, doubles comparatifs, \eng{assist} vs \eng{attend}, accord \eng{family has/have}).
\end{itemize}
\end{aretenir}

\courssub{Vocabulary Synthesis}

\begin{definition}[Mots-clés du chapitre — Catégorisés]
Le lexique suivant doit être actif (reconnu à l'oral/écrit et produit dans une composition). Vérifiez la prononciation (accent tonique souligné) et la traduction.
\end{definition}

\begin{center}
\begin{tabularx}{\linewidth}{|X|X|X|}
\hline
\rowcolor{popblueL}
\textbf{People / Personnes} & \textbf{Events \& Places / Événements \& Lieux} & \textbf{Actions \& Qualities / Actions \& Qualités} \\
\hline
\eng{relative} (parent) & \eng{wedding} (mariage) & \eng{to gather} (se rassembler) \\
\eng{in-law} (beau-parent) & \eng{funeral} (funérailles) & \eng{to celebrate} (célébrer) \\
\eng{bride} (mariée) & \eng{reunion} (réunion de famille) & \eng{to attend} (assister à) \\
\eng{groom} (marié) & \eng{ceremony} (cérémonie) & \eng{to bless} (bénir) \\
\eng{elder} (aîné, doyen) & \eng{celebration} (célébration) & \eng{traditional} (traditionnel) \\
\eng{nephew / niece} (neveu / nièce) & \eng{community} (communauté) & \eng{joyful} (joyeux) \\
\eng{guest} (invité) & \eng{neighbourhood} (quartier) & \eng{unforgettable} (inoubliable) \\
\eng{neighbour} (voisin) & \eng{ancestral home} (maison ancestrale) & \eng{solemn} (solennel) \\
\hline
\end{tabularx}
\end{center}

\begin{exemplebox}[Prononciation repère (approx. français)]
\begin{itemize}
\item \eng{relative} : « ré-la-tive » ; \eng{ceremony} : « sé-ré-mo-ni » ; \eng{neighbour} : « néi-beu » (gh muet) ;
\item \eng{funeral} : « fiou-né-ral » ; \eng{elder} : « é-lde » ; \eng{bride/groom} : « braïd / groum ».
\end{itemize}
\end{exemplebox}

\courssub{Grammar Synthesis}

\begin{propriete}[Present Perfect vs Past Simple : le choix du temps]
\begin{center}
\begin{tabularx}{\linewidth}{|X|X|X|}
\hline
\rowcolor{popblueL}
\textbf{Critère} & \textbf{Present Perfect} & \textbf{Past Simple} \\
\hline
\emph{Forme} & \eng{have/has + V-ed / 3\textsuperscript{e} col.} & \eng{V-ed / 2\textsuperscript{e} col.} \\
\hline
\emph{Sens principal} & Lien avec le \textbf{présent} (résultat, expérience vécue \emph{maintenant}) & Action \textbf{terminée} dans un passé daté \\
\hline
\emph{Marqueurs} & \eng{just, already, yet, ever, never, since, for, recently} & \eng{yesterday, last week, in 2022, two days ago, when I was...} \\
\hline
\emph{Exemple} & \eng{I have attended many weddings.} (expérience) & \eng{I attended a wedding last Saturday.} (daté) \\
\hline
\emph{Attention} & \eng{The family \textbf{has} arrived.} (collectif = singulier) & \eng{The guests \textbf{have} arrived.} (pluriel) \\
\hline
\end{tabularx}
\end{center}
\end{propriete}

\begin{propriete}[Comparatif et Superlatif : règles et exceptions]
\begin{center}
\begin{tabular}{|l|l|l|l|}
\hline
\rowcolor{popblueL}
\textbf{Type} & \textbf{Adjectif} & \textbf{Comparatif} & \textbf{Superlatif} \\
\hline
Court (1 syll.) & \eng{big, hot, loud} & \eng{bigger, hotter, louder than} & \eng{the biggest, the hottest, the loudest} \\
\hline
En \emph{-y} (2 syll.) & \eng{happy, noisy, busy} & \eng{happier, noisier, busier than} & \eng{the happiest, the noisiest, the busiest} \\
\hline
Long ($\ge$ 2 syll.) & \eng{beautiful, expensive, traditional} & \eng{more beautiful than} & \eng{the most beautiful} \\
\hline
Irréguliers & \eng{good, bad, far, much/many} & \eng{better, worse, farther/further, more} & \eng{the best, the worst, the farthest/furthest, the most} \\
\hline
\end{tabular}
\end{center}
\end{propriete}

\begin{attention}[Pièges « francophones » à bannir]
\begin{itemize}
\item \eng{*more bigger}, \eng{*more happier} $\rightarrow$ choix exclusif : \emph{-er} \textbf{ou} \emph{more}.
\item \eng{*The most best} $\rightarrow$ \eng{the best}.
\item \eng{*I have seen him yesterday} $\rightarrow$ \eng{I saw him yesterday} (date passée = Past Simple).
\item \eng{*We assisted the ceremony} $\rightarrow$ \eng{We \textbf{attended} the ceremony} (\eng{to assist} = aider).
\end{itemize}
\end{attention}

\courssub{Reading : Integration Text}

\begin{textebox}[The Student Voice — Call for Articles]
\textbf{CALL FOR ARTICLES — Issue 4: “Family and Social Life”}

Dear students,

Our next issue will celebrate family life in Cameroon. We want YOUR stories!

Write a short article (80--100 words) about a family ceremony you have attended or imagined: a traditional wedding, a family reunion, a birth celebration or a funeral. Tell us who was there, what happened, and why the event was special.

Your article must include:
\begin{itemize}
\item a catchy title;
\item at least 8 words from our “Family and Social Life” word bank;
\item 2 sentences in the present perfect;
\item 1 comparative and 1 superlative.
\end{itemize}

Deadline: Friday. The best articles will be displayed on the classroom wall and read aloud in class.

The Editorial Team
\end{textebox}

\begin{exemplebox}[Glossaire du document]
\begin{itemize}
\item \eng{call for articles} (n.) : appel à articles ;
\item \eng{issue} (n.) : numéro (d'un magazine) ;
\item \eng{catchy} (adj.) : accrocheur ;
\item \eng{word bank} (n.) : banque de mots, liste de vocabulaire ;
\item \eng{deadline} (n.) : date limite ;
\item \eng{displayed} (v., \emph{to display}) : affiché ;
\item \eng{read aloud} : lu à voix haute.
\end{itemize}
\end{exemplebox}

\begin{exercice}[Compréhension de l'appel à articles]
Répondez en anglais complet.
\begin{enumerate}
\item What is the theme of Issue 4?
\item What is the required word count?
\item List the four language constraints.
\item Find a synonym in the text for “published” (displayed).
\item Translate: “We want YOUR stories!”
\end{enumerate}
\end{exercice}

\begin{corrige}[Exercice Compréhension]
\begin{enumerate}
\item The theme is “Family and Social Life”.
\item 80 to 100 words.
\item 1) Catchy title. 2) At least 8 words from the word bank. 3) 2 present perfect sentences. 4) 1 comparative and 1 superlative.
\item Displayed.
\item « Nous voulons VOS histoires ! »
\end{enumerate}
\end{corrige}

\courssub{Progressive Exam-Type Practice}

\begin{experience}[Tâche finale : Rédiger pour \emph{The Student Voice}]
\textbf{Situation :} Votre classe de 1ère A au Lycée Bilingue de Bafoussam publie le magazine \emph{The Student Voice}. En binôme, rédigez un article d'une page (80--100 mots) pour le numéro 4.

\textbf{Consignes étape par étape :}
\begin{enumerate}
\item \textbf{Choix du sujet} (2 min) : Mariage traditionnel (Bafoussam), réunion familiale (Douala), funérailles (village), naissance. Décidez : réel ou imaginé ?
\item \textbf{Brainstorming lexical} (5 min) : Sélectionnez 10 mots de la \emph{Word Bank} (p. 3) que vous caserez. Notez-les.
\item \textbf{Plan en 3 parties} (3 min) :
      \begin{itemize}
      \item \emph{Introduction} : Quelle cérémonie ? Où ? Quand ? (Past Simple si date).
      \item \emph{Développement} : Qui ? Quoi ? (Present Perfect pour expériences/résultats).
      \item \emph{Conclusion} : Pourquoi spécial ? (Comparatif + Superlatif + avis personnel).
      \end{itemize}
\item \textbf{Rédaction brouillon} (15 min) : Écrivez 80--100 mots. \textbf{Soulignez} au crayon : 2 Present Perfect, 1 Comparatif, 1 Superlatif, 8 mots lexique.
\item \textbf{Correction croisée} (10 min) : Échangez avec un autre binôme. Utilisez la \textbf{Grille Can-do} (ci-dessous) pour évaluer. Cochez \textbf{Yes} / \textbf{Not yet}. Entourez les erreurs (grammaire, orthographe, mots manquants).
\item \textbf{Remédiation} (10 min) : Récupérez votre copie. Corrigez \textbf{toutes} les erreurs signalées. Réécrivez la version finale au propre sur une feuille A4 (titre centré, paragraphes visibles).
\item \textbf{Affichage \& Lecture} (5 min) : Affichez au mur. Les 3 meilleurs binômes lisent à voix haute. Travaillez l'intonation : $\nearrow$ énumérations, $\searrow$ fin de phrase.
\end{enumerate}
\end{experience}

\begin{methode}[Grille d'auto-évaluation et d'évaluation par les pairs — « Can-do »]
Pour chaque critère, cochez \textbf{Yes} ou \textbf{Not yet}. Si \textbf{Not yet} $\rightarrow$ reportez-vous à la leçon indiquée et refaites l'exercice de remédiation.
\begin{center}
\begin{tabularx}{\linewidth}{|X|c|c|l|}
\hline
\rowcolor{popblueL}
\textbf{Critère} & \textbf{Yes} & \textbf{Not yet} & \textbf{Remédiation (Leçon / Exercice)} \\
\hline
I can write a catchy title in English. & $\square$ & $\square$ & Leçon 2 (Writing a headline) / Ex. 3 \\
\hline
I can use at least 8 words about family and social life correctly. & $\square$ & $\square$ & Leçon 1 (Vocabulary) / Word Bank p. 3 \\
\hline
I can write 2 correct sentences in the present perfect (\eng{have/has + V3}). & $\square$ & $\square$ & Leçon 5 (Grammar) / Ex. 2, 4 \\
\hline
I can form a comparative (\eng{-er / more ... than}) and a superlative (\eng{the -est / the most}). & $\square$ & $\square$ & Leçon 7 (Grammar) / Ex. 1, 5 \\
\hline
I can organise my article in three clear parts (Intro / Body / Conclusion). & $\square$ & $\square$ & Leçon 9 (Writing structure) \\
\hline
I can respect the word limit (80--100 words). & $\square$ & $\square$ & Entraînement : compter les mots / Ex. 6 \\
\hline
I can read my article aloud with correct intonation (stress, rhythm). & $\square$ & $\square$ & Leçon 11 (Speaking) / Recording practice \\
\hline
\end{tabularx}
\end{center}
\end{methode}

\courssub{Self-Evaluation and Remediation}

\begin{exoresolu}[Article modèle : “A Wedding to Remember”]
Rédigez un article de 80 à 100 mots sur une cérémonie familiale, en respectant les quatre contraintes de l'appel à articles.
\tcblower
\textbf{Production modèle :}

\emph{\textbf{A Wedding to Remember}}

\emph{Last Saturday, our whole family gathered in Bafoussam for my cousin's traditional wedding. The bride wore the most beautiful \emph{toghu} dress I have ever seen, and the groom looked happier than ever. All our relatives, neighbours and even our in-laws from Douala attended the ceremony. The elders have blessed the couple with gifts and advice. The celebration was bigger than last year's reunion, and the food was the best in the neighbourhood. I have never danced so much in my life! It was truly a day of joy, love and community.} (98 mots)

\textbf{Commentaire détaillé des choix de langue :}
\begin{itemize}
\item \textbf{Titre accrocheur} : \eng{A Wedding to Remember} — structure nom + infinitif, typique des titres de presse (elliptique, percutant).
\item \textbf{Lexique du chapitre (11 mots)} : \eng{family, gathered, wedding, bride, groom, relatives, neighbours, in-laws, ceremony, elders, community} — contrainte des 8 mots largement remplie.
\item \textbf{2 Present Perfect (minimum)} :
      \begin{itemize}
      \item \eng{I have ever seen} : expérience de vie + \eng{ever} (superlatif \eng{the most beautiful}).
      \item \eng{The elders have blessed} : action récente avec résultat présent (le couple est béni).
      \item \eng{I have never danced} : expérience de vie + \eng{never}.
      \end{itemize}
\item \textbf{Comparatifs} : \eng{happier than ever} (adj. en \emph{-y} $\rightarrow$ \emph{-ier}) ; \eng{bigger than last year's reunion} (adj. court, redoublement consonne).
\item \textbf{Superlatifs} : \eng{the most beautiful} (adj. long) ; \eng{the best} (irrégulier \eng{good}).
\item \textbf{Cohérence temporelle} : \eng{Last Saturday} $\rightarrow$ Past Simple (\eng{gathered, wore, looked, attended}). Le Present Perfect n'apparaît que pour les expériences \textbf{sans date}.
\item \textbf{Structure} : 3 paragraphes logiques (Intro / Déroulement / Impression finale).
\end{itemize}
\end{exoresolu}

\begin{savaistu}[Culture camerounaise : Le mariage traditionnel Bamiléké]
Dans l'Ouest (Bafoussam, Dschang, Bafang), le mariage traditionnel (\eng{traditional wedding}) est une cérémonie codifiée. La dot (\eng{bride price}) est négociée entre les familles. La mariée porte le \eng{toghu} (velours noir brodé de fils colorés), le marié le \eng{ndop} ou un boubou. Les \eng{elders} (notables, \eng{fon} ou chefs de quartier) bénissent le couple (\eng{to bless}) avec de l'eau, du sel, du vin de palme. La communauté entière (\eng{community, neighbours, in-laws}) festoye pendant plusieurs jours. C'est un événement \eng{joyful}, \eng{solemn} et \eng{unforgettable}.
\end{savaistu}

\begin{aretenir}[Bilan final — Ce que je dois maîtriser pour l'évaluation]
\begin{itemize}
\item \textbf{Vocabulaire} : 20+ mots actifs (noms, verbes, adjectifs) sur la famille et les cérémonies.
\item \textbf{Grammaire} : Distinction nette \textbf{Past Simple} (fait daté) / \textbf{Present Perfect} (expérience, résultat présent). Formation sans erreur du \textbf{comparatif} et du \textbf{superlatif} (règles + irréguliers).
\item \textbf{Écrit} : Article de magazine structuré (Titre, 3 parties, 80--100 mots), cohérent, lexicale riche.
\item \textbf{Oral} : Lecture fluide, accentuation correcte (\eng{cele-BRA-tion, com-MU-ni-ty}), intonation expressive.
\item \textbf{Méthode} : \textbf{Toujours} relire sa production avec la grille Can-do avant de la rendre. C'est le secret de la progression.
\end{itemize}
\end{aretenir}

\begin{attention}[Check-list finale avant de rendre votre copie — « Last Minute Check »]
\begin{enumerate}
\item Titre présent et accrocheur ? \eng{\checkmark}
\item 80--100 mots exacts ? (Comptez !) \eng{\checkmark}
\item 8 mots \textbf{soulignés} du Word Bank ? \eng{\checkmark}
\item 2 phrases au Present Perfect \textbf{soulignées} ? (Pas de date passée avec !) \eng{\checkmark}
\item 1 Comparatif \textbf{souligné} (bien formé : \emph{-er} OU \emph{more}) ? \eng{\checkmark}
\item 1 Superlatif \textbf{souligné} (bien formé : \emph{the -est} OU \emph{the most}) ? \eng{\checkmark}
\item 3 paragraphes distincts (Intro / Body / Conclusion) ? \eng{\checkmark}
\item Orthographe vérifiée (\eng{neighbour, ceremony, beautiful, wedding}) ? \eng{\checkmark}
\end{enumerate}
Si une case n'est pas cochée $\rightarrow$ Corrigez MAINTENANT.
\end{attention}

\newpage

\coursec{Méthodes et exercices résolus}

\coursec{Integration Activities / Evaluation / Remediation — Chapter 1: Family and Social Life}

\courssub{Warm-up : Diagnostic rapide}

\begin{experience}[Diagnostic de départ — 10 min]
\textbf{Consigne.} En 5 minutes, sans cahier, réponds sur une feuille volante :
\begin{enumerate}
\item Cite 5 mots de vocabulaire du chapitre \eng{Family and Social Life} (liens de parenté, types de famille, verbes de relation).
\item Conjugue au présent simple : \eng{to look after} (3\textsuperscript{e} pers. sg.), \eng{to get on with} (1\textsuperscript{e} pers. pl.), \eng{to bring up} (3\textsuperscript{e} pers. sg.).
\item Écris une phrase correcte avec \eng{used to} (habitude passée) et une avec \eng{be used to} (habitude présente).
\item Complète : \eng{This book is \ldots\ (I)} \quad \eng{This is \ldots\ book (I)}.
\item Traduis : « Ma belle-sœur s'entend bien avec mes parents. »
\end{enumerate}
\textbf{Auto-correction guidée.} L'enseignant affiche les réponses attendues. Chaque élève note son score /5 et identifie sa zone de faiblesse (V, G, S, W). Cette grille oriente la révision personnelle de la séance.
\end{experience}

\courssub{Synthèse du vocabulaire — Chapter 1}

\begin{aretenir}[Vocabulaire thématique : Family and Social Life]
\textbf{1. Liens de parenté (Family ties)}
\begin{center}
\begin{tabularx}{\linewidth}{|X|X|X|}\hline
\rowcolor{popblueL} \textbf{English} & \textbf{Français} & \textbf{Note} \\ \hline
\eng{parents} & parents (père et mère) & pluriel seulement \\ \hline
\eng{siblings} & frères et sœurs & terme générique, neutre \\ \hline
\eng{spouse} & conjoint(e) & mari ou femme \\ \hline
\eng{in-laws} & beaux-parents, beaux-frères/sœurs & \eng{mother-in-law, brother-in-law} (traits d'union) \\ \hline
\eng{relatives} & parents (au sens large), famille & synonymes : \eng{extended family, kin} \\ \hline
\eng{stepfather / stepmother / stepbrother / stepsister} & beau-père / belle-mère / demi-frère / demi-sœur (par remariage) & \eng{step-} = par alliance \\ \hline
\eng{half-brother / half-sister} & demi-frère / demi-sœur (un parent commun) & \eng{half-} = par le sang \\ \hline
\eng{widow / widower} & veuve / veuf & statut marital, pas lien de parenté \\ \hline
\end{tabularx}
\end{center}

\textbf{2. Types de famille (Family structures)}
\begin{itemize}
\item \eng{nuclear family} : famille nucléaire (père, mère, enfants).
\item \eng{extended family} : famille élargie (grands-parents, oncles, tantes, cousins).
\item \eng{single-parent family / lone-parent family} : famille monoparentale.
\item \eng{blended family / reconstituted family} : famille recomposée.
\item \eng{childless couple} : couple sans enfants.
\end{itemize}

\textbf{3. Relations et verbes clés (Relations \& key verbs) — \cle{Apprendre avec leur préposition}}
\begin{center}
\begin{tabularx}{\linewidth}{|X|X|X|}\hline
\rowcolor{popblueL} \textbf{Verbe + préposition} & \textbf{Sens} & \textbf{Exemple} \\ \hline
\eng{to get on (well) with sb} & bien s'entendre qqn & \eng{I get on well with my sister.} \\ \hline
\eng{to look after sb/sth} & s'occuper de qqn/qch & \eng{She looks after her nephews.} \\ \hline
\eng{to bring up sb} & élever qqn (enfant) & \eng{They brought up three children.} \\ \hline
\eng{to take after sb} & ressembler à qqn (caractère/physique) & \eng{He takes after his father.} \\ \hline
\eng{to depend on sb/sth} & dépendre de & \eng{Children depend on their parents.} \\ \hline
\eng{to be responsible for sb/sth} & être responsable de & \eng{Parents are responsible for their kids.} \\ \hline
\eng{to be married to sb} & être marié à qqn & \textbf{Pas \eng{married with}!} \\ \hline
\eng{to agree with sb on sth} & être d'accord avec qqn sur qch & \textbf{Pas \eng{I am agree}!} \\ \hline
\end{tabularx}
\end{center}

\textbf{4. Événements familiaux (Family events)}
\eng{birth} (naissance), \eng{wedding} (mariage), \eng{naming ceremony} (cérémonie de nomination / \eng{outdooring} au Ghana), \eng{funeral} (enterrement), \eng{family gathering / reunion} (retrouvailles familiales), \eng{celebration} (célébration).

\textbf{5. Adjectifs utiles}
\eng{close} (proche), \eng{distant} (lointain), \eng{strict} (strict), \eng{lenient} (indulgent), \eng{supportive} (qui soutient), \eng{generation gap} (fossé des générations), \eng{mutual respect} (respect mutuel).
\end{aretenir}

\courssub{Synthèse de grammaire — Chapter 1}

\begin{propriete}[Points de grammaire essentiels du chapitre]
\textbf{1. Présent simple vs Présent continu (révision)}
\begin{center}
\begin{tabular}{|l|l|l|}\hline
\rowcolor{popblueL} \textbf{Aspect} & \textbf{Present Simple} & \textbf{Present Continuous} \\ \hline
\emph{Forme} & \eng{V / V-s (3\textsuperscript{e} sg)} & \eng{be + V-ing} \\ \hline
\emph{Emploi} & Habitudes, vérités générales, états permanents & Actions en cours, situations temporaires, projets proches \\ \hline
\emph{Mots-indices} & \eng{always, usually, often, every day, on Sundays} & \eng{now, at the moment, currently, this week} \\ \hline
\emph{Verbes d'état} & \eng{like, love, hate, know, believe, have (possession), be} $\rightarrow$ \textbf{pas de continu} & \eng{I have a car.} (pas \eng{I am having}) \\ \hline
\end{tabular}
\end{center}

\textbf{2. \eng{Used to} (habitude passée révolue) vs \eng{Be used to} (habitude actuelle)}
\begin{center}
\begin{tabular}{|l|l|l|}\hline
\rowcolor{popblueL} \textbf{Structure} & \textbf{Sens} & \textbf{Exemple} \\ \hline
\eng{used to + BV} & Habitude passée qui n'existe plus & \eng{We used to live in Buea.} (maintenant non) \\ \hline
\eng{be used to + V-ing / nom} & Être habitué à (présent/futur) & \eng{I am used to waking up early.} \\ \hline
\eng{get used to + V-ing / nom} & S'habituer à (processus) & \eng{He will get used to the noise.} \\ \hline
\end{tabular}
\end{center}
\end{propriete}

\begin{attention}[Piège classique] \eng{We used to visiting} $\times$ $\rightarrow$ \eng{We used to visit} $\checkmark$. \eng{Be used to} + \eng{-ing} = être habitué à ; \eng{Used to} + BV = autrefois.
\end{attention}

\textbf{3. Possessifs : Adjectifs vs Pronoms}
\begin{center}
\begin{tabular}{|c|c|c|c|}\hline
\rowcolor{popblueL} \textbf{Personne} & \textbf{Adjectif possessif (+ nom)} & \textbf{Pronom possessif (seul)} & \textbf{Exemple pronom} \\ \hline
1\textsuperscript{e} sg & \eng{my} & \eng{mine} & \eng{This pen is mine.} \\ \hline
2\textsuperscript{e} sg & \eng{your} & \eng{yours} & \eng{Is this yours?} \\ \hline
3\textsuperscript{e} sg (m) & \eng{his} & \eng{his} & \eng{The car is his.} \\ \hline
3\textsuperscript{e} sg (f) & \eng{her} & \eng{hers} & \eng{The bag is hers.} \\ \hline
3\textsuperscript{e} sg (n) & \eng{its} & \eng{its} (rare) & — \\ \hline
1\textsuperscript{e} pl & \eng{our} & \eng{ours} & \eng{This house is ours.} \\ \hline
2\textsuperscript{e} pl & \eng{your} & \eng{yours} & \eng{The decision is yours.} \\ \hline
3\textsuperscript{e} pl & \eng{their} & \eng{theirs} & \eng{The children are theirs.} \\ \hline
\end{tabular}
\end{center}
\cle{Règle d'or} : \eng{her/his/your/our/their/its/my} + \textbf{nom} ; \eng{hers/his/yours/ours/theirs/mine} \textbf{seul} (pas de nom après).

\textbf{4. Comparatif / Superlatif (révision rapide)}
\begin{itemize}
\item Court (1 syllabe) : \eng{fast $\rightarrow$ faster $\rightarrow$ the fastest}.
\item Long ($\ge$ 2 syllabes) : \eng{beautiful $\rightarrow$ more beautiful $\rightarrow$ the most beautiful}.
\item Irréguliers : \eng{good $\rightarrow$ better $\rightarrow$ best} ; \eng{bad $\rightarrow$ worse $\rightarrow$ worst} ; \eng{far $\rightarrow$ further/farther $\rightarrow$ furthest/farthest}.
\item \eng{as ... as} (égalité) ; \eng{not as ... as} / \eng{less ... than} (infériorité).
\end{itemize}

\textbf{5. Structure de la phrase : Ordre des mots (SVOMPT)}
\cle{Sujet} + \cle{Verbe} + \cle{Objet} + \cle{Manière} + \cle{Lieu} + \cle{Temps}.
\begin{attention}[Double sujet interdit] \eng{My father he works} $\times$ $\rightarrow$ \eng{My father works} $\checkmark$ OU \eng{He works} $\checkmark$. Un seul sujet par proposition.
\end{attention}


\courssub{Lecture : Texte d'intégration}

\begin{textebox}[Integration Text — Family Dynamics in Modern Cameroon]
In many Cameroonian families, the extended family plays a central role. Grandparents, uncles, aunts and cousins are not distant relatives: they take part in everyday decisions. When a child is born, the whole family contributes to the naming ceremony. However, in big cities like Douala and Yaoundé, some young couples now prefer to live alone with their children, in a nuclear family. Sociologists say this change is linked to modern working life and the high cost of housing. Despite this trend, the bond between siblings remains strong, and family gatherings during holidays are still cherished moments.
\end{textebox}

\begin{exercice}[Compréhension de l'intégration text — Type examen]
\begin{enumerate}
\item True or False? Justify by quoting the text: “Cousins are considered distant relatives.”
\item What event does the whole family contribute to when a child is born?
\item Find in the text a word or phrase meaning “connected to”.
\item What does “this change” (line 5) refer to?
\item According to the text, why do some young couples choose the nuclear family model in big cities? Give two reasons.
\item Translate into French: “Despite this trend, the bond between siblings remains strong.”
\end{enumerate}
\end{exercice}

\begin{corrige}[Exercice — Compréhension]
\begin{enumerate}
\item \textbf{False.} The text says: “Grandparents, uncles, aunts and cousins are \textbf{not} distant relatives.”
\item The whole family contributes to the \textbf{naming ceremony}.
\item The phrase is “\textbf{linked to}” (“this change is linked to modern working life”).
\item “This change” refers to \textbf{young couples preferring to live alone with their children, in a nuclear family} (instead of the extended family).
\item Reasons: (1) \textbf{modern working life} (mobility, schedules), (2) \textbf{high cost of housing} in Douala/Yaoundé.
\item « Malgré cette tendance, le lien entre frères et sœurs reste fort. » (\eng{bond} = lien ; \eng{siblings} = frères et sœurs ; \eng{remains} = reste).
\end{enumerate}
\end{corrige}

\courssub{Savoir-faire 1 : Mobiliser ses acquis dans une situation d'intégration}

\begin{methode}[Réussir une tâche d'intégration (integration task)]
Une tâche d'intégration te demande de \textbf{réutiliser tout le chapitre} (vocabulaire de la famille et de la vie sociale, grammaire, expressions de communication) dans \textbf{une seule production réelle} : lettre, dialogue, discours, article. Voici la démarche en 6 étapes.
\begin{enumerate}
\item \textbf{Lis la consigne deux fois} et souligne : le \cle{type de texte} (letter, speech, dialogue), le \cle{destinataire} (a friend, the headmaster), le \cle{sujet} et le \cle{nombre de mots} exigé.
\item \textbf{Fais un brainstorming de 3 minutes} : liste le vocabulaire du chapitre utile au sujet (\eng{relatives, nuclear family, extended family, to get on with, to look after, generation gap, single parent\ldots}).
\item \textbf{Choisis les outils grammaticaux du chapitre} : possessifs (\eng{my, mine, yours, his, hers, ours, theirs}), présent simple et présent continu, \eng{used to} pour les habitudes passées, comparatifs pour comparer deux générations.
\item \textbf{Planifie} : introduction (1 phrase), développement (2 ou 3 paragraphes), conclusion (1 phrase). Pour une lettre : adresse, date, formule d'appel (\eng{Dear\ldots}), formule de politesse finale (\eng{Yours sincerely} / \eng{Best wishes}).
\item \textbf{Rédige} en phrases simples et correctes plutôt qu'en phrases longues et risquées. Une idée = une phrase.
\item \textbf{Relis-toi} avec la grille C-O-P-S : \cle{C}apitals (majuscules), \cle{O}rder of words (sujet–verbe–complément), \cle{P}unctuation, \cle{S}pelling. Vérifie chaque verbe : sujet ? temps ? \eng{-s} à la 3\textsuperscript{e} personne ?
\end{enumerate}
\end{methode}

\begin{exoresolu}[Integration task — A letter about family life]
\textbf{Énoncé.} Your pen friend from London wants to know about family life in Cameroon. Write a letter of about 120 words in which you describe your family, the role of the extended family, and one family activity you enjoy. Use at least five words from the chapter vocabulary and the structure \eng{used to}.
\tcblower
\textbf{Solution détaillée.}
\begin{enumerate}
\item \textbf{Analyse de la consigne} : type = lettre amicale ; destinataire = un correspondant anglais ; contenu imposé = (a) description de la famille, (b) rôle de la famille élargie, (c) une activité appréciée ; contraintes = environ 120 mots, 5 mots du chapitre, \eng{used to}.
\item \textbf{Brainstorming} : \eng{nuclear family, extended family, relatives, to get on well with, to look after, family gathering, to share, naming ceremony}.
\item \textbf{Grammaire} : présent simple pour les habitudes (\eng{we meet every Sunday}), \eng{used to} pour le passé (\eng{we used to visit our grandmother}), possessifs (\eng{my, our}).
\end{enumerate}
\textbf{Réponse rédigée (modèle) :}

\eng{Dear Tom,}

\eng{Thank you for your letter. You asked me about family life in Cameroon, so here it is! I live in Yaoundé with my nuclear family: my parents, my two sisters and me. We get on very well together.}

\eng{The extended family is very important here. Our relatives often visit us, and when someone is sick or has a problem, the whole family helps. My grandmother used to live with us and she used to tell us wonderful stories every evening.}

\eng{My favourite family activity is our Sunday gathering: we cook together, share a big meal and talk about everything. I really enjoy these moments.}

\eng{Write soon and tell me about your family.}

\eng{Yours sincerely,}

\eng{Alain}

\textbf{Conclusion} : le texte respecte le format lettre (appel, paragraphes, formule finale), contient 6 mots du chapitre (\eng{nuclear family, extended family, relatives, get on well with, gathering, used to}), deux occurrences de \eng{used to}, et fait environ 120 mots. Toutes les exigences de la consigne sont satisfaites : c'est cela, l'intégration.
\end{exoresolu}

\courssub{Savoir-faire 2 : S'auto-évaluer et remédier aux difficultés}

\begin{methode}[S'auto-évaluer efficacement après un exercice]
L'auto-évaluation n'est pas « regarder la correction » : c'est un \textbf{diagnostic} suivi d'une \textbf{remédiation}.
\begin{enumerate}
\item \textbf{Fais l'exercice sans aide} (ni cahier, ni dictionnaire) et chronomètre-toi : cela reproduit les conditions de l'examen.
\item \textbf{Corrige en changeant de couleur} : entoure chaque erreur, ne l'efface pas.
\item \textbf{Classe chaque erreur} dans une catégorie : \cle{V} = vocabulaire (mot inconnu, faux ami), \cle{G} = grammaire (temps, accord, possessif), \cle{S} = spelling (orthographe), \cle{W} = word order (ordre des mots).
\item \textbf{Compte tes erreurs par catégorie} : la catégorie dominante révèle ta difficulté principale.
\item \textbf{Remédie} : relis la règle correspondante du chapitre (voir Synthèse de grammaire ci-dessus), recopie 3 exemples corrects, puis refais uniquement les questions ratées, deux jours plus tard.
\item \textbf{Note ta progression} dans un carnet : « possessifs : 4 erreurs le lundi, 1 erreur le jeudi ». Tant qu'une catégorie dépasse 2 erreurs, elle reste prioritaire.
\end{enumerate}
\end{methode}

\begin{exoresolu}[Auto-évaluation d'une copie fautive]
\textbf{Énoncé.} Voici la production d'un élève. 1) Repère et classe les 5 erreurs (V, G, S ou W). 2) Réécris le texte correctement.

\eng{My family is very united. I am agree with my parents on everything. My brother have twenty years and he works in Douala. We used to visiting our uncle every holidays. I get on well with all my familly.}
\tcblower
\textbf{Solution détaillée.}
\begin{enumerate}
\item \eng{I am agree} $\rightarrow$ \eng{I agree} : erreur \textbf{G} (calque du français « je suis d'accord » ; \eng{agree} est un verbe, pas un adjectif : pas de \eng{be}).
\item \eng{My brother have} $\rightarrow$ \eng{My brother has} : erreur \textbf{G} (3\textsuperscript{e} personne du singulier : \eng{have} devient \eng{has}).
\item \eng{twenty years} $\rightarrow$ \eng{twenty years old} : erreur \textbf{G} (en anglais, l'âge exige \eng{to be + number + years old}).
\item \eng{used to visiting} $\rightarrow$ \eng{used to visit} : erreur \textbf{G} (après \eng{used to}, on met la base verbale ; le gérondif en \eng{-ing} convient après \eng{to be used to}, qui signifie « être habitué à »).
\item \eng{familly} $\rightarrow$ \eng{family} : erreur \textbf{S} (un seul \eng{l} ; attention, l'adjectif \eng{familiar} prend aussi un seul \eng{l}).
\end{enumerate}
\textbf{Texte réécrit :} \eng{My family is very united. I agree with my parents on everything. My brother is twenty years old and he works in Douala. We used to visit our uncle every holiday. I get on well with all my family.}

\textbf{Bilan} : 4 erreurs de grammaire (G) sur 5 $\rightarrow$ la remédiation prioritaire de cet élève est la conjugaison (3\textsuperscript{e} personne, \eng{used to}, \eng{be + age}). Il doit refaire les exercices de grammaire du chapitre (présent simple, \eng{used to}) avant de rédiger à nouveau.
\end{exoresolu}

\courssub{Savoir-faire 3 : Réussir les questions de grammaire de type examen}

\begin{methode}[Traiter les items de grammaire (transformation, complétion, correction)]
\begin{enumerate}
\item \textbf{Identifie la notion testée} avant de répondre : possessif ? temps ? comparatif ? Relis mentalement la règle (voir Synthèse de grammaire).
\item \textbf{Pour la complétion} (fill in the blanks) : regarde les indices autour du trou — un sujet à la 3\textsuperscript{e} personne appelle le \eng{-s} ; un nom après le trou appelle un adjectif possessif (\eng{her}), un verbe ou rien appelle un pronom possessif (\eng{hers}).
\item \textbf{Pour la transformation} (rewrite without changing the meaning) : repère ce qui doit changer (actif/passif, affirmatif/négatif, \eng{used to} $\leftrightarrow$ prétérit) et \textbf{conserve le sens et le temps}.
\item \textbf{Pour la correction de phrase} : lis la phrase à voix basse ; vérifie dans l'ordre : sujet–verbe (accord), temps, ordre des mots, préposition.
\item \textbf{Vérifie ta réponse en la réinsérant} dans la phrase complète : si la phrase finale est correcte et naturelle, c'est bon.
\end{enumerate}
\end{methode}

\begin{exoresolu}[Grammar practice — Type examen]
\textbf{Énoncé.}
\begin{enumerate}
\item Complete: This is my aunt's house. It is \ldots\ (she).
\item Rewrite with \eng{used to}: My father played football every weekend when he was young.
\item Correct the mistake: My sister she lives in Buea with her husband.
\item Transform to comparative: My mother is careful. My father is careful. (equality)
\item Put in the correct form: Look! The children \ldots\ (play) in the garden right now.
\end{enumerate}
\tcblower
\textbf{Solution détaillée.}
\begin{enumerate}
\item \eng{It is \textbf{hers}.} Il n'y a pas de nom après le trou : il faut donc un \textbf{pronom possessif} (\eng{hers}), pas un adjectif possessif (\eng{her} + nom). Astuce : \eng{her} + nom, \eng{hers} tout seul.
\item \eng{My father \textbf{used to play} football every weekend when he was young.} \eng{Used to} + base verbale exprime une habitude passée révolue ; le sens (« jouait chaque week-end ») et le passé sont conservés. Attention à ne pas écrire \eng{used to played}.
\item \eng{My sister lives in Buea with her husband.} L'erreur est le \textbf{double sujet} \eng{My sister she} : en anglais, un seul sujet par proposition. On supprime soit \eng{my sister}, soit \eng{she}.
\item \eng{My mother is \textbf{as careful as} my father.} (Égalité : \eng{as + adj + as}).
\item \eng{Look! The children \textbf{are playing} in the garden right now.} Indices : \eng{Look!} + \eng{right now} $\rightarrow$ présent continu (\eng{be + V-ing}).
\end{enumerate}
\textbf{Conclusion} : chaque item teste une notion précise du chapitre (possessifs, \eng{used to}, structure de la phrase, comparatif, présent continu). Nommer la règle avant de répondre évite la moitié des erreurs.
\end{exoresolu}

\courssub{Savoir-faire 4 : Réussir les questions de compréhension écrite (reading)}

\begin{methode}[Répondre aux questions de compréhension]
\begin{enumerate}
\item \textbf{Lis d'abord les questions}, puis le texte : tu sauras quoi chercher.
\item \textbf{Repère les mots-clés} de chaque question (noms propres, dates, verbes d'action) et localise-les dans le texte.
\item \textbf{Pour les questions vrai/faux} (True or False, justify) : la justification est \textbf{obligatoire} ; cite la ligne du texte entre guillemets anglais “ … ”.
\item \textbf{Pour les questions ouvertes} : réponds par une \textbf{phrase complète} qui reprend les termes de la question ; ne recopie jamais un paragraphe entier.
\item \textbf{Pour les questions de vocabulaire} (find in the text a word meaning\ldots) : le mot demandé est \textbf{dans le texte}, pas dans ta tête ; vérifie qu'il a bien le sens du synonyme donné.
\item \textbf{Pour les questions de référence} (what does « it » refer to?) : cherche le dernier nom singulier neutre avant le pronom.
\end{enumerate}
\end{methode}

\begin{exoresolu}[Reading comprehension — Integration text (reprise)]
\textbf{Énoncé.} Read the text (Integration Text above) and answer the questions.

1) True or False? Justify: Cousins are considered distant relatives. — 2) What event does the whole family contribute to? — 3) Find in the text a word meaning “connected to”. — 4) What does “this change” refer to? — 5) Give two reasons for the shift to nuclear families in big cities.
\tcblower
\textbf{Solution détaillée.}
\begin{enumerate}
\item \eng{\textbf{False.} The text says: “Grandparents, uncles, aunts and cousins are \textbf{not} distant relatives.”} La justification cite la phrase exacte du texte : sans elle, la réponse est à moitié notée.
\item \eng{The whole family contributes to the \textbf{naming ceremony} (when a child is born).} La réponse reprend les termes de la question et reste courte.
\item \eng{The word is “\textbf{linked}” (“this change is linked to modern working life”).} \eng{Linked to} = relié à, connecté à. Le mot est bien extrait du texte, comme l'exige la consigne.
\item \eng{“This change” refers to \textbf{young couples preferring to live alone with their children, in a nuclear family}.} On remonte à l'idée exprimée juste avant : le passage de la famille élargie à la famille nucléaire.
\item \eng{(1) Modern working life. (2) The high cost of housing.} Les deux raisons sont explicitement données dans la phrase : “Sociologists say this change is linked to modern working life and the high cost of housing.”
\end{enumerate}
\textbf{Conclusion} : toutes les réponses sont des phrases complètes, ancrées dans le texte, sans recopie inutile. C'est la méthode qui garantit le maximum de points en compréhension écrite.
\end{exoresolu}

\courssub{Savoir-faire 5 : Consolider le vocabulaire du chapitre}

\begin{methode}[Réviser et réutiliser le vocabulaire thématique]
\begin{enumerate}
\item \textbf{Organise le lexique en familles de sens} : les liens familiaux (\eng{parents, siblings, spouse, in-laws, relatives}), les types de famille (\eng{nuclear, extended, single-parent, blended family}), les relations (\eng{to get on with, to quarrel, to support, to look after, to bring up}), les événements (\eng{birth, wedding, naming ceremony, funeral}).
\item \textbf{Apprends les mots avec leur préposition} : \eng{to get on \textbf{with} somebody, to be married \textbf{to}, to depend \textbf{on}, to be responsible \textbf{for}, to agree \textbf{with} sb \textbf{on} sth}.
\item \textbf{Apprends les collocations} : \eng{a close friend, a family gathering, to raise children, mutual respect, generation gap, naming ceremony}.
\item \textbf{Teste-toi dans les deux sens} : anglais $\rightarrow$ français, puis français $\rightarrow$ anglais (le sens actif est celui de la production écrite).
\item \textbf{Réutilise chaque mot dans une phrase personnelle} : un mot utilisé est un mot mémorisé.
\end{enumerate}
\end{methode}

\begin{exoresolu}[Vocabulary practice — Type examen]
\textbf{Énoncé.}
\begin{enumerate}
\item Complete with a word from the chapter: A family with a mother, a father and their children is a \ldots\ family.
\item Choose the correct preposition: She has been married \ldots\ (to / with) my uncle for ten years.
\item Find the odd one out: \eng{nephew — niece — widow — cousin}.
\item Translate: « Ma belle-mère s'occupe bien de ses petits-enfants. »
\item Complete the collocation: We had a big family \ldots\ last Christmas. Everyone came.
\end{enumerate}
\tcblower
\textbf{Solution détaillée.}
\begin{enumerate}
\item \eng{A \textbf{nuclear} family.} Par opposition à \eng{extended family} (famille élargie, avec grands-parents, oncles, tantes).
\item \eng{married \textbf{to} my uncle.} En anglais, on est « marié à » quelqu'un : \eng{to be married to}. Jamais \eng{married with} (calque du français « marié avec »).
\item \eng{\textbf{Widow}.} \eng{Nephew, niece, cousin} désignent des liens de parenté ; \eng{widow} (veuve) désigne un statut marital, pas un lien de parenté.
\item \eng{My mother-in-law looks after her grandchildren well.} Points de vigilance : \eng{mother-in-law} avec traits d'union ; \eng{to look after} = s'occuper de ; \eng{grandchildren} en un seul mot ; l'adverbe \eng{well} se place après le complément.
\item \eng{family \textbf{gathering}} (ou \eng{reunion}). Collocation fixe du chapitre.
\end{enumerate}
\textbf{Conclusion} : ces formats (complétion, préposition, intrus, traduction, collocation) sont les formats classiques d'évaluation du lexique. Maîtriser les prépositions et les collocations est la clé pour le Bac.
\end{exoresolu}

\begin{attention}[Les 5 erreurs qui coûtent des points au Bac]
\begin{enumerate}
\item \textbf{Le calque « je suis d'accord »} : \eng{I am agree} est faux. \eng{Agree} est un verbe : écris \eng{I agree with you.} Même piège avec \eng{I am depend on} $\rightarrow$ \eng{I depend on}.
\item \textbf{Le double sujet} : \eng{My father he works in Douala} est faux. Un seul sujet par proposition : \eng{My father works in Douala.}
\item \textbf{La confusion adjectif / pronom possessif} : \eng{This book is her} est faux ; sans nom derrière, il faut le pronom : \eng{This book is hers.} Rappel : \eng{its} (possessif) sans apostrophe, \eng{it's} = \eng{it is}.
\item \textbf{\eng{Used to} + gérondif} : \eng{We used to living together} est faux. Après \eng{used to}, base verbale : \eng{We used to live together.} Ne confonds pas avec \eng{to be used to living} (= être habitué à vivre).
\item \textbf{Les prépositions et l'orthographe du lexique} : \eng{married with} $\rightarrow$ \eng{married to} ; \eng{to discuss about} $\rightarrow$ \eng{to discuss something} (pas de préposition) ; \eng{familly} $\rightarrow$ \eng{family} (un seul \eng{l}) ; \eng{wich} $\rightarrow$ \eng{which} ; \eng{belive} $\rightarrow$ \eng{believe}.
\end{enumerate}
\end{attention}

\begin{savaistu}[Culture : La famille au Cameroun et dans le monde anglophone]
Au Cameroun, la \eng{extended family} est une réalité sociale forte : les \eng{naming ceremonies} (baptêmes traditionnels, \eng{outdooring} chez les Bassa, \eng{kombo} chez les Bamiléké) mobilisent tout le clan. Au Royaume-Uni et aux USA, la \eng{nuclear family} est la norme statistique, mais les \eng{single-parent families} et \eng{blended families} augmentent. Le \eng{generation gap} s'observe partout : les jeunes (Gen Z) communiquent par réseaux sociaux, les aînés privilégient la visite physique. Le respect des \eng{elders} (aînés) reste une valeur partagée entre cultures africaines et valeurs traditionnelles britanniques (\eng{respect for elders}).
\end{savaistu}

\begin{exercice}[Entraînement final — Mini-Bac Blanc (30 min)]
\textbf{Partie A : Grammaire (7 pts)}
\begin{enumerate}
\item Put the verbs in brackets in the correct tense (Present Simple or Continuous): Every morning, my mother \ldots\ (wake) up early. Today, she \ldots\ (prepare) a special meal for the naming ceremony.
\item Rewrite using \eng{used to}: My grandparents lived in a big compound. They farmed cocoa.
\item Choose the right form: This pen is (my / mine). That car is (their / theirs).
\item Correct: My uncle he is a doctor. He works in Bamenda.
\item Comparative: My sister is \ldots\ (old) than me. But I am \ldots\ (tall) as her.
\end{enumerate}

\textbf{Partie B : Vocabulaire (6 pts)}
\begin{enumerate}
\item Complete: The son of my brother is my \ldots\ .
\item Translate: « Mes beaux-parents habitent à Bafoussam. »
\item Odd one out: \eng{husband — wife — spouse — cousin}.
\item Preposition: He takes \ldots\ his father (he looks like him).
\item Collocation: We need to \ldots\ (organiser) a family meeting.
\item Spelling: Correct \eng{neice, wich, recieve, arguement}.
\end{enumerate}

\textbf{Partie C : Expression écrite (7 pts)}
\textbf{Sujet.} “Some people think that the extended family is outdated in modern society. Do you agree?” Write an article of 150–180 words for your school magazine. Use at least 6 words from the chapter and two different grammar points (e.g., \eng{used to}, comparatives, present simple/continuous).
\end{exercice}

\begin{corrige}[Mini-Bac Blanc — Éléments de correction]
\textbf{A. Grammaire}
1. \eng{wakes up} (habitude) / \eng{is preparing} (action en cours, aujourd'hui).
2. \eng{My grandparents used to live in a big compound. They used to farm cocoa.}
3. \eng{mine} (pronom, seul) / \eng{theirs} (pronom, seul).
4. \eng{My uncle is a doctor. He works in Bamenda.} (suppression double sujet).
5. \eng{older than} / \eng{as tall as}.

\textbf{B. Vocabulaire}
1. \eng{nephew}.
2. \eng{My in-laws live in Bafoussam.} (\eng{in-laws} pluriel, \eng{live} verbe).
3. \eng{cousin} (lien de parenté vs statut marital pour les 3 autres).
4. \eng{after} (\eng{take after} = ressembler à).
5. \eng{organize} / \eng{hold} / \eng{arrange} a family meeting.
6. \eng{niece, which, receive, argument}.

\textbf{C. Expression écrite — Pistes}
\emph{Introduction} : Define extended vs nuclear. Thesis: “I partly agree: the form changes, the function remains.”
\emph{Body 1} : Modern constraints (work, housing) $\rightarrow$ nuclear families practical. \eng{Used to}: “Grandparents used to live with us.”
\emph{Body 2} : But emotional/financial support vital. \eng{Comparative}: “Support is more crucial than ever.” \eng{Vocab}: \eng{relatives, gathering, naming ceremony, mutual respect, generation gap, to look after}.
\emph{Conclusion} : Balance needed. Call to action: “Let’s keep the bond alive.”
\emph{Grille COPS} : Vérifier \eng{-s} (3\textsuperscript{e} sg), \eng{used to + BV}, possessifs, prépositions (\eng{married to, agree with}), orthographe \eng{family, receive, argument}.
\end{corrige}

\newpage

\coursec{Exercices}

\coursec{Integration Activities / Evaluation / Remediation (Chapter 1)}

\courssub{Warm-up : Diagnostic rapide}

\begin{exercice}[QCM : Grammaire et vocabulaire du chapitre]
Complète chaque phrase en choisissant la bonne option (a, b ou c). Une seule réponse est correcte.
\begin{enumerate}
    \item My aunt \underline{\hspace{2cm}} in Yaoundé \underline{\hspace{1cm}} 2010.
    \begin{enumerate}
        \item lives / for
        \item has lived / since
        \item is living / since
    \end{enumerate}
    \item My father's brother's daughter is my \underline{\hspace{2cm}}.
    \begin{enumerate}
        \item niece
        \item nephew
        \item cousin
    \end{enumerate}
    \item \underline{\hspace{1cm}} book is this? It's \underline{\hspace{1cm}}.
    \begin{enumerate}
        \item Who / my sister
        \item Whose / my sister's
        \item Who's / my sisters'
    \end{enumerate}
    \item This traditional wedding was \underline{\hspace{2cm}} the last one.
    \begin{enumerate}
        \item more colourful than
        \item most colourful than
        \item colourfuller than
    \end{enumerate}
    \item My parents \underline{\hspace{2cm}} twenty years ago.
    \begin{enumerate}
        \item get married
        \item have got married
        \item got married
    \end{enumerate}
\end{enumerate}
\end{exercice}

\begin{exercice}[Vrai ou Faux justifié]
Lis le texte suivant et dis si les affirmations sont vraies (V) ou fausses (F). Justifie chaque réponse par une phrase exacte du texte.
\begin{textebox}[Family Changes in Modern Cameroon]
In many Cameroonian families today, the extended family system is evolving. Grandparents used to live with their children and grandchildren under the same roof. Nowadays, young couples often prefer to set up their own nuclear families in cities like Douala or Buea. However, strong ties remain: during holidays, everyone gathers in the village to celebrate weddings, births or funerals. The \cle{bride price} (dowry) ceremony is still a key step before marriage, uniting two clans.
\end{textebox}
\begin{enumerate}
    \item Grandparents never live with their grandchildren anymore.
    \item Young couples in cities usually live in nuclear families.
    \item Family ties have completely disappeared in modern times.
    \item The bride price ceremony takes place after the wedding.
    \item Holidays are moments when the extended family meets.
\end{enumerate}
\end{exercice}

\courssub{Synthèse de vocabulaire : La famille}

\begin{aretenir}[Vocabulaire clé : Family and Social Life]
\begin{center}
\begin{tabularx}{\linewidth}{|X|X|X|}
\hline
\rowcolor{popblueL} \textbf{Catégorie} & \textbf{English} & \textbf{Français} \\
\hline
Nuclear family & father, mother, parents, son, daughter, children, brother, sister, siblings & père, mère, parents, fils, fille, enfants, frère, sœur, frères et sœurs \\
\hline
Extended family & grandfather, grandmother, grandparents, grandson, granddaughter, uncle, aunt, cousin, nephew, niece, great-grandfather, great-grandmother & grand-père, grand-mère, grands-parents, petit-fils, petite-fille, oncle, tante, cousin(e), neveu, nièce, arrière-grand-père, arrière-grand-mère \\
\hline
In-laws / Step-family & father-in-law, mother-in-law, parents-in-law, brother-in-law, sister-in-law, stepfather, stepmother, stepbrother, stepsister, half-brother, half-sister, spouse & beau-père, belle-mère, beaux-parents, beau-frère, belle-sœur, beau-père (par remariage), belle-mère (par remariage), demi-frère, demi-sœur, conjoint(e) \\
\hline
Events & wedding, marriage, engagement, bride, groom, bride price, dowry, ceremony, funeral, birth celebration, reunion, knock-door (custom) & mariage, noces, fiançailles, mariée, marié, dot, cérémonie, enterrement, fête de naissance, retrouvailles, \eng{knock-door} (coutume) \\
\hline
Verbs / Adjectives & to get married, to marry, to divorce, to separate, to adopt, to raise, to support, to gather, to rely on, close-knit, nuclear, extended, widowed, engaged & se marier, divorcer, séparer, adopter, élever, soutenir, se réunir, compter sur, soudée, nucléaire, élargie, veuf/veuve, fiancé(e) \\
\hline
\end{tabularx}
\end{center}
\end{aretenir}

\begin{exercice}[Vocabulaire : L'arbre généalogique]
Complète l'arbre généalogique ci-dessous avec les mots de la liste. Chaque mot s'utilise une seule fois.
\begin{center}
\textbf{Liste :} \eng{stepfather, niece, in-laws, spouse, siblings, nephew, half-sister, great-grandmother}
\end{center}
\begin{enumerate}
    \item My mother's new husband is my \underline{\hspace{3cm}}.
    \item My brother's daughter is my \underline{\hspace{3cm}}.
    \item My husband's parents are my \underline{\hspace{3cm}}.
    \item My \underline{\hspace{3cm}} and I share only one parent (same father, different mothers).
    \item My \underline{\hspace{3cm}} is the mother of my grandmother.
    \item My sister's son is my \underline{\hspace{3cm}}.
    \item My \underline{\hspace{3cm}} (my husband or wife) works in a bank.
    \item My brothers and sisters are my \underline{\hspace{3cm}}.
\end{enumerate}
\end{exercice}

\courssub{Synthèse de grammaire : Les temps du récit et du vécu}

\begin{propriete}[Récapitulatif des temps étudiés au Chapitre 1]
\begin{center}
\begin{tabularx}{\linewidth}{|l|X|X|l|}
\hline
\rowcolor{popblueL} \textbf{Temps} & \textbf{Forme} & \textbf{Emploi principal} & \textbf{Marqueurs} \\
\hline
Present Simple & V / V-s (3e pers. sg.) & Habitudes, vérités générales, situations permanentes & usually, often, every day, always \\
\hline
Present Continuous & be + V-ing & Action en cours, situation temporaire, futur programmé & now, at the moment, look!, listen! \\
\hline
Past Simple & V-ed / 2e colonne (irr.) & Action terminée dans le passé, datée & yesterday, last week, ... ago, in 2010 \\
\hline
Present Perfect Simple & have/has + V-en / 3e col. & Lien passé-présent : expérience, durée (since/for), résultat présent & since, for, just, already, yet, ever, never \\
\hline
Present Perfect Continuous & have/has been + V-ing & Action commencée dans le passé et qui continue (durée) & since, for, how long \\
\hline
Past Perfect & had + V-en / 3e col. & Antériorité par rapport à un autre passé (le "passé du passé") & before, after, by the time, when \\
\hline
\end{tabularx}
\end{center}
\end{propriete}

\begin{attention}[Pièges fréquents pour francophones]
\begin{itemize}
    \item \textbf{Since / For :} \eng{Since} + point de départ (2010, Monday, childhood) ; \eng{For} + durée (ten years, two hours).
    \item \textbf{Présent Perfect vs Past Simple :} Ne pas utiliser le Present Perfect avec un marqueur de temps passé daté (\eng{yesterday, last year}). \eng{*I have seen him yesterday} $\rightarrow$ \eng{I saw him yesterday}.
    \item \textbf{Verbes d'état (stative) :} \eng{know, like, love, hate, want, have (possession), be, understand} $\rightarrow$ généralement pas de \eng{-ing}. \eng{*I am knowing him} $\rightarrow$ \eng{I have known him}.
    \item \textbf{Comparatif de supériorité :} Court adjectif (1 syllabe) $\rightarrow$ \eng{-er than} (\eng{taller}) ; Long adjectif ($\ge$ 2 syll.) $\rightarrow$ \eng{more ... than} (\eng{more beautiful}). \textbf{Pas de "more taller"}.
    \item \textbf{Passif :} \eng{be} (au temps requis) + \eng{Participe Passé}. \eng{They organised it} $\rightarrow$ \eng{It was organised}.
    \item \textbf{Discours rapporté :} Backshift (recul des temps) + changement pronoms/adverbes. \eng{"I am happy"} $\rightarrow$ \eng{He said he was happy}.
\end{itemize}
\end{attention}


\begin{exercice}[Conjugaison : Temps du récit et du vécu]
Mets les verbes entre parenthèses au temps qui convient (Present Simple, Present Continuous, Past Simple, Present Perfect Simple ou Continuous, Past Perfect).
\begin{enumerate}
    \item Look! My little brother \eng{(play)} with his new toys right now.
    \item We \eng{(know)} each other \eng{(since)} childhood.
    \item Last December, my sister \eng{(get)} engaged during the Christmas party.
    \item How long \eng{you / wait} for the \cle{bride}? She is late!
    \item My grandfather \eng{(never / fly)} in a plane before his trip to London last year.
    \item Usually, the family \eng{(gather)} for Sunday lunch, but today we \eng{(eat)} at a restaurant.
\end{enumerate}
\end{exercice}

\courssub{Lecture : Texte d'intégration}

\begin{exercice}[Compréhension de texte : Family Changes in Modern Cameroon]
Reprends le texte de l'exercice 2 (Vrai/Faux) et réponds aux questions suivantes en anglais complet.
\begin{enumerate}
    \item What did grandparents use to do in the past?
    \item What do young couples prefer nowadays?
    \item What happens during holidays?
    \item When does the bride price ceremony take place?
    \item Find in the text a synonym for "changing" (l. 1).
    \item Find in the text the opposite of "weak" (l. 5).
\end{enumerate}
\end{exercice}

\courssub{Entraînement progressif (Type Examen)}

\begin{exercice}[Traduction : Du français vers l'anglais]
Traduis les phrases suivantes en anglais. Sois attentif aux temps, aux comparatifs et au vocabulaire de la famille.
\begin{enumerate}
    \item Mes beaux-parents habitent à Bamenda depuis dix ans.
    \item Ce mariage était plus traditionnel que le précédent.
    \item Ma demi-sœur est plus âgée que moi.
    \item Ils se sont mariés il y a deux ans, mais ils n'ont pas encore d'enfants.
    \item Depuis quand ton oncle est-il veuf ?
\end{enumerate}
\end{exercice}

\begin{exercice}[Transformation de phrases]
Réécris chaque phrase en commençant par les mots donnés, sans en changer le sens.
\begin{enumerate}
    \item My father is older than my mother. \\
    My mother is not \underline{\hspace{5cm}}.
    \item This is the most beautiful traditional outfit I have ever seen. \\
    I have never \underline{\hspace{5cm}}.
    \item "Where does your cousin live?" he asked me. \\
    He asked me \underline{\hspace{5cm}}.
    \item They organised the dowry ceremony last week. \\
    The dowry ceremony \underline{\hspace{5cm}}.
    \item My nephew is too young to get married. \\
    My nephew is not \underline{\hspace{5cm}}.
\end{enumerate}
\end{exercice}

\begin{exercice}[Texte à trous : Grammaire et vocabulaire mixtes]
Complète le texte avec un mot approprié (verbe conjugué, préposition, pronom, adjectif, mot de vocabulaire).
\begin{textebox}[A Family Reunion in the Village]
Every year, during the long holidays, my extended family \eng{(1)} \underline{\hspace{1cm}} in our ancestral village. It \eng{(2)} \underline{\hspace{1cm}} a tradition that \eng{(3)} \underline{\hspace{1cm}} for generations. This year, my \eng{(4)} \underline{\hspace{1cm}} (my mother's brother) arrived \eng{(5)} \underline{\hspace{1cm}} his \eng{(6)} \underline{\hspace{1cm}} (wife) and their three children. We \eng{(7)} \underline{\hspace{1cm}} not \eng{(8)} \underline{\hspace{1cm}} them for two years. The \eng{(9)} \underline{\hspace{1cm}} (old) cousin, Paul, \eng{(10)} \underline{\hspace{1cm}} just \eng{(11)} \underline{\hspace{1cm}} his exams. He is \eng{(12)} \underline{\hspace{1cm}} intelligent \eng{(13)} \underline{\hspace{1cm}} his younger brother, but less serious. In the evening, the \eng{(14)} \underline{\hspace{1cm}} (grandparents) tell stories about our \eng{(15)} \underline{\hspace{1cm}} (ancestors).
\end{textebox}
\end{exercice}

\begin{exercice}[Appariements : Fonctions de communication et structures]
Associe chaque situation (1--6) à l'énoncé anglais approprié (A--F).
\begin{center}
\begin{tabular}{|l|l|}
\hline
\textbf{Situation} & \textbf{Énoncé} \\
\hline
1. Présenter sa fiancée à son père. & A. "How long have you been married?" \\
2. Demander depuis quand un couple est marié. & B. "Dad, this is Mary, my fiancée." \\
3. Comparer deux cérémonies de dot. & C. "My sister gets on well with her in-laws." \\
4. Dire qu'on s'entend bien avec sa belle-famille. & D. "The first ceremony was more formal than the second." \\
5. Demander des nouvelles d'un neveu. & E. "What does your nephew look like?" \\
6. Décrire physiquement un membre de la famille. & F. "How is your nephew doing at school?" \\
\hline
\end{tabular}
\end{center}
\end{exercice}

\courssub{Préparation au Baccalauréat \& Auto-évaluation}

\begin{exercice}[Compréhension écrite et questions de langue — Type Bac]
Lis attentivement le texte suivant et réponds aux questions en anglais. Respecte le nombre de mots indiqué pour les réponses longues.
\begin{textebox}[The Changing Face of the African Family]
The structure of the African family has undergone significant transformations over the past fifty years. Traditionally, the extended family was the cornerstone of society. Several generations lived together in a compound, sharing resources, raising children collectively and caring for the elderly. The concept of “it takes a village to raise a child” was a daily reality. Uncles, aunts and grandparents played active roles in education and discipline.

Today, urbanisation and the search for employment have driven young adults to cities such as Yaoundé, Lagos or Nairobi. The nuclear family — parents and children only — has become the dominant model in urban areas. High rents and the cost of living make it difficult to accommodate relatives. Consequently, the support network has weakened. Working parents often rely on housemaids or day-care centres, and the elderly are sometimes left alone in villages.

However, technology bridges the gap. Video calls allow daily contact with grandparents. Financial remittances sent via mobile money sustain rural homes. Moreover, major events like weddings, funerals and \cle{bride price} ceremonies still trigger massive reunions. The \cle{clan} identity remains strong; one's surname opens doors and implies mutual obligation. The African family has not disappeared; it has adapted, stretching its bonds across kilometres rather than metres.
\end{textebox}

\textbf{Questions :}
\begin{enumerate}
    \item \textbf{True or False? Justify with a quote from the text.}
    \begin{enumerate}
        \item In the past, only parents disciplined the children.
        \item Urbanisation has strengthened the extended family system in cities.
        \item Technology helps maintain family ties today.
    \end{enumerate}
    \item \textbf{Multiple Choice. Choose the correct answer.}
    \begin{enumerate}
        \item The phrase “cornerstone of society” (l. 2) means:
        \begin{itemize}
            \item[a)] a stone used for building corners
            \item[b)] the most important part
            \item[c)] a traditional house
        \end{itemize}
        \item According to the text, why do working parents in cities use day-care centres?
        \begin{itemize}
            \item[a)] They dislike their relatives.
            \item[b)] The support network has weakened.
            \item[c)] It is a cultural tradition.
        \end{itemize}
    \end{enumerate}
    \item \textbf{Short Answers (Max 20 words each).}
    \begin{enumerate}
        \item Name two major events that still bring the extended family together.
        \item How do city dwellers support their rural relatives financially?
    \end{enumerate}
    \item \textbf{Language Work.}
    \begin{enumerate}
        \item Find in the text a synonym for “old people” (paragraph 2).
        \item Find in the text the opposite of “strengthened” (paragraph 2).
        \item Rewrite in the Passive Voice: “Urbanisation has driven young adults to cities.”
        \item Put into Reported Speech: “The clan identity remains strong,” the sociologist said.
    \end{enumerate}
\end{enumerate}
\end{exercice}

\begin{exercice}[Grammaire en contexte et Traduction — Type Bac]
\textbf{Partie A : Conjugaison (10 points)}
Complète le texte en conjuguant les verbes entre parenthèses au temps approprié (Past Simple, Present Perfect Simple/Continuous, Past Perfect, etc.).
\begin{textebox}[A Traditional Wedding in Bafut]
Last Saturday, the Fon's palace \eng{(1. be)} \underline{\hspace{2cm}} full of guests. My cousin Achu \eng{(2. just / return)} \underline{\hspace{2cm}} from Germany where he \eng{(3. study)} \underline{\hspace{2cm}} engineering for four years. He \eng{(4. not / see)} \underline{\hspace{2cm}} his betrothed, Njwen, \eng{since} the \cle{knock-door} ceremony two years ago. By the time the \cle{bride price} \eng{(5. be)} \underline{\hspace{2cm}} paid, the sun \eng{(6. already / set)} \underline{\hspace{2cm}}. The \cle{quarter-head} \eng{(7. announce)} \underline{\hspace{2cm}} that the couple \eng{(8. become)} \underline{\hspace{2cm}} husband and wife. Since then, they \eng{(9. live)} \underline{\hspace{2cm}} happily in their new home.
\end{textebox}

\textbf{Partie B : Traduction (5 points)}
Traduis en anglais :
\begin{enumerate}
    \item Si j'avais su, je serais venu à la cérémonie de dot.
    \item Ma grand-mère me racontait des histoires quand j'étais petit.
\end{enumerate}
\end{exercice}

\begin{exercice}[Expression écrite et Remédiation — Type Bac]
\textbf{Partie A : Composition (10 points)}
Traite l'un des deux sujets suivants (100--120 mots). Ton texte sera évalué sur : la pertinence du contenu (4 pts), la correction grammaticale et lexicale (4 pts), la cohérence et l'organisation (2 pts).
\begin{itemize}
    \item \textbf{Sujet 1 :} “Describe an important social event in your family (wedding, birth celebration, funeral, reunion) and explain why it matters to you.”
    \item \textbf{Sujet 2 :} “Some people think the extended family system is outdated in modern Africa. Do you agree or disagree? Give reasons and examples.”
\end{itemize}

\textbf{Partie B : Remédiation — Chasse aux erreurs (5 points)}
Le paragraphe suivant contient \textbf{8 erreurs} typiques de francophones (grammaire, vocabulaire, ordre des mots, orthographe). Souligne-les, écris la correction et explique brièvement la règle en français.
\begin{textebox}[Student's Draft]
My family is composed by my father, my mother and three childs. We are living in Yaoundé since 2015. My father works in a society of cacao. He is more tall than my mother. Yesterday, he has bought a new car. He said me: “I am very happy today.” We didn't saw him so happy since a long time. He promised to bring us to the village for the holidays.
\end{textebox}
\textbf{Consigne :} Présente tes réponses sous forme de tableau : \textit{Erreur | Correction | Explication}.
\end{exercice}

\begin{savaistu}[Le système des chefferies et la dot au Cameroun]
Au Cameroun, le mariage traditionnel (\eng{traditional marriage}) implique souvent le versement de la \cle{bride price} (dot) par la famille du fiancé à celle de la fiancée. Ce n'est pas un "achat" de la femme, mais une alliance entre deux clans (\eng{clans}). La cérémonie du \cle{knock-door} (frapper à la porte) marque la demande officielle. Dans les régions de l'Ouest et du Nord-Ouest (Bamileke, Bafut, etc.), le \cle{Fon} (chef traditionnel) joue un rôle central. La dot peut inclure de l'argent, des tissus, de l'huile de palme, du bétail, du vin de palme. Le non-paiement peut entraîner des conflits coutumiers. Aujourd'hui, beaucoup de couples font le mariage civil (mairie) et religieux (église) \textbf{après} la dot coutumière.
\end{savaistu}

\begin{methode}[Réussir la composition (Writing) au Bac]
\begin{enumerate}
    \item \textbf{Analyser le sujet :} Identifier le type (description, narration, argumentation), le temps requis (principalement Past Simple pour narration, Present pour description/opinion), le vocabulaire clé.
    \item \textbf{Planifier (5 min) :} Faire un plan en 3 parties : Introduction (accroche + thèse/sujet), Développement (2--3 paragraphes avec arguments/exemples/link words : \eng{Firstly, Moreover, However, For instance}), Conclusion (bilan + ouverture).
    \item \textbf{Rédiger (20 min) :} Écrire des phrases complètes, variées (simples, composées, complexes). Utiliser le vocabulaire du chapitre (\eng{extended family, bride price, gathering, support network}). Vérifier les accords (SV), les temps, l'orthographe.
    \item \textbf{Relire (5 min) :} Chasser les erreurs types (voir Partie B) : \eng{composed of, children, since/for, taller, told me, didn't see, for a long time, cocoa, company}.
\end{enumerate}
\end{methode}

\newpage

\coursec{Corrigés des exercices}

\coursec{Corrigés des activités d'intégration et d'évaluation — Chapter 1 : Family and Social Life}
\courssub{Exercices de révision, type examen et remédiation (1ère A)}

\begin{corrige}[Exercice 1 — QCM : Grammaire et vocabulaire]
\begin{enumerate}
    \item \textbf{b — has lived / since.} \eng{My aunt has lived in Yaoundé since 2010.} « Ma tante habite à Yaoundé depuis 2010. » Le repère \eng{since 2010} (point de départ, durée jusqu'à aujourd'hui) impose le \cle{Present Perfect}. \eng{For} exigerait une durée (\eng{for ten years}) ; \eng{is living since} est impossible (le Present Continuous ne se combine pas avec \eng{since}).
    \item \textbf{c — cousin.} \eng{My father's brother's daughter is my cousin.} « La fille du frère de mon père est ma cousine. » \eng{Niece} = nièce (fille du frère/de la sœur) ; \eng{nephew} = neveu.
    \item \textbf{b — Whose / my sister's.} \eng{Whose book is this? It's my sister's.} « À qui est ce livre ? C'est celui de ma sœur. » \eng{Whose} = à qui (possession) ; le génitif \eng{'s} marque le possesseur. \eng{Who's} = \eng{who is} (contraction), à ne pas confondre.
    \item \textbf{a — more colourful than.} \eng{This traditional wedding was more colourful than the last one.} « Ce mariage traditionnel était plus coloré que le précédent. » Adjectif long (3 syllabes) $\rightarrow$ \eng{more ... than}. \eng{Colourfuller} n'existe pas.
    \item \textbf{c — got married.} \eng{My parents got married twenty years ago.} « Mes parents se sont mariés il y a vingt ans. » Le marqueur \eng{ago} (passé révolu, daté) impose le \cle{Past Simple} ; le Present Perfect est impossible avec \eng{ago}.
\end{enumerate}
\textbf{Points de vigilance :} revoir la paire \eng{since/for}, l'interdiction du Present Perfect avec un repère daté, et la confusion \eng{whose/who's}.
\end{corrige}

\begin{corrige}[Exercice 2 — Vrai ou Faux justifié]
\begin{enumerate}
    \item \textbf{False.} Justification : “Nowadays, young couples often prefer to set up their own nuclear families.” Le texte dit que les jeunes couples préfèrent \emph{souvent} vivre en famille nucléaire, pas que les grands-parents ne vivent \emph{jamais} avec leurs petits-enfants. \eng{Often} $\neq$ \eng{never}.
    \item \textbf{True.} Justification : “Nowadays, young couples often prefer to set up their own nuclear families in cities like Douala or Buea.”
    \item \textbf{False.} Justification : “However, strong ties remain.” Les liens familiaux demeurent forts ; ils n'ont pas disparu.
    \item \textbf{False.} Justification : “The bride price ceremony is still a key step \textbf{before} marriage.” La dot a lieu \emph{avant} le mariage, pas après.
    \item \textbf{True.} Justification : “During holidays, everyone gathers in the village to celebrate weddings, births or funerals.”
\end{enumerate}
\textbf{Méthode :} toujours citer la phrase exacte du texte entre guillemets anglais ; repérer les modificateurs (\eng{often, never, before, completely}) qui rendent une affirmation fausse.
\end{corrige}

\begin{corrige}[Exercice 3 — L'arbre généalogique]
\begin{enumerate}
    \item \textbf{stepfather} — \eng{My mother's new husband is my stepfather.} « Le nouveau mari de ma mère est mon beau-père. »
    \item \textbf{niece} — \eng{My brother's daughter is my niece.} « La fille de mon frère est ma nièce. »
    \item \textbf{in-laws} — \eng{My husband's parents are my in-laws.} « Les parents de mon mari sont mes beaux-parents. »
    \item \textbf{half-sister} — \eng{My half-sister and I share only one parent.} « Ma demi-sœur et moi n'avons qu'un seul parent en commun. »
    \item \textbf{great-grandmother} — \eng{My great-grandmother is the mother of my grandmother.} « Mon arrière-grand-mère est la mère de ma grand-mère. »
    \item \textbf{nephew} — \eng{My sister's son is my nephew.} « Le fils de ma sœur est mon neveu. »
    \item \textbf{spouse} — \eng{My spouse works in a bank.} « Mon conjoint / ma conjointe travaille dans une banque. »
    \item \textbf{siblings} — \eng{My brothers and sisters are my siblings.} « Mes frères et sœurs sont mes frères et sœurs (\eng{siblings}). »
\end{enumerate}
\textbf{Point de vigilance :} \eng{step-} = par remariage (aucun lien de sang) ; \eng{half-} = un seul parent commun ; \eng{-in-law} = par alliance.
\end{corrige}

\begin{corrige}[Exercice 4 — Conjugaison]
\begin{enumerate}
    \item \eng{Look! My little brother \textbf{is playing} with his new toys right now.} « Regarde ! Mon petit frère joue avec ses nouveaux jouets en ce moment. » \cle{Present Continuous} : \eng{Look!} et \eng{right now} signalent une action en cours.
    \item \eng{We \textbf{have known} each other since childhood.} « Nous nous connaissons depuis l'enfance. » \cle{Present Perfect Simple} : \eng{know} est un verbe d'état (pas de \eng{-ing}) et \eng{since} exprime la durée.
    \item \eng{Last December, my sister \textbf{got} engaged during the Christmas party.} « En décembre dernier, ma sœur s'est fiancée pendant la fête de Noël. » \cle{Past Simple} : action datée (\eng{Last December}).
    \item \eng{How long \textbf{have you been waiting} for the bride?} « Depuis combien de temps attends-tu la mariée ? » \cle{Present Perfect Continuous} : action commencée dans le passé, toujours en cours, avec résultat visible (\eng{She is late!}).
    \item \eng{My grandfather \textbf{had never flown} in a plane before his trip to London last year.} « Mon grand-père n'avait jamais pris l'avion avant son voyage à Londres l'année dernière. » \cle{Past Perfect} : antériorité (\eng{before}) par rapport à un événement passé. Participe passé irrégulier : \eng{fly — flew — flown}.
    \item \eng{Usually, the family \textbf{gathers} for Sunday lunch, but today we \textbf{are eating} at a restaurant.} « D'habitude, la famille se réunit pour le déjeuner du dimanche, mais aujourd'hui nous mangeons au restaurant. » \cle{Present Simple} (\eng{usually} : habitude) puis \cle{Present Continuous} (\eng{today} : exception temporaire).
\end{enumerate}
\end{corrige}

\begin{corrige}[Exercice 5 — Compréhension de texte]
\begin{enumerate}
    \item \eng{Grandparents used to live with their children and grandchildren under the same roof.} « Les grands-parents avaient l'habitude de vivre avec leurs enfants et petits-enfants sous le même toit. » (\eng{used to} + base verbale = habitude passée révolue.)
    \item \eng{Nowadays, young couples prefer to set up their own nuclear families in cities like Douala or Buea.} « Aujourd'hui, les jeunes couples préfèrent fonder leur propre famille nucléaire en ville. »
    \item \eng{During holidays, everyone gathers in the village to celebrate weddings, births or funerals.} « Pendant les vacances, tout le monde se réunit au village pour célébrer les mariages, les naissances ou les funérailles. »
    \item \eng{The bride price ceremony takes place before marriage.} « La cérémonie de la dot a lieu avant le mariage. »
    \item \eng{“Evolving”} (l. 2) est le synonyme de \eng{changing}.
    \item \eng{“Strong”} (dans \eng{strong ties}) est le contraire de \eng{weak}.
\end{enumerate}
\textbf{Méthode :} pour les questions de type \eng{find a synonym / the opposite}, recopier uniquement le mot du texte, sans le modifier.
\end{corrige}

\begin{corrige}[Exercice 6 — Traduction (français $\rightarrow$ anglais)]
\begin{enumerate}
    \item \eng{My in-laws have lived in Bamenda for ten years.} (ou \eng{have been living}). « Depuis dix ans » = durée $\rightarrow$ \cle{Present Perfect} + \eng{for}. Faux ami : « beaux-parents » = \eng{parents-in-law / in-laws}, jamais \eng{beautiful parents}.
    \item \eng{This wedding was more traditional than the previous one.} Adjectif long $\rightarrow$ \eng{more ... than} ; \eng{one} évite la répétition de \eng{wedding}.
    \item \eng{My half-sister is older than me.} (ou \eng{than I am}). « Demi-sœur » = \eng{half-sister} ; comparatif de \eng{old} : \eng{older} (\eng{elder} possible devant un nom : \eng{my elder sister}).
    \item \eng{They got married two years ago, but they haven't had any children yet.} « Il y a deux ans » = \eng{ago} $\rightarrow$ Past Simple ; « pas encore » = \eng{not ... yet} $\rightarrow$ Present Perfect.
    \item \eng{How long has your uncle been a widower?} (ou \eng{has your uncle been widowed?}). « Veuf » = \eng{widower} (homme) / \eng{widow} (femme) ; « depuis quand » = \eng{how long} + Present Perfect.
\end{enumerate}
\textbf{Barème indicatif :} 1 point par phrase (0,5 pour le temps correct, 0,5 pour le vocabulaire et l'ordre des mots).
\end{corrige}

\begin{corrige}[Exercice 7 — Transformation de phrases]
\begin{enumerate}
    \item \eng{My mother is not \textbf{as old as my father}.} « Ma mère n'est pas aussi âgée que mon père. » Comparatif d'égalité nié : \eng{not as ... as}.
    \item \eng{I have never \textbf{seen such a beautiful traditional outfit}.} (ou \eng{seen a more beautiful traditional outfit}). « Je n'ai jamais vu une tenue traditionnelle aussi belle. » Superlatif $\rightarrow$ \eng{never ... such a / a more ...}.
    \item \eng{He asked me \textbf{where my cousin lived}.} Discours rapporté : question en \eng{wh-} $\rightarrow$ \eng{asked me where} + ordre sujet-verbe (pas d'inversion) + backshift (\eng{does ... live} $\rightarrow$ \eng{lived}).
    \item \eng{The dowry ceremony \textbf{was organised last week}.} Passif : \eng{be} au Past Simple (\eng{was}) + participe passé (\eng{organised}) ; le complément d'agent \eng{by them} est omis car inutile.
    \item \eng{My nephew is not \textbf{old enough to get married}.} « Mon neveu n'est pas assez âgé pour se marier. » Transformation \eng{too + adj + to} $\rightarrow$ \eng{not + adj + enough + to} ; \eng{enough} se place \emph{après} l'adjectif.
\end{enumerate}
\textbf{Points de vigilance :} ne jamais garder l'inversion dans le discours rapporté (\eng{*where did my cousin live}) ; ne pas placer \eng{enough} avant l'adjectif.
\end{corrige}

\begin{corrige}[Exercice 8 — Texte à trous]
\begin{enumerate}
    \item \textbf{gathers} (ou \eng{meets}) — Present Simple, habitude annuelle (\eng{every year}).
    \item \textbf{is} — \eng{It is a tradition...}
    \item \textbf{has lasted} (ou \eng{has been lasting}) — Present Perfect : durée \eng{for generations} jusqu'à aujourd'hui.
    \item \textbf{uncle} — le frère de ma mère.
    \item \textbf{with} — \eng{arrived with his wife} (accompagnement).
    \item \textbf{wife} — épouse.
    \item \textbf{have} — auxiliaire du Present Perfect.
    \item \textbf{seen} — \eng{We have not seen them for two years.} Participe passé : \eng{see — saw — seen}.
    \item \textbf{eldest} (ou \eng{oldest}) — superlatif : le cousin le plus âgé.
    \item \textbf{has} — auxiliaire du Present Perfect (3e personne).
    \item \textbf{passed} (ou \eng{taken}) — \eng{has just passed his exams} « vient de réussir ses examens » ; \eng{just} impose le Present Perfect.
    \item \textbf{more} — comparatif d'adjectif long.
    \item \textbf{than} — \eng{more intelligent than his younger brother}.
    \item \textbf{grandparents} — les grands-parents racontent des histoires.
    \item \textbf{ancestors} — \eng{our ancestors} « nos ancêtres ».
\end{enumerate}
\textbf{Texte reconstitué :} \eng{Every year, during the long holidays, my extended family gathers in our ancestral village. It is a tradition that has lasted for generations. This year, my uncle arrived with his wife and their three children. We have not seen them for two years. The eldest cousin, Paul, has just passed his exams. He is more intelligent than his younger brother, but less serious. In the evening, the grandparents tell stories about our ancestors.}
\end{corrige}

\begin{corrige}[Exercice 9 — Appariements]
\begin{center}
\begin{tabular}{|c|c|l|}
\hline
\rowcolor{popblueL} \textbf{Situation} & \textbf{Réponse} & \textbf{Justification} \\
\hline
1 & B & Présentation : \eng{Dad, this is Mary, my fiancée.} (\eng{this is...} pour présenter quelqu'un). \\
\hline
2 & A & Durée du mariage : \eng{How long have you been married?} (Present Perfect + \eng{how long}). \\
\hline
3 & D & Comparaison : \eng{The first ceremony was more formal than the second.} (comparatif). \\
\hline
4 & C & Bonnes relations : \eng{My sister gets on well with her in-laws.} (\eng{to get on well with} = bien s'entendre avec). \\
\hline
5 & F & Nouvelles : \eng{How is your nephew doing at school?} (\eng{How is ... doing?} = comment va...). \\
\hline
6 & E & Description physique : \eng{What does your nephew look like?} (\eng{look like} = ressembler à). \\
\hline
\end{tabular}
\end{center}
\textbf{Point de vigilance :} \eng{What does he look like?} (apparence) $\neq$ \eng{What is he like?} (caractère) $\neq$ \eng{What does he like?} (goûts).
\end{corrige}

\begin{corrige}[Exercice 10 — Compréhension écrite Type Bac]
\textbf{1. True or False (justifier par une citation) :}
\begin{enumerate}
    \item \textbf{False.} “Uncles, aunts and grandparents played active roles in education and discipline.” Dans le passé, toute la famille élargie participait à la discipline, pas seulement les parents.
    \item \textbf{False.} “The nuclear family — parents and children only — has become the dominant model in urban areas” et “the support network has weakened.” L'urbanisation a \emph{affaibli}, et non renforcé, la famille élargie.
    \item \textbf{True.} “Video calls allow daily contact with grandparents.” La technologie maintient les liens familiaux.
\end{enumerate}
\textbf{2. Multiple Choice :}
\begin{enumerate}
    \item \textbf{b)} \eng{the most important part} — \eng{cornerstone} = « pierre angulaire », au sens figuré : élément essentiel.
    \item \textbf{b)} \eng{The support network has weakened.} Le texte : “Consequently, the support network has weakened. Working parents often rely on housemaids or day-care centres.”
\end{enumerate}
\textbf{3. Short answers (max 20 mots) :}
\begin{enumerate}
    \item \eng{Weddings, funerals and bride price ceremonies still bring the extended family together.} (deux événements suffisent).
    \item \eng{They send financial remittances via mobile money to sustain rural homes.}
\end{enumerate}
\textbf{4. Language work :}
\begin{enumerate}
    \item Synonyme de \eng{old people} : \eng{the elderly} (paragraphe 2).
    \item Contraire de \eng{strengthened} : \eng{weakened} (paragraphe 2).
    \item Passif : \eng{Young adults have been driven to cities by urbanisation.} (\eng{have been} + participe passé \eng{driven} ; \eng{drive — drove — driven}).
    \item Discours rapporté : \eng{The sociologist said (that) the clan identity remained strong.} (backshift : \eng{remains} $\rightarrow$ \eng{remained}).
\end{enumerate}
\textbf{Barème indicatif (sur 10) :} Q1 : 3 pts (0,5 réponse + 0,5 citation) ; Q2 : 2 pts ; Q3 : 2 pts ; Q4 : 3 pts (0,75 par item).
\end{corrige}

\begin{corrige}[Exercice 11 — Grammaire en contexte et Traduction Type Bac]
\textbf{Partie A — Conjugaison :}
\begin{enumerate}
    \item \textbf{was} — \eng{the Fon's palace was full of guests} : Past Simple, récit daté (\eng{Last Saturday}).
    \item \textbf{had just returned} — Past Perfect : le retour d'Allemagne est antérieur au samedi du mariage (\eng{just} + antériorité).
    \item \textbf{had been studying} (ou \eng{had studied}) — Past Perfect Continuous : durée (\eng{for four years}) avant un moment passé.
    \item \textbf{had not seen} — Past Perfect avec \eng{since} + point de départ passé (\eng{since the knock-door ceremony two years ago}) : il n'avait pas vu Njwen depuis cette cérémonie.
    \item \textbf{was} — \eng{by the time the bride price was paid} : Past Simple dans la subordonnée de temps.
    \item \textbf{had already set} — Past Perfect : le soleil s'était déjà couché (antériorité). \eng{set — set — set} (irrégulier, invariable).
    \item \textbf{announced} — Past Simple : action du récit.
    \item \textbf{had become} — discours rapporté avec backshift (\eng{announced that the couple had become...}).
    \item \textbf{have lived} (ou \eng{have been living}) — \eng{Since then} « depuis lors » jusqu'à aujourd'hui $\rightarrow$ Present Perfect.
\end{enumerate}
\textbf{Barème indicatif :} 1 point par verbe correctement conjugué (10 pts).

\textbf{Partie B — Traduction :}
\begin{enumerate}
    \item \eng{If I had known, I would have come to the bride price ceremony.} Conditionnel de type 3 (hypothèse irréelle du passé) : \eng{If + Past Perfect, ... would have + participe passé}.
    \item \eng{My grandmother used to tell me stories when I was a child.} (ou \eng{would tell me}). Habitude passée révolue : \eng{used to} + base verbale ; \eng{when I was little} également accepté.
\end{enumerate}
\textbf{Barème indicatif :} 2,5 pts par phrase (structure du conditionnel / de l'habitude passée : 1,5 pt ; vocabulaire : 1 pt).
\end{corrige}

\begin{corrige}[Exercice 12 — Expression écrite et Remédiation Type Bac]
\textbf{Partie A — Proposition de rédaction modèle (Sujet 2) :}
\begin{textebox}[Model Answer — Opinion Essay]
I strongly disagree with the idea that the extended family is outdated in modern Africa. In my opinion, it remains the backbone of our society.

Firstly, the extended family provides financial and emotional support. When my uncle lost his job in Douala, the whole family contributed money until he found a new one. Moreover, grandparents play an essential role in raising children: they transmit values, languages and traditions that no day-care centre can teach.

Admittedly, urbanisation has weakened this system because young people move to cities and live in nuclear families. However, technology now bridges the gap: video calls and mobile money keep relatives connected across long distances.

To conclude, the extended family has not disappeared; it has simply adapted. As the proverb says, “it takes a village to raise a child”, and this remains true today.
\end{textebox}
(Environ 120 mots. « Je ne suis absolument pas d'accord avec l'idée que la famille élargie est dépassée dans l'Afrique moderne. À mon avis, elle reste la colonne vertébrale de notre société. D'abord, elle apporte un soutien financier et affectif... Certes, l'urbanisation a affaibli ce système... Cependant, la technologie comble désormais la distance... En conclusion, la famille élargie n'a pas disparu ; elle s'est simplement adaptée. »)

\textbf{Barème indicatif (10 pts) :} contenu et pertinence : 4 pts ; grammaire et vocabulaire : 4 pts ; cohérence, organisation et connecteurs (\eng{firstly, moreover, however, to conclude}) : 2 pts.

\textbf{Partie B — Chasse aux erreurs (10 erreurs) :}
\begin{center}
\begin{tabularx}{\linewidth}{|X|X|X|}
\hline
\rowcolor{popblueL} \textbf{Erreur} & \textbf{Correction} & \textbf{Explication} \\
\hline
is composed \textbf{by} & is composed \textbf{of} & \eng{To be composed of} = être composé de. Préposition fixe. \\
\hline
three \textbf{childs} & three \textbf{children} & Pluriel irrégulier : \eng{child} $\rightarrow$ \eng{children}. \\
\hline
\textbf{are living} since 2015 & \textbf{have lived / have been living} since 2015 & Durée depuis un point de départ $\rightarrow$ Present Perfect, jamais Present Continuous avec \eng{since}. \\
\hline
a \textbf{society} of cacao & a cocoa \textbf{company} & Faux ami : \eng{society} = société (la collectivité) ; « entreprise » = \eng{company/firm}. Orthographe : \eng{cocoa}. \\
\hline
\textbf{more tall} than & \textbf{taller} than & Adjectif d'une syllabe $\rightarrow$ comparatif en \eng{-er}. Jamais \eng{more} + \eng{-er}. \\
\hline
he \textbf{has bought} (yesterday) & he \textbf{bought} & \eng{Yesterday} = passé daté $\rightarrow$ Past Simple obligatoire. \\
\hline
He \textbf{said me} & He \textbf{told me} & \eng{Tell} + personne (sans \eng{to}) ; \eng{say} + \eng{to} + personne : \eng{He said to me...} \\
\hline
We \textbf{didn't saw} & We \textbf{didn't see} & Après l'auxiliaire \eng{did}, la base verbale (sans \eng{-ed}). \\
\hline
\textbf{since} a long time & \textbf{for} a long time & \eng{For} + durée ; \eng{since} + point de départ. \\
\hline
to \textbf{bring} us to the village & to \textbf{take} us to the village & \eng{Take} = emmener (vers un lieu éloigné) ; \eng{bring} = amener (vers le locuteur). \\
\hline
\end{tabularx}
\end{center}
\textbf{Barème indicatif (5 pts) :} 0,5 pt par erreur repérée et corrigée (l'explication de la règle est un bonus méthodologique).

\textbf{Texte corrigé complet :} \eng{My family is composed of my father, my mother and three children. We have lived in Yaoundé since 2015. My father works in a cocoa company. He is taller than my mother. Yesterday, he bought a new car. He told me: “I am very happy today.” We haven't seen him so happy for a long time. He promised to take us to the village for the holidays.}
\end{corrige}

\newpage

\coursec{Fiche bilan}

\begin{popfigure}
\begin{tikzpicture}[
  mindmap,
  concept color=popblue,
  level 1 concept/.append style={level distance=3.5cm, sibling angle=360/6},
  level 2 concept/.append style={level distance=2.5cm},
  every node/.style={font=\small, align=center},
  grow cyclic,
  branch/.style={concept color=#1, edge from parent/.style={draw=#1, thick}}
]
\node[concept, text=white, minimum size=2.2cm, align=center]{Family \&\\ Social Life}
  child[branch=poporange] { node[concept, text=white, minimum size=1.8cm, align=center]{Vocabulary\\ \textit{Key Lexis}}
    child { node[concept, text=white, minimum size=1.4cm] {Family ties} }
    child { node[concept, text=white, minimum size=1.4cm] {Life events} }
    child { node[concept, text=white, minimum size=1.4cm, align=center]{Adjectives\\ feelings} }
  }
  child[branch=popgreen] { node[concept, text=white, minimum size=1.8cm, align=center]{Grammar\\ \textit{Core Points}}
    child { node[concept, text=white, minimum size=1.4cm, align=center]{Present\\ Simple/Cont.} }
    child { node[concept, text=white, minimum size=1.4cm, align=center]{Past Simple\\ vs Continuous} }
    child { node[concept, text=white, minimum size=1.4cm, align=center]{Relative\\ clauses} }
    child { node[concept, text=white, minimum size=1.4cm, align=center]{Modals\\ obligation} }
  }
  child[branch=poppurple] { node[concept, text=white, minimum size=1.8cm, align=center]{Reading\\ \textit{Integration}}
    child { node[concept, text=white, minimum size=1.4cm] {Skimming} }
    child { node[concept, text=white, minimum size=1.4cm] {Scanning} }
    child { node[concept, text=white, minimum size=1.4cm] {Inference} }
  }
  child[branch=poppink] { node[concept, text=white, minimum size=1.8cm, align=center]{Writing\\ \textit{Production}}
    child { node[concept, text=white, minimum size=1.4cm, align=center]{Informal\\ letter} }
    child { node[concept, text=white, minimum size=1.4cm] {Description} }
    child { node[concept, text=white, minimum size=1.4cm] {Narrative} }
  }
  child[branch=popgold] { node[concept, text=white, minimum size=1.8cm, align=center]{Speaking/\\ Listening}
    child { node[concept, text=white, minimum size=1.4cm] {Role-play} }
    child { node[concept, text=white, minimum size=1.4cm] {Interview} }
    child { node[concept, text=white, minimum size=1.4cm] {Pronunciation} }
  }
  child[branch=popdark] { node[concept, text=white, minimum size=1.8cm, align=center]{Exam Prep\\ \textit{Strategy}}
    child { node[concept, text=white, minimum size=1.4cm] {Integration} }
    child { node[concept, text=white, minimum size=1.4cm] {Self-eval} }
    child { node[concept, text=white, minimum size=1.4cm] {Remediation} }
  };
\end{tikzpicture}
\legende{Carte mentale de synthèse — Chapter 1 : Family and Social Life}
\end{popfigure}

\begin{bilan}

\end{bilan}

\courssub{Vocabulary Synthesis — Lexique clé (English / Français)}

\begin{center}
\begin{tabularx}{\linewidth}{|X|X|X|}
\hline
\rowcolor{popblueL} \textbf{Category} & \textbf{English} & \textbf{Français} \\
\hline
Family ties & \eng{nuclear family, extended family, sibling, spouse, in-laws, stepfather, half-brother, adopted child} & famille nucléaire, famille élargie, frère/sœur, conjoint, beaux-parents, beau-père, demi-frère, enfant adopté \\
\hline
Life events & \eng{birth, wedding, funeral, graduation, engagement, anniversary, naming ceremony, bride price} & naissance, mariage, enterrement, remise de diplômes, fiançailles, anniversaire, cérémonie de baptême, dot \\
\hline
Adjectives (feelings) & \eng{proud, jealous, supportive, strict, caring, close-knit, dysfunctional, overprotective} & fier, jaloux, solidaire, strict, attentionné, soudé, dysfonctionnel, surprotecteur \\
\hline
Social life & \eng{peer pressure, generation gap, household chores, upbringing, role model, gathering, reunion, inheritance} & pression des pairs, fossé des générations, tâches ménagères, éducation, modèle, réunion, retrouvailles, héritage \\
\hline
\end{tabularx}
\end{center}

\courssub{Grammar Synthesis — Règles à connaître par cœur}

\begin{center}
\begin{tabularx}{\linewidth}{|X|X|X|}
\hline
\rowcolor{popgreenL} \textbf{Point} & \textbf{Forme / Structure} & \textbf{Exemple traduit} \\
\hline
Present Simple vs Continuous & \textbf{Simple:} Habits, facts (\eng{V/Vs}). \textbf{Cont:} Now, temporary (\eng{am/is/are + V-ing}). & \eng{She lives in Yaoundé.} (Elle habite…)\\
& & \eng{She is staying with her aunt this week.} (Elle séjourne… cette semaine) \\
\hline
Past Simple vs Continuous & \textbf{Simple:} Finished action (\eng{V-ed/irreg}). \textbf{Cont:} Action in progress (\eng{was/were + V-ing}). & \eng{While I was cooking, the phone rang.} (Pendant que je cuisinais, le téléphone a sonné) \\
\hline
Relative Clauses & \textbf{Who} (people), \textbf{Which} (things), \textbf{Where} (place), \textbf{Whose} (possession). No comma if defining. & \eng{The man who called is my uncle.} (L'homme qui a appelé est mon oncle) \\
\hline
Modals (Obligation/Advice) & \textbf{Must} (internal), \textbf{Have to} (external), \textbf{Should} (advice). \textbf{Don't have to} = absence d'obligation. & \eng{You must respect elders.} / \eng{You have to wear a uniform.} / \eng{You should listen.} \\
\hline
Possessives & \textbf{'s} (singular), \textbf{s'} (plural), \textbf{of} (things). Double genitive: \eng{a friend of mine}. & \eng{My parents' house.} (La maison de mes parents) / \eng{A cousin of mine.} (Un de mes cousins) \\
\hline
\end{tabularx}
\end{center}

\courssub{Méthodes Express — Réflexes d'examen}

\begin{itemize}
\item \textbf{Reading Integration:} 1) Lire le titre + images (prédiction). 2) \textit{Skimming} : idée générale (1ère phrase par paragraphe). 3) \textit{Scanning} : mots-clés (dates, noms, \cle{who/where/when}). 4) Inférence : déduire le sens des mots inconnus par le contexte.
\item \textbf{Writing (Informal Letter):} 1) Adresse + Date (en haut à droite). 2) Salutation (\eng{Dear [Prénom],}). 3) Opening (\eng{Thanks for your letter...}). 4) Corps : 2-3 paragraphes (nouvelles, description, invitation). 5) Closing (\eng{Best wishes, / Love,}). 6) Signature (prénom seulement).
\item \textbf{Grammar Gap-fill:} Identifier le marqueur temporel (\eng{now, yesterday, since, while}) $\rightarrow$ Choisir le temps $\rightarrow$ Vérifier l'accord (3e pers. sg \eng{s}, auxiliaire \eng{be/do/have}).
\item \textbf{Speaking (Role-play):} Utiliser les \cle{gambits} : \eng{“What do you think about…?”}, \eng{“I agree/disagree because…”}, \eng{“Could you repeat, please?”}. Regarder l'interlocuteur, pas la feuille.
\end{itemize}



\begin{aretenir}[Checklist avant le Bac — Chapter 1]
$\square$ Je sais décrire mon arbre généalogique (vocabulaire \eng{step-, half-, in-laws}) sans confondre \eng{brother-in-law} et \eng{half-brother}.\\
$\square$ Je distingue \eng{Present Simple} (habitude) et \eng{Present Continuous} (en cours / futur proche programmé) et j'applique la règle du \eng{s} à la 3e personne.\\
$\square$ Je raconte un souvenir d'enfance en combinant correctement \eng{Past Continuous} (action longue) et \eng{Past Simple} (action courte qui l'interrompt) avec \eng{when/while}.\\
$\square$ Je relie deux phrases simples en une phrase complexe avec \eng{who/which/where/whose} sans oublier de supprimer le pronom redondant (\eng{*The man who I know him* $\rightarrow$ The man I know / The man who knows me}).\\
$\square$ Je choisis le modal adapté : \eng{must} (règle interne/forte), \eng{have to} (règle externe/loi), \eng{should} (conseil), et je forme la négation \eng{don't have to} (pas \eng{mustn't} = interdiction).\\
$\square$ J'écris une lettre informelle structurée (adresse, date, formules d'appel/de politesse, paragraphes distincts, signature).\\
$\square$ Je comprends un dialogue authentique sur la vie familiale (accents variés) et je relève les informations explicites et implicites.\\
$\square$ Je participe à un jeu de rôle \eng{“Family meeting”} ou \eng{“Generations debate”} en exprimant accord/désaccord (\eng{I see your point, but…}) et en utilisant le vocabulaire des sentiments.\\
$\square$ Je repère mes erreurs types (faux-amis : \eng{actually} $\neq$ actuellement, \eng{parents} $\neq$ parents [= père et mère], \eng{eventually} $\neq$ éventuellement) et je les corrige seul.\\
$\square$ Je gère mon temps à l'examen : 15 min lecture, 45 min écriture/grammaire, 10 min relecture (accords, orthographe, ponctuation).
\end{aretenir}

\findecours

\end{document}
